\documentclass[11pt,openany]{book}
% compilar com XeLaTeX (fonte da marca Nexus BI em fontes/)
\usepackage{mathpazo}
\usepackage[no-math]{fontspec}
\setmainfont{TeX Gyre Pagella}
\newfontfamily\barlow{BarlowSemiCondensed}[Path=fontes/, Extension=.ttf, Ligatures=TeX,
  UprightFont=*-Regular, BoldFont=*-Bold,
  FontFace={m}{n}{*-Medium}, FontFace={sb}{n}{*-SemiBold}]
\usepackage[brazil]{babel}
\usepackage{microtype}
\usepackage{amsmath,amssymb,mathtools}
\usepackage[a4paper,margin=2.3cm,includefoot]{geometry}
\usepackage{xcolor}
\usepackage[most]{tcolorbox}
\usepackage{graphicx}
\usepackage[export]{adjustbox}
\usepackage{enumitem}
\usepackage{booktabs,array,colortbl}
\usepackage{fancyhdr}
\usepackage{titlesec}
\usepackage{textcomp}
\usepackage{multicol}
\usepackage{pgfplots}
\usepackage{imakeidx}
\makeindex[title=Índice remissivo,columns=2,intoc,options=-s indice.ist]
\pgfplotsset{compat=1.18}
\usepgfplotslibrary{groupplots,fillbetween}
\usetikzlibrary{arrows.meta,calc}
\pgfplotsset{painel/.style={axis lines=middle,width=5.4cm,height=4.4cm,
  tick label style={font=\scriptsize},label style={font=\footnotesize},
  xlabel={$x$},ylabel={},xlabel style={anchor=west},ylabel style={anchor=south},
  title style={font=\footnotesize,yshift=-2pt},samples=200,
  every axis plot/.append style={thick,azul},clip=true},
  }
\tikzset{assint/.style={dashed,laranja,thin}}
\newcommand{\figtikz}[2]{\begin{center}#1\\[2pt]
  \begin{minipage}{0.9\linewidth}\small\textbf{Figura \thechapter.\stepcounter{figure}\arabic{figure}}\quad #2\end{minipage}\end{center}}
\usepackage[hidelinks]{hyperref}
\hypersetup{colorlinks=true,linkcolor=azul,urlcolor=azul!80!black}

% Identidade visual Nexus BI (ver Manual de Identidade Visual):
%   azul    -> grafite #353535 (cor institucional; texto, títulos, curvas principais)
%   verde   -> verde-nexus #AFE213 (cor institucional; preenchimentos e destaques)
%   laranja -> verde-escuro #6E8F0B (derivado do verde-nexus para linhas e símbolos
%              sobre branco, onde o verde-nexus não tem contraste suficiente)
% Os nomes antigos foram mantidos para não reescrever as figuras.
\definecolor{azul}{HTML}{353535}
\definecolor{verde}{HTML}{AFE213}
\definecolor{laranja}{HTML}{6E8F0B}
\definecolor{cinza}{HTML}{595959}
\colorlet{grafite}{azul}\colorlet{verdenexus}{verde}\colorlet{verdeescuro}{laranja}

\setlength{\parindent}{0pt}
\setlength{\parskip}{0.55em}
\setlist{itemsep=2pt,topsep=3pt}

\titleformat{\chapter}[display]{\normalfont\barlow\bfseries\color{azul}}
  {\large\color{laranja}CAPÍTULO \thechapter}{2pt}{\Huge}[{\color{verde}\titlerule[2.5pt]}]
\titlespacing*{\chapter}{0pt}{0pt}{24pt}
\titleformat{\section}{\normalfont\barlow\Large\bfseries\color{azul}}{\color{laranja}\thesection}{0.6em}{}
\titleformat{\subsection}{\normalfont\barlow\large\bfseries\color{azul}}{}{0em}{}
\titleformat*{\subsubsection}{\normalfont\normalsize\bfseries}
\setcounter{secnumdepth}{1}

\pagestyle{fancy}\fancyhf{}
\fancyhead[LE]{\small\barlow\color{cinza}\leftmark}
\fancyhead[RO]{\small\barlow\color{cinza}\rightmark}
\fancyfoot[C]{\small\thepage}
\fancyfoot[R]{\scriptsize\color{cinza}Access for free at openstax.org}
\fancypagestyle{plain}{\fancyhf{}\fancyfoot[C]{\small\thepage}
  \fancyfoot[R]{\scriptsize\color{cinza}Access for free at openstax.org}\renewcommand{\headrulewidth}{0pt}}
\renewcommand{\headrulewidth}{0.3pt}
\let\cleardoublepage\clearpage

% ---------- caixas
\tcbset{base/.style={enhanced,breakable,arc=2pt,boxrule=0.6pt,fonttitle=\barlow\bfseries,
  left=6pt,right=6pt,top=4pt,bottom=4pt,before skip=8pt,after skip=8pt}}
\newtcolorbox{definicao}[1][]{base,colback=azul!4,colframe=azul,
  title={\ifx&#1&Definição\else #1\fi}}
\newtcolorbox{regra}[1]{base,colback=azul!4,colframe=azul,title={#1}}
\newtcolorbox{teorema}[2][]{base,colback=verde!7,colframe=verde,coltitle=azul,
  title={Teorema #2\ifx&#1&\else: #1\fi}}
\newtcolorbox{estrategia}[1]{base,colback=white,colframe=azul,coltitle=verde,
  title={Estratégia: #1}}
\newtcolorbox{exemplo}[2]{base,colback=white,colframe=cinza,coltitle=black,colbacktitle=cinza!15,
  title={Exemplo #1 --- #2}}
\newtcolorbox{nota}{base,colback=red!3,colframe=red!55!black,title={Atenção}}
\newcommand{\solucao}{\par\smallskip{\barlow\color{laranja}\textbf{Solução}}\par}
\newcommand{\fig}[3][0.62]{\begin{center}\includegraphics[width=#1\linewidth,max height=7.2cm,keepaspectratio]{img/#2}\\[2pt]
  \begin{minipage}{0.9\linewidth}\small\textbf{Figura \thechapter.\stepcounter{figure}\arabic{figure}}\quad #3\end{minipage}\end{center}}
\newcommand{\figdupla}[4][0.46]{\begin{center}
  \includegraphics[width=#1\linewidth,max height=6cm,keepaspectratio]{img/#2}\hfill
  \includegraphics[width=#1\linewidth,max height=6cm,keepaspectratio]{img/#3}\\[2pt]
  \begin{minipage}{0.9\linewidth}\small\textbf{Figura \thechapter.\stepcounter{figure}\arabic{figure}}\quad #4\end{minipage}\end{center}}

\newcommand{\R}{\mathbb{R}}
\newcommand{\nivel}[1]{{\ifnum#1=4\color{laranja}\else\color{azul}\fi
  \foreach \i in {1,...,#1}{\textbullet}}}
\newcommand{\exnum}[2]{\makebox[1.9em][r]{\nivel{#1}}\,\makebox[2.3em][r]{\textbf{#2}}}
\newcommand{\estrela}{\textcolor{laranja}{$\bigstar$}}
% notação brasileira das funções trigonométricas e hiperbólicas
\DeclareMathOperator{\sen}{sen}
\DeclareMathOperator{\tg}{tg}
\DeclareMathOperator{\cotg}{cotg}
\DeclareMathOperator{\cossec}{cossec}
\DeclareMathOperator{\arcsen}{arcsen}
\DeclareMathOperator{\arctg}{arctg}
\DeclareMathOperator{\arcsec}{arcsec}
\DeclareMathOperator{\arccotg}{arccotg}
\DeclareMathOperator{\arccossec}{arccossec}
\DeclareMathOperator{\senh}{senh}
\DeclareMathOperator{\tgh}{tgh}
\DeclareMathOperator{\cotgh}{cotgh}
\DeclareMathOperator{\sech}{sech}
\DeclareMathOperator{\cossech}{cossech}
\DeclareMathOperator{\arcsenh}{arcsenh}
\DeclareMathOperator{\arccosh}{arccosh}
\DeclareMathOperator{\arctgh}{arctgh}
\newcommand{\Set}[2]{\{\,#1 \mid #2\,\}}

\begin{document}
\frontmatter
% Capa do livro (TikZ, página inteira)
\definecolor{noite}{HTML}{0A1A33}
\definecolor{marinho}{HTML}{153A63}
\definecolor{ouro}{HTML}{E8A33D}
\begin{titlepage}
\thispagestyle{empty}
\begin{tikzpicture}[remember picture,overlay,
  declare function={curva(\x)=8.6+3.2*sin(deg(0.42*\x-0.9))+0.9*sin(deg(1.1*\x));}]
% fundo
\shade[top color=noite,bottom color=marinho] (current page.south west) rectangle (current page.north east);
\begin{scope}[shift={(current page.south west)}]
  % grade discreta
  \draw[white,opacity=0.05,line width=0.3pt,step=1cm] (0,0) grid (21,29.7);
  % retângulos de Riemann: cada vez mais finos da esquerda para a direita
  \foreach \a/\w/\op in {1.0/1.6/0.10, 2.6/1.6/0.12, 4.2/1.2/0.14, 5.4/1.2/0.16,
                         6.6/0.8/0.18, 7.4/0.8/0.20, 8.2/0.8/0.22, 9.0/0.5/0.24,
                         9.5/0.5/0.26, 10.0/0.5/0.28, 10.5/0.5/0.30, 11.0/0.3/0.32,
                         11.3/0.3/0.34, 11.6/0.3/0.36, 11.9/0.3/0.38, 12.2/0.3/0.40}{
    \pgfmathsetmacro{\m}{\a+\w/2}
    \fill[white,opacity=\op] (\a,2.2) rectangle ({\a+\w},{curva(\m)});
    \draw[white,opacity=\op+0.15,line width=0.3pt] (\a,2.2) rectangle ({\a+\w},{curva(\m)});
  }
  % área exata (preenchimento contínuo) à direita
  \fill[ouro,opacity=0.30] (12.5,2.2) -- plot[domain=12.5:20,samples=80] (\x,{curva(\x)}) -- (20,2.2) -- cycle;
  % eixo
  \draw[white,opacity=0.5,line width=0.6pt,-{Stealth}] (0.4,2.2) -- (20.6,2.2);
  % curva
  \draw[white,line width=1.8pt] plot[domain=0.5:20.5,samples=200] (\x,{curva(\x)});
  \begin{scope}\clip (0.4,2.25) rectangle (20.6,20);
  % secantes convergindo para a tangente em x0
  \pgfmathsetmacro{\xo}{16.4}
  \pgfmathsetmacro{\yo}{curva(\xo)}
  \pgfmathsetmacro{\dm}{(curva(\xo+0.001)-curva(\xo-0.001))/0.002}
  \foreach \h/\op in {4.2/0.18, 2.8/0.26, 1.7/0.36, 0.9/0.48}{
    \pgfmathsetmacro{\ms}{(curva(\xo+\h)-\yo)/\h}
    \draw[ouro,opacity=\op,line width=0.8pt] ({\xo-3.2},{\yo-3.2*\ms}) -- ({\xo+\h+1.2},{\yo+(\h+1.2)*\ms});
    \fill[ouro,opacity=\op+0.2] ({\xo+\h},{curva(\xo+\h)}) circle (2.2pt);
  }
  \draw[ouro,line width=1.6pt] ({\xo-2.6},{\yo-2.6*\dm}) -- ({\xo+2.6},{\yo+2.6*\dm});
  \fill[ouro] (\xo,\yo) circle (3.6pt);
  \fill[noite] (\xo,\yo) circle (1.6pt);
  \end{scope}
  % fórmulas em marca-d'água
  \node[white,opacity=0.16,scale=1.5] at (5.2,14.6)
    {$\displaystyle f'(x)=\lim_{h\to0}\frac{f(x+h)-f(x)}{h}$};
  \node[white,opacity=0.16,scale=1.5] at (14.6,14.6)
    {$\displaystyle\int_a^b f(x)\,dx=F(b)-F(a)$};
  % título
  \node[anchor=north west,text=ouro,font=\large\scshape] at (2,27.4)
    {E\,d\,i\,ç\,ã\,o\ \ c\,o\,n\,d\,e\,n\,s\,a\,d\,a\ \ e\,m\ \ p\,o\,r\,t\,u\,g\,u\,ê\,s};
  \draw[ouro,line width=1.2pt] (2,26.55) -- (6.5,26.55);
  \node[anchor=north west,text=white,font=\bfseries,scale=6.4] at (1.75,26.2) {Cálculo\,\textcolor{ouro}{I}};
  \node[anchor=north west,text=white,opacity=0.9,font=\Large] at (2,21.6)
    {Funções \textperiodcentered{} Limites \textperiodcentered{} Derivadas \textperiodcentered{} Integrais};
  % rodapé da capa
  \draw[ouro,line width=0.8pt] (2,1.55) -- (19,1.55);
  \node[anchor=west,text=white,opacity=0.85,font=\small] at (2,1.05)
    {Baseado em \emph{Calculus Volume 1}, de Gilbert Strang e Edwin ``Jed'' Herman};
  \node[anchor=east,text=white,opacity=0.6,font=\scriptsize] at (19,1.05)
    {Access for free at openstax.org};
\end{scope}
\end{tikzpicture}
\end{titlepage}


\thispagestyle{plain}
\section*{Créditos e licença}
Este livro é uma obra derivada de \emph{Calculus Volume 1}, de Gilbert Strang e Edwin
``Jed'' Herman, publicado pela OpenStax (Rice University) em 2016 e disponível
gratuitamente em \url{https://openstax.org/details/books/calculus-volume-1}.
A obra original é licenciada sob a
\href{https://creativecommons.org/licenses/by-nc-sa/4.0/}{Creative Commons
Atribuição-NãoComercial-CompartilhaIgual 4.0 Internacional (CC BY-NC-SA 4.0)}.
O capítulo~7 (Técnicas de integração) é adaptado do capítulo~3 de \emph{Calculus Volume 2},
de Gilbert Strang e Edwin ``Jed'' Herman (OpenStax, 2016,
\url{https://openstax.org/details/books/calculus-volume-2}), sob a mesma licença.

\textbf{Alterações feitas nesta versão:} tradução para o português; condensação do texto
(redução de explicações, exemplos e figuras); reorganização de algumas passagens; remoção dos
pontos de verificação, dos projetos e dos exercícios do original. Todas as figuras foram
redesenhadas em português. As listas de exercícios ao fim de cada capítulo foram
elaboradas para esta edição, inspiradas nos temas e exemplos do original; os exercícios
marcados com \estrela{} têm solução completa e animação no repositório do livro
(\url{https://github.com/AleffAzevedo/calculo-1-pt}). Foram acrescentados um apêndice de
fórmulas e um índice remissivo. No texto, os decimais usam vírgula, no padrão brasileiro.

Esta versão não é endossada pela OpenStax nem pela Rice University. Os nomes e logotipos
OpenStax e Rice University são marcas registradas e não fazem parte da licença.
Esta adaptação é distribuída sob a mesma licença CC BY-NC-SA 4.0, apenas para uso não comercial.

\textbf{Projeto gráfico:} identidade visual da Nexus BI (cores grafite e verde-nexus, padrão
gráfico e logotipo); títulos em Barlow Semi Condensed, fonte distribuída sob a SIL Open Font
License. O logotipo e a marca Nexus BI pertencem a seus titulares e não estão cobertos pela
licença CC BY-NC-SA deste livro.

\textbf{Versão 1.1}, novembro de 2026: notação brasileira para as funções trigonométricas
($\sen$, $\tg$, $\cotg$, $\cossec$, $\arcsen$, $\arctg$, $\senh$, $\tgh$\dots) e identidade
visual Nexus BI.

\tableofcontents

\mainmatter
% =====================================================================
\chapter{Funções e gráficos}

O Cálculo é a matemática que descreve como as funções variam. Neste capítulo revisamos
todas as funções necessárias ao estudo do Cálculo: polinomiais, racionais, trigonométricas,
exponenciais, logarítmicas e hiperbólicas. Vemos como avaliá-las, as propriedades dos seus
gráficos e as técnicas algébricas para resolver equações que as envolvem. É a base de tudo o
que vem depois; convém dominá-la antes da introdução formal do Cálculo, no próximo capítulo.

% ---------------------------------------------------------------------
\section{Revisão de funções}

\subsection{Funções}

Dados dois conjuntos $A$ e $B$, um conjunto de pares ordenados $(x,y)$ com $x\in A$ e
$y\in B$ é uma \emph{relação} de $A$ em $B$. Uma função é um tipo especial de relação, em
que cada elemento do primeiro conjunto está associado a \emph{exatamente um} elemento do
segundo. O elemento do primeiro conjunto é a \emph{entrada}; o do segundo, a \emph{saída}.
Quando conhecemos a entrada, a saída fica determinada; por isso dizemos que a saída é
função da entrada. A área de um quadrado, por exemplo, é função do comprimento do lado.

\begin{definicao}
Uma \textbf{função}\index{funcao@função} $f$ consiste em um conjunto de entradas, um conjunto de saídas e uma
regra que associa cada entrada a exatamente uma saída. O conjunto de entradas é o
\textbf{domínio}\index{dominio@domínio} da função; o conjunto de saídas é a \textbf{imagem}\index{imagem}.
\end{definicao}

Para uma função $f$ com domínio $D$, costumamos usar $x$ para a entrada e $y$ para a saída,
e escrevemos $y=f(x)$ (lê-se ``$y$ é igual a $f$ de $x$''). Chamamos $x$ de \emph{variável
independente} e $y$ de \emph{variável dependente}. Uma entrada só pode ser enviada a uma
saída, mas duas entradas diferentes podem ser enviadas à mesma saída
(Figura~\thechapter.1).

\figtikz{\begin{tikzpicture}[font=\small]
\draw[rounded corners=8pt,fill=azul!10,draw=azul] (0,0) rectangle (2,3.6);
\draw[rounded corners=8pt,fill=azul!10,draw=azul] (4.2,0) rectangle (6.2,3.6);
\node at (1,-0.35) {Domínio}; \node at (5.2,-0.35) {Imagem};
\foreach \y/\t in {3/1,2.2/2,1.4/3,0.6/4} \node (a\t) at (1,\y) {\t};
\foreach \y/\t in {2.8/2,1.8/4,0.8/6} \node (b\t) at (5.2,\y) {\t};
\draw[-{Stealth},laranja,thick] (a1) -- (b6);
\draw[-{Stealth},laranja,thick] (a2) -- (b4);
\draw[-{Stealth},azul,thick] (a3) -- (b2);
\draw[-{Stealth},red!70!black,thick] (a4) -- (b2);
\draw[-{Stealth},thick] (1.6,3.9) to[bend left=20] node[above]{$f$} (4.6,3.9);
\end{tikzpicture}}{Uma função associa cada elemento do domínio a exatamente um elemento
da imagem. Entradas diferentes (aqui, 3 e 4) podem ter a mesma saída.}

O \textbf{gráfico}\index{grafico de uma funcao@gráfico de uma função} de $f$ é o conjunto dos pontos $(x,f(x))$ do plano, com $x$ no domínio.
Quando uma função é dada apenas por uma fórmula, sem domínio explícito, o domínio é o
conjunto de todos os reais $x$ para os quais $f(x)$ é um número real; é o chamado
\emph{domínio natural}. Assim, o domínio de $f(x)=x^2$ é $\R$, e o de $f(x)=\sqrt{x}$ é o
conjunto dos reais não negativos.

\textbf{Notação.} Usamos a notação de conjuntos, como $\Set{x}{1<x<5}$, e a de intervalos:
\[
(1,5)=\Set{x}{1<x<5},\quad [1,5]=\Set{x}{1\le x\le 5},\quad
[0,\infty)=\Set{x}{x\ge 0},\quad (-\infty,\infty)=\R .
\]
O símbolo $\infty$ não é um número real; indica apenas que o conjunto não tem limitação
naquela direção.

Algumas funções são \textbf{definidas por partes}\index{funcao@função!definida por partes}, com fórmulas diferentes em diferentes
partes do domínio. Por exemplo,
\[
f(x)=\begin{cases}3x+1, & x\ge 2\\ x^2, & x<2\end{cases}
\]
dá $f(5)=3(5)+1=16$ e $f(-1)=(-1)^2=1$.

\begin{exemplo}{1.1}{Determinando domínio e imagem}
Determine o domínio e a imagem de (a) $f(x)=(x-4)^2+5$; (b) $f(x)=\sqrt{3x+2}-1$;
(c) $f(x)=\dfrac{3}{x-2}$.
\solucao
(a) $f(x)$ é real para todo $x$, logo o domínio é $(-\infty,\infty)$. Como $(x-4)^2\ge0$,
temos $f(x)\ge5$. Para todo $y\ge5$, a equação $(x-4)^2+5=y$ tem solução
$x=4\pm\sqrt{y-5}$; logo a imagem é $\Set{y}{y\ge5}$.

(b) Precisamos de $3x+2\ge0$, ou seja, domínio $\Set{x}{x\ge -2/3}$. Como
$\sqrt{3x+2}\ge0$, temos $f(x)\ge-1$. Dado $y\ge-1$, de $\sqrt{3x+2}=y+1$ vem
$x=\tfrac13(y+1)^2-\tfrac23\ge-\tfrac23$, que está no domínio. A imagem é $\Set{y}{y\ge-1}$.

(c) O denominador não pode se anular: domínio $\Set{x}{x\neq2}$. De $3/(x-2)=y$ vem
$x=3/y+2$, possível sempre que $y\neq0$. A imagem é $\Set{y}{y\neq0}$.
\end{exemplo}

\subsection{Representações de uma função}

Uma função pode ser representada por uma \emph{tabela}, um \emph{gráfico} ou uma
\emph{fórmula}, e muitas vezes usamos as três juntas: marcamos no gráfico os valores de uma
tabela, ou construímos uma tabela a partir de uma fórmula.

Os valores de $x$ com $f(x)=0$ são os \textbf{zeros}\index{zeros de uma funcao@zeros de uma função} da função; eles indicam onde o gráfico
corta o eixo $x$. O \textbf{intercepto $y$}, se existir, é o ponto $(0,f(0))$. Como cada
entrada tem uma única saída, cada reta vertical $x=c$ corta o gráfico no máximo uma vez.

\begin{regra}{Regra: teste da reta vertical}\index{teste da reta vertical}
Dada uma função $f$, toda reta vertical corta o gráfico de $f$ no máximo uma vez. Se alguma
reta vertical corta um conjunto de pontos mais de uma vez, esse conjunto não é gráfico de
uma função.
\end{regra}

\figtikz{\begin{tikzpicture}
\begin{groupplot}[group style={group size=2 by 1,horizontal sep=1.2cm},painel,width=5.6cm,height=4.8cm,xtick=\empty,ytick=\empty]
\nextgroupplot[xmin=-2.4,xmax=2.6,ymin=-2.2,ymax=2.4,title={(a) é função}]
\foreach \a in {-1.4,0.3,1.6}{\edef\tmp{\noexpand\draw[laranja,thin] (axis cs:\a,-2.2) -- (axis cs:\a,2.4);}\tmp}
\addplot[domain=-2.2:2.4]{0.25*x^4-0.15*x^3-1.1*x^2+0.4*x+0.3};
\nextgroupplot[xmin=-2.2,xmax=2.2,ymin=-2.2,ymax=2.4,title={(b) não é função}]
\foreach \a in {-0.6,0.5}{\edef\tmp{\noexpand\draw[laranja,thin] (axis cs:\a,-2.2) -- (axis cs:\a,2.4);}\tmp}
\addplot[variable=t,domain=-1.75:1.75] ({t^3-2*t},{t});
\end{groupplot}\end{tikzpicture}}{(a) Gráfico de uma função: toda reta vertical o
corta no máximo uma vez. (b) Não é gráfico de função: algumas retas verticais o cortam mais
de uma vez.}

Dizemos que $f$ é \textbf{crescente}\index{funcao@função!crescente} em um intervalo $I$ se $f(x_1)\le f(x_2)$ sempre que
$x_1<x_2$ em $I$, e \textbf{estritamente crescente} se $f(x_1)<f(x_2)$. Analogamente, $f$
é \textbf{decrescente}\index{funcao@função!decrescente} se $f(x_1)\ge f(x_2)$ e \textbf{estritamente decrescente} se
$f(x_1)>f(x_2)$ sempre que $x_1<x_2$. Por exemplo, $f(x)=3x$ é crescente em $\R$, e
$f(x)=-x^3$ é estritamente decrescente em $\R$.

\subsection{Combinação e composição de funções}

Dadas $f$ e $g$, definimos a soma, a diferença, o produto e o quociente:
\[
(f\pm g)(x)=f(x)\pm g(x),\qquad (f\cdot g)(x)=f(x)\,g(x),\qquad
\Big(\frac fg\Big)(x)=\frac{f(x)}{g(x)}\ \ \text{se } g(x)\neq0 .
\]
Por exemplo, com $f(x)=2x-3$ e $g(x)=x^2-1$, temos $(f/g)(x)=\dfrac{2x-3}{x^2-1}$, com
domínio $\Set{x}{x\neq\pm1}$.

Também podemos aplicar uma função ao resultado de outra. Se o custo $C$ de aquecer um
prédio depende da temperatura $T$, e a temperatura depende da hora $t$, então o custo como
função do tempo é $C(T(t))$.

\begin{definicao}
Sejam $f$ uma função com domínio $A$ e imagem $B$, e $g$ uma função com domínio $D$ e
imagem $E$. Se $B\subseteq D$, a \textbf{função composta}\index{funcao@função!composta} $g\circ f$ é a função com
domínio $A$ dada por
\[
(g\circ f)(x)=g\big(f(x)\big).
\]
\end{definicao}

\begin{exemplo}{1.2}{Composição de funções dadas por fórmulas}
Sejam $f(x)=x^2+1$ e $g(x)=1/x$. Encontre $g\circ f$ e $f\circ g$, com domínio e imagem.
\solucao
$(g\circ f)(x)=g(x^2+1)=\dfrac{1}{x^2+1}$. Como $x^2+1\neq0$ para todo $x$, o domínio é
$\R$. Como $0<1/(x^2+1)\le1$, e para cada $y\in(0,1]$ a equação tem solução
$x=\pm\sqrt{1/y-1}$, a imagem é $(0,1]$.

$(f\circ g)(x)=f(1/x)=\dfrac{1}{x^2}+1$, com domínio $\Set{x}{x\neq0}$. Resolvendo
$(1/x)^2+1=y$, obtemos $x=\pm1/\sqrt{y-1}$, real se e somente se $y>1$. A imagem é
$\Set{y}{y>1}$.

Note que $f\circ g\neq g\circ f$: a ordem da composição importa.
\end{exemplo}

\subsection{Simetria}

Se $f(-x)=f(x)$, o ponto $(-x,y)$ está no gráfico sempre que $(x,y)$ está, e o gráfico é
simétrico em relação ao eixo $y$. Se $f(-x)=-f(x)$, o gráfico é simétrico em relação à
origem: girá-lo $180^\circ$ em torno da origem não o altera.

\begin{definicao}
Se $f(-x)=f(x)$ para todo $x$ do domínio, $f$ é uma \textbf{função par}\index{funcao@função!par} (simétrica em
relação ao eixo $y$). Se $f(-x)=-f(x)$ para todo $x$ do domínio, $f$ é uma
\textbf{função ímpar}\index{funcao@função!impar@ímpar} (simétrica em relação à origem).
\end{definicao}

\figtikz{\begin{tikzpicture}
\begin{groupplot}[group style={group size=2 by 1,horizontal sep=1.2cm},painel,width=5.6cm,height=4.8cm]
\nextgroupplot[xmin=-2.6,xmax=2.6,ymin=-5,ymax=5,xtick={-2,-1,1,2},ytick={-4,-3,2,4},title={(a) $x^4-2x^2-3$: par}]
\addplot[domain=-2.25:2.25]{x^4-2*x^2-3};
\nextgroupplot[xmin=-3,xmax=3,ymin=-5,ymax=5,xtick={-2,2},ytick={-3,3},title={(b) $x^3-4x$: ímpar}]
\addplot[domain=-2.75:2.75]{x^3-4*x};
\end{groupplot}\end{tikzpicture}}{(a) $f(x)=x^4-2x^2-3$ é par: simetria em relação
ao eixo $y$. (b) $f(x)=x^3-4x$ é ímpar: simetria em relação à origem.}

\begin{exemplo}{1.3}{Funções pares e ímpares}
Classifique: (a) $f(x)=-5x^4+7x^2-2$; (b) $f(x)=2x^5-4x+5$; (c) $f(x)=\dfrac{3x}{x^2+1}$.
\solucao
Calculamos $f(-x)$ e comparamos com $f(x)$ e $-f(x)$.
(a) $f(-x)=-5x^4+7x^2-2=f(x)$: par.
(b) $f(-x)=-2x^5+4x+5$, que não é $f(x)$ nem $-f(x)=-2x^5+4x-5$: nem par nem ímpar.
(c) $f(-x)=\dfrac{-3x}{x^2+1}=-f(x)$: ímpar.
\end{exemplo}

A \textbf{função valor absoluto}\index{valor absoluto} é
\[
|x|=\begin{cases}-x, & x<0\\ x, & x\ge0.\end{cases}
\]
Sua imagem é $\Set{y}{y\ge0}$, e ela é par. Por exemplo, $f(x)=2|x-3|+4$ tem domínio $\R$
e, como $|x-3|\ge0$, imagem $\Set{y}{y\ge4}$.

% ---------------------------------------------------------------------
\section{Classes básicas de funções}

\subsection{Funções afins e inclinação}

Uma \textbf{função afim}\index{funcao@função!afim} (ou linear) tem a forma $f(x)=mx+b$, e seu gráfico é uma reta.
Para quaisquer dois pontos $(x_1,y_1)$ e $(x_2,y_2)$ dessa reta, a razão entre a variação
de $y$ e a de $x$ é sempre a mesma.

\begin{definicao}
A \textbf{inclinação}\index{inclinacao@inclinação} da reta que passa por $(x_1,y_1)$ e $(x_2,y_2)$ é
\[
m=\frac{y_2-y_1}{x_2-x_1}=\frac{\Delta y}{\Delta x}.
\]
A reta que passa por $(x_1,y_1)$ com inclinação $m$ tem \textbf{equação ponto-inclinação}\index{equacao da reta@equação da reta}
$y-y_1=m(x-x_1)$. A reta com inclinação $m$ e intercepto $(0,b)$ tem equação
$y=mx+b$ (forma \textbf{reduzida}). A \textbf{forma geral} $ax+by=c$, com $a$ e $b$ não
ambos nulos, inclui também as retas verticais $x=k$.
\end{definicao}

\begin{exemplo}{1.4}{Inclinação e equação de uma reta}
Encontre a reta que passa por $(11,-4)$ e $(-4,5)$.
\solucao
$m=\dfrac{5-(-4)}{-4-11}=-\dfrac{9}{15}=-\dfrac35$. Pelo ponto $(11,-4)$:
$y+4=-\tfrac35(x-11)$, ou seja, $y=-\tfrac35x+\tfrac{13}{5}$.
\end{exemplo}

\subsection{Polinômios}

Um \textbf{polinômio}\index{polinomio@polinômio} é uma função da forma
$f(x)=a_nx^n+a_{n-1}x^{n-1}+\dots+a_1x+a_0$, com $n$ inteiro não negativo e $a_n\neq0$; $n$
é o \emph{grau}. Polinômios de grau 0, 1, 2 e 3 são chamados constantes, afins,
quadráticos e cúbicos. As \textbf{funções potência}\index{funcao@função!potencia@potência} $f(x)=ax^b$ incluem $x^n$, que é par
para $n$ par e ímpar para $n$ ímpar.

\figtikz{\begin{tikzpicture}
\begin{groupplot}[group style={group size=2 by 1,horizontal sep=1.2cm},painel,width=5.6cm,height=4.8cm,xmin=-1.7,xmax=1.7,xtick={-1,1}]
\nextgroupplot[ymin=-0.4,ymax=3.2,ytick={1,2},title={(a) $n$ par}]
\addplot[laranja,domain=-1.7:1.7]{x^2} node[pos=0.97,left,font=\scriptsize]{$x^2$};
\addplot[domain=-1.33:1.33]{x^4} node[pos=0.98,right,font=\scriptsize]{$x^4$};
\nextgroupplot[ymin=-3.2,ymax=3.2,ytick={-1,1},title={(b) $n$ ímpar}]
\addplot[laranja,domain=-1.47:1.47]{x^3} node[pos=0.98,right,font=\scriptsize]{$x^3$};
\addplot[domain=-1.26:1.26]{x^5} node[pos=0.98,left,font=\scriptsize]{$x^5$};
\end{groupplot}\end{tikzpicture}}{(a) Para $n$ par, $ax^n$ é par. (b) Para $n$
ímpar, $ax^n$ é ímpar.}

\textbf{Comportamento no infinito.} Para $|x|$ grande, o termo de maior grau domina. Se
$a_n>0$ e $n$ é par, $f(x)\to\infty$ quando $x\to\pm\infty$; se $n$ é ímpar,
$f(x)\to\infty$ quando $x\to\infty$ e $f(x)\to-\infty$ quando $x\to-\infty$. Se $a_n<0$,
os sinais se invertem.

\textbf{Zeros.} Os zeros de polinômios de grau 1 e 2 se encontram diretamente; para grau
maior, fatoramos quando possível.

\begin{regra}{Regra: fórmula quadrática}\index{formula quadratica@fórmula quadrática}
As soluções de $ax^2+bx+c=0$, com $a\neq0$, são
\[
x=\frac{-b\pm\sqrt{b^2-4ac}}{2a}.
\]
Se o discriminante $b^2-4ac$ é positivo, há duas soluções reais; se é nulo, uma; se é
negativo, nenhuma.
\end{regra}

\subsection{Funções algébricas e transcendentes}

Uma \textbf{função racional}\index{funcao@função!racional} é um quociente de polinômios, $f(x)=p(x)/q(x)$, definida onde
$q(x)\neq0$. As \textbf{funções algébricas}\index{funcao@função!algebrica@algébrica} são as obtidas de polinômios por soma,
subtração, multiplicação, divisão e extração de raízes; por exemplo,
$f(x)=\sqrt{x^2-4}$ e $g(x)=\sqrt[3]{x}/(1-x^2)$. Para $f(x)=\sqrt[n]{x}$, o domínio é
$[0,\infty)$ se $n$ é par e $\R$ se $n$ é ímpar. As funções que não são algébricas são as
\textbf{transcendentes}\index{funcao@função!transcendente}: trigonométricas, exponenciais, logarítmicas e hiperbólicas.

\subsection{Funções definidas por partes}

\begin{exemplo}{1.5}{Gráfico de uma função definida por partes}
Esboce $f(x)=\begin{cases}x+3, & x<1\\ (x-2)^2, & x\ge1.\end{cases}$
\solucao
Desenhamos a reta $y=x+3$ em $(-\infty,1)$ e a parábola $y=(x-2)^2$ em $[1,\infty)$. Como
$f(1)=(1-2)^2=1$, marcamos uma bola cheia em $(1,1)$. A reta se aproxima de $(1,4)$ sem
atingi-lo: marcamos ali uma bola vazia.
\figtikz{\begin{tikzpicture}
\begin{axis}[painel,width=6.2cm,height=5cm,xmin=-4.5,xmax=4.8,ymin=-1.8,ymax=6,xtick={-3,1,2,4},ytick={1,4}]
\addplot[domain=-4.5:1]{x+3} node[pos=0.25,above left,font=\scriptsize]{$x+3$};
\addplot[laranja,domain=1:4.4]{(x-2)^2} node[pos=0.9,right,font=\scriptsize]{$(x-2)^2$};
\draw[azul,fill=white] (axis cs:1,4) circle (2pt); \fill[laranja] (axis cs:1,1) circle (2pt);
\end{axis}\end{tikzpicture}}{Função afim para $x<1$ e quadrática para $x\ge1$.}
\end{exemplo}

\subsection{Transformações de funções}\index{transformacoes de graficos@transformações de gráficos}

A partir do gráfico de $y=f(x)$ obtemos outros sem novos cálculos. Para $c>0$:

\begin{center}
\begin{tabular}{@{}ll@{}}\toprule
\textbf{Função} & \textbf{Efeito sobre o gráfico de $f$}\\\midrule
$f(x)+c$ / $f(x)-c$ & translação vertical de $c$ para cima / para baixo\\
$f(x+c)$ / $f(x-c)$ & translação horizontal de $c$ para a esquerda / para a direita\\
$c\,f(x)$ & alongamento vertical se $c>1$; compressão vertical se $0<c<1$\\
$f(cx)$ & compressão horizontal se $c>1$; alongamento horizontal se $0<c<1$\\
$-f(x)$ & reflexão em relação ao eixo $x$\\
$f(-x)$ & reflexão em relação ao eixo $y$\\\bottomrule
\end{tabular}
\end{center}

\figtikz{\begin{tikzpicture}
\begin{groupplot}[group style={group size=2 by 1,horizontal sep=1.4cm},painel,width=6.2cm,height=5cm]
\nextgroupplot[xmin=-4,xmax=4,ymin=-5.5,ymax=9,title={Translações verticais}]
\addplot[cinza,dashed,domain=-4:4]{x^2};
\addplot[domain=-4:4]{x^2+4} node[pos=0.85,right,font=\scriptsize]{$x^2+4$};
\addplot[laranja,domain=-4:4]{x^2-4} node[pos=0.93,right,font=\scriptsize]{$x^2-4$};
\nextgroupplot[xmin=-7,xmax=7,ymin=-1,ymax=6,title={Translações horizontais}]
\addplot[cinza,dashed,domain=-7:7]{abs(x)};
\addplot[domain=-7:7]{abs(x+3)} node[pos=0.04,above right,font=\scriptsize]{$|x+3|$};
\addplot[laranja,domain=-7:7]{abs(x-3)} node[pos=0.96,above left,font=\scriptsize]{$|x-3|$};
\end{groupplot}\end{tikzpicture}}{Translações dos gráficos de $x^2$ e $|x|$ (tracejados).
À esquerda, verticais: $x^2\pm4$. À direita, horizontais: $|x+3|$ desloca para a esquerda
e $|x-3|$, para a direita.}

\begin{exemplo}{1.6}{Transformando uma função}
Esboce (a) $f(x)=-|x+2|-3$ e (b) $f(x)=3\sqrt{-x}+1$.
\solucao
(a) Partindo de $y=|x|$: translação de 2 unidades para a esquerda, reflexão em relação ao
eixo $x$ e translação de 3 unidades para baixo.
(b) Partindo de $y=\sqrt{x}$: reflexão em relação ao eixo $y$, alongamento vertical por um
fator 3 e translação de 1 unidade para cima.
\figtikz{\begin{tikzpicture}
\begin{groupplot}[group style={group size=2 by 1,horizontal sep=1.2cm},painel,width=5.8cm,height=5cm]
\nextgroupplot[xmin=-6.5,xmax=4,ymin=-7.5,ymax=4,xtick={-2,2},ytick={-3,3},title={(a) $-|x+2|-3$}]
\addplot[cinza,dashed,domain=-4:4]{abs(x)} node[pos=0.95,left,font=\scriptsize]{$|x|$};
\addplot[domain=-6.5:2.5]{-abs(x+2)-3};
\fill[azul] (axis cs:-2,-3) circle (1.6pt);
\nextgroupplot[xmin=-6.5,xmax=4.5,ymin=-0.8,ymax=9,xtick={-4,-1,1,4},ytick={1,4,7},title={(b) $3\sqrt{-x}+1$}]
\addplot[cinza,dashed,domain=0:4.5,samples=100]{sqrt(x)} node[pos=0.9,above,font=\scriptsize]{$\sqrt x$};
\addplot[domain=-6.5:0,samples=100]{3*sqrt(-x)+1};
\fill[azul] (axis cs:0,1) circle (1.6pt); \fill[azul] (axis cs:-4,7) circle (1.6pt);
\end{groupplot}\end{tikzpicture}}{As funções (a)
e (b) como sequências de três transformações de $|x|$ e de $\sqrt x$.}
\end{exemplo}

% ---------------------------------------------------------------------
\section{Funções trigonométricas}\index{funcoes trigonometricas@funções trigonométricas}

\subsection{Radianos}

No Cálculo, os ângulos são medidos em \textbf{radianos}\index{radiano}. A medida em radianos de um ângulo
$\theta$ é o comprimento $s$ do arco que ele determina no círculo unitário. Como a
circunferência unitária mede $2\pi$, temos $2\pi$ rad $=360^\circ$, ou seja,
\[
1\ \text{rad}=\frac{180^\circ}{\pi}\approx57{,}3^\circ,\qquad 1^\circ=\frac{\pi}{180}\ \text{rad}.
\]
Assim, $225^\circ=225\cdot\pi/180=5\pi/4$ rad, e $5\pi/3$ rad $=300^\circ$.

\subsection{As seis funções trigonométricas}

\begin{definicao}
Seja $P=(x,y)$ um ponto do círculo unitário centrado na origem $O$, e $\theta$ o ângulo
com lado inicial no semieixo $x$ positivo e lado final $OP$. Então
\[
\sen\theta=y,\quad \cos\theta=x,\quad \tg\theta=\frac yx,\quad
\cossec\theta=\frac1y,\quad \sec\theta=\frac1x,\quad \cotg\theta=\frac xy .
\]
Se $x=0$, $\sec\theta$ e $\tg\theta$ não estão definidas; se $y=0$, $\cotg\theta$ e
$\cossec\theta$ não estão definidas.
\end{definicao}

\figtikz{\begin{tikzpicture}[scale=1.9,font=\small]
\draw[->] (-1.3,0) -- (1.4,0) node[right]{$x$}; \draw[->] (0,-1.3) -- (0,1.4) node[above]{$y$};
\draw[thick] (0,0) circle (1);
\draw[thick,azul] (0,0) -- node[above left,pos=0.55]{$1$} (60:1); \fill[azul] (60:1) circle (1.3pt) node[above right]{$P=(\cos\theta,\sen\theta)$};
\draw[laranja,dashed] (60:1) -- (0.5,0) node[below]{$\cos\theta$};
\draw[->] (0.3,0) arc (0:60:0.3); \node at (25:0.45) {$\theta$};
\draw[laranja,thin] (0.5,0.45) -- (1.08,0.62) node[right,font=\footnotesize]{$\sen\theta$};
\end{tikzpicture}}{O ângulo $\theta$ na posição padrão: o ponto $P$ do círculo unitário tem coordenadas $(\cos\theta,\sen\theta)$.}

Em um círculo de raio $r$, as coordenadas de um ponto satisfazem $x=r\cos\theta$ e
$y=r\sen\theta$. Em um triângulo retângulo com ângulo agudo $\theta$, isso dá as razões
conhecidas: $\sen\theta=$ cateto oposto/hipotenusa, $\cos\theta=$ cateto
adjacente/hipotenusa e $\tg\theta=$ oposto/adjacente.

\begin{center}
\renewcommand{\arraystretch}{1.3}
\begin{tabular}{@{}c|ccccc@{}}\toprule
$\theta$ & $0$ & $\pi/6$ & $\pi/4$ & $\pi/3$ & $\pi/2$\\\midrule
$\sen\theta$ & $0$ & $1/2$ & $\sqrt2/2$ & $\sqrt3/2$ & $1$\\
$\cos\theta$ & $1$ & $\sqrt3/2$ & $\sqrt2/2$ & $1/2$ & $0$\\
$\tg\theta$ & $0$ & $\sqrt3/3$ & $1$ & $\sqrt3$ & não definida\\\bottomrule
\end{tabular}
\end{center}

\subsection{Identidades trigonométricas}

\begin{regra}{Regra: identidades trigonométricas}\index{identidades trigonometricas@identidades trigonométricas}
\textbf{Recíprocas:} $\cotg\theta=\dfrac1{\tg\theta}$,\; $\cossec\theta=\dfrac1{\sen\theta}$,\;
$\sec\theta=\dfrac1{\cos\theta}$.

\textbf{De razão:} $\tg\theta=\dfrac{\sen\theta}{\cos\theta}$,\;
$\cotg\theta=\dfrac{\cos\theta}{\sen\theta}$.

\textbf{Pitagóricas:} $\sen^2\theta+\cos^2\theta=1$,\quad $1+\tg^2\theta=\sec^2\theta$,\quad
$1+\cotg^2\theta=\cossec^2\theta$.

\textbf{Adição e subtração:}
$\sen(\alpha\pm\beta)=\sen\alpha\cos\beta\pm\cos\alpha\sen\beta$,\quad
$\cos(\alpha\pm\beta)=\cos\alpha\cos\beta\mp\sen\alpha\sen\beta$.

\textbf{Ângulo duplo:} $\sen(2\theta)=2\sen\theta\cos\theta$,\quad
$\cos(2\theta)=2\cos^2\theta-1=1-2\sen^2\theta=\cos^2\theta-\sen^2\theta$.
\end{regra}

\begin{exemplo}{1.7}{Resolvendo equações trigonométricas}
Encontre todas as soluções de (a) $1+\cos(2\theta)=\cos\theta$ e (b) $\sen(2\theta)=\tg\theta$.
\solucao
(a) Pela fórmula do ângulo duplo, a equação equivale a $2\cos^2\theta-\cos\theta=0$.
Devemos \emph{fatorar}, e não dividir por $\cos\theta$, que pode ser zero:
$\cos\theta\,(2\cos\theta-1)=0$. Logo $\cos\theta=0$ ou $\cos\theta=1/2$, e as soluções são
\[
\theta=\frac\pi2+n\pi,\qquad \theta=\pm\frac\pi3+2n\pi,\qquad n\in\mathbb Z.
\]
(b) Escrevendo $2\sen\theta\cos\theta=\sen\theta/\cos\theta$ e multiplicando por
$\cos\theta$, obtemos $\sen\theta\,(2\cos^2\theta-1)=0$. Como multiplicamos por
$\cos\theta$, que poderia ser nulo, é preciso descartar soluções estranhas no fim. Temos
$\sen\theta=0$ ou $\cos^2\theta=1/2$, e as soluções são
$\theta=n\pi$ e $\theta=\dfrac\pi4+\dfrac{n\pi}2$, $n\in\mathbb Z$.
\end{exemplo}

\subsection{Gráficos e períodos}

Uma função $f$ é \textbf{periódica}\index{funcao@função!periodica@periódica} se existe $P>0$ tal que $f(x+P)=f(x)$ para todo $x$; o
menor desses $P$ é o \textbf{período}. Seno, cosseno, secante e cossecante têm período
$2\pi$; tangente e cotangente têm período $\pi$.

\figtikz{\begin{tikzpicture}
\begin{groupplot}[group style={group size=2 by 3,horizontal sep=1.3cm,vertical sep=1.3cm},painel,width=6cm,height=3.6cm]
\nextgroupplot[title={$y=\sen x$},xmin=-6.6,xmax=6.6,ymin=-1.6,ymax=1.6,xtick={-6.2832,-3.1416,3.1416,6.2832},xticklabels={$-2\pi$,$-\pi$,$\pi$,$2\pi$}]
\addplot[domain=-6.6:6.6]{sin(deg(x))};\nextgroupplot[title={$y=\cos x$},xmin=-6.6,xmax=6.6,ymin=-1.6,ymax=1.6,xtick={-6.2832,-3.1416,3.1416,6.2832},xticklabels={$-2\pi$,$-\pi$,$\pi$,$2\pi$}]
\addplot[domain=-6.6:6.6]{cos(deg(x))};\nextgroupplot[title={$y=\tg x$},xmin=-6.6,xmax=6.6,ymin=-4,ymax=4,xtick={-6.2832,-3.1416,3.1416,6.2832},xticklabels={$-2\pi$,$-\pi$,$\pi$,$2\pi$}]
\addplot[domain=-6.560:-4.752,restrict y to domain=-6:6]{tan(deg(x))};
\addplot[domain=-4.672:-1.611,restrict y to domain=-6:6]{tan(deg(x))};
\addplot[domain=-1.531:1.531,restrict y to domain=-6:6]{tan(deg(x))};
\addplot[domain=1.611:4.672,restrict y to domain=-6:6]{tan(deg(x))};
\addplot[domain=4.752:6.560,restrict y to domain=-6:6]{tan(deg(x))};
\draw[assint] (axis cs:-4.7124,-4) -- (axis cs:-4.7124,4);
\draw[assint] (axis cs:-1.5708,-4) -- (axis cs:-1.5708,4);
\draw[assint] (axis cs:1.5708,-4) -- (axis cs:1.5708,4);
\draw[assint] (axis cs:4.7124,-4) -- (axis cs:4.7124,4);
\nextgroupplot[title={$y=\cotg x$},xmin=-6.6,xmax=6.6,ymin=-4,ymax=4,xtick={-6.2832,-3.1416,3.1416,6.2832},xticklabels={$-2\pi$,$-\pi$,$\pi$,$2\pi$}]
\addplot[domain=-6.560:-6.323,restrict y to domain=-6:6]{cos(deg(x))/sin(deg(x))};
\addplot[domain=-6.243:-3.182,restrict y to domain=-6:6]{cos(deg(x))/sin(deg(x))};
\addplot[domain=-3.102:-0.040,restrict y to domain=-6:6]{cos(deg(x))/sin(deg(x))};
\addplot[domain=0.040:3.102,restrict y to domain=-6:6]{cos(deg(x))/sin(deg(x))};
\addplot[domain=3.182:6.243,restrict y to domain=-6:6]{cos(deg(x))/sin(deg(x))};
\addplot[domain=6.323:6.560,restrict y to domain=-6:6]{cos(deg(x))/sin(deg(x))};
\draw[assint] (axis cs:-6.2832,-4) -- (axis cs:-6.2832,4);
\draw[assint] (axis cs:-3.1416,-4) -- (axis cs:-3.1416,4);
\draw[assint] (axis cs:0.0000,-4) -- (axis cs:0.0000,4);
\draw[assint] (axis cs:3.1416,-4) -- (axis cs:3.1416,4);
\draw[assint] (axis cs:6.2832,-4) -- (axis cs:6.2832,4);
\nextgroupplot[title={$y=\sec x$},xmin=-6.6,xmax=6.6,ymin=-4,ymax=4,xtick={-6.2832,-3.1416,3.1416,6.2832},xticklabels={$-2\pi$,$-\pi$,$\pi$,$2\pi$}]
\addplot[domain=-6.560:-4.752,restrict y to domain=-6:6]{1/cos(deg(x))};
\addplot[domain=-4.672:-1.611,restrict y to domain=-6:6]{1/cos(deg(x))};
\addplot[domain=-1.531:1.531,restrict y to domain=-6:6]{1/cos(deg(x))};
\addplot[domain=1.611:4.672,restrict y to domain=-6:6]{1/cos(deg(x))};
\addplot[domain=4.752:6.560,restrict y to domain=-6:6]{1/cos(deg(x))};
\draw[assint] (axis cs:-4.7124,-4) -- (axis cs:-4.7124,4);
\draw[assint] (axis cs:-1.5708,-4) -- (axis cs:-1.5708,4);
\draw[assint] (axis cs:1.5708,-4) -- (axis cs:1.5708,4);
\draw[assint] (axis cs:4.7124,-4) -- (axis cs:4.7124,4);
\nextgroupplot[title={$y=\cossec x$},xmin=-6.6,xmax=6.6,ymin=-4,ymax=4,xtick={-6.2832,-3.1416,3.1416,6.2832},xticklabels={$-2\pi$,$-\pi$,$\pi$,$2\pi$}]
\addplot[domain=-6.560:-6.323,restrict y to domain=-6:6]{1/sin(deg(x))};
\addplot[domain=-6.243:-3.182,restrict y to domain=-6:6]{1/sin(deg(x))};
\addplot[domain=-3.102:-0.040,restrict y to domain=-6:6]{1/sin(deg(x))};
\addplot[domain=0.040:3.102,restrict y to domain=-6:6]{1/sin(deg(x))};
\addplot[domain=3.182:6.243,restrict y to domain=-6:6]{1/sin(deg(x))};
\addplot[domain=6.323:6.560,restrict y to domain=-6:6]{1/sin(deg(x))};
\draw[assint] (axis cs:-6.2832,-4) -- (axis cs:-6.2832,4);
\draw[assint] (axis cs:-3.1416,-4) -- (axis cs:-3.1416,4);
\draw[assint] (axis cs:0.0000,-4) -- (axis cs:0.0000,4);
\draw[assint] (axis cs:3.1416,-4) -- (axis cs:3.1416,4);
\draw[assint] (axis cs:6.2832,-4) -- (axis cs:6.2832,4);
\end{groupplot}\end{tikzpicture}}{As seis funções trigonométricas. Seno, cosseno, secante e
cossecante têm período $2\pi$; tangente e cotangente, período $\pi$. Tracejado: assíntotas
verticais.}

Para $f(x)=A\sen\big(B(x-\alpha)\big)+C$ (ou com cosseno), $|A|$ é a \textbf{amplitude}\index{amplitude},
o período é $2\pi/|B|$, $\alpha$ é a translação horizontal (\emph{defasagem}) e $C$ é a
translação vertical.

\begin{exemplo}{1.8}{Esboçando uma senoide transformada}
Esboce $f(x)=3\sen\!\big(2(x-\tfrac\pi4)\big)+1$.
\solucao
Compressão horizontal por fator 2, defasagem de $\pi/4$ para a direita, alongamento
vertical por fator 3 e translação de 1 unidade para cima. Amplitude 3, período $\pi$.
\figtikz{\begin{tikzpicture}
\begin{axis}[painel,width=10cm,height=4.8cm,xmin=-4.9,xmax=7.2,ymin=-2.6,ymax=4.8,
  xtick={-4.712,-3.1416,-1.5708,0.7854,1.5708,3.1416,4.712,6.2832},xticklabels={$-\frac{3\pi}2$,$-\pi$,$-\frac\pi2$,$\frac\pi4$,$\frac\pi2$,$\pi$,$\frac{3\pi}2$,$2\pi$},ytick={-2,1,4}]
\draw[assint] (axis cs:-4.9,1) -- (axis cs:7.2,1);
\addplot[domain=-4.9:7.2,samples=300]{3*sin(deg(2*(x-pi/4)))+1};
\end{axis}\end{tikzpicture}}{Gráfico de $f(x)=3\sen(2(x-\pi/4))+1$: amplitude $3$, período $\pi$, deslocamento de fase $\frac\pi4$ e linha média $y=1$ (tracejada).}
\end{exemplo}

% ---------------------------------------------------------------------
\section{Funções inversas}

\subsection{Existência da inversa}

Uma função ``desfaz'' a outra se, ao aplicarmos uma e depois a outra, voltamos à entrada.

\begin{definicao}
Dada $f$ com domínio $D$ e imagem $R$, sua \textbf{função inversa}\index{funcao@função!inversa} (se existir) é a
função $f^{-1}$ com domínio $R$ e imagem $D$ tal que $f^{-1}(y)=x$ se $f(x)=y$. Ou seja,
\[
f^{-1}\big(f(x)\big)=x\ \text{ para todo } x\in D,\qquad
f\big(f^{-1}(y)\big)=y\ \text{ para todo } y\in R.
\]
\end{definicao}

\begin{nota}
$f^{-1}(x)$ não significa $1/f(x)$. A notação indica a função inversa.
\end{nota}

Para que $f^{-1}$ seja função, cada saída de $f$ deve vir de uma única entrada.

\begin{definicao}
$f$ é \textbf{injetora}\index{funcao@função!injetora} se $f(x_1)\neq f(x_2)$ sempre que $x_1\neq x_2$. Uma função tem
inversa se, e somente se, é injetora.
\end{definicao}

\begin{regra}{Regra: teste da reta horizontal}\index{teste da reta horizontal}
Uma função $f$ é injetora se, e somente se, toda reta horizontal corta o gráfico de $f$ no
máximo uma vez.
\end{regra}

\figtikz{\begin{tikzpicture}
\begin{groupplot}[group style={group size=2 by 1,horizontal sep=1.2cm},painel,width=5.6cm,height=4.8cm,xmin=-2.4,xmax=2.4,xtick={-1,1},ytick={2}]
\nextgroupplot[ymin=-0.5,ymax=4.4,title={(a) $x^2$}]
\draw[laranja] (axis cs:-2.4,2) -- (axis cs:2.4,2);
\addplot[domain=-2.1:2.1]{x^2};
\fill (axis cs:-1.414,2) circle (1.6pt); \fill (axis cs:1.414,2) circle (1.6pt);
\nextgroupplot[ymin=-4,ymax=4,title={(b) $x^3$}]
\draw[laranja] (axis cs:-2.4,2) -- (axis cs:2.4,2);
\addplot[domain=-1.6:1.6]{x^3};
\fill (axis cs:1.26,2) circle (1.6pt);
\end{groupplot}\end{tikzpicture}}{(a) $x^2$ não é injetora: falha no teste da reta
horizontal. (b) $x^3$ é injetora.}

\begin{estrategia}{encontrando a inversa}
\begin{enumerate}[label=\arabic*.]
\item Resolva $y=f(x)$ para $x$, em função de $y$.
\item Troque as variáveis e escreva $y=f^{-1}(x)$.
\end{enumerate}
\end{estrategia}

\begin{exemplo}{1.9}{Encontrando uma inversa}
Encontre a inversa de $f(x)=3x-4$.
\solucao
De $y=3x-4$ vem $x=\tfrac13y+\tfrac43$; logo $f^{-1}(x)=\tfrac13x+\tfrac43$, com domínio e
imagem $\R$. Verificação:
$f^{-1}(f(x))=\tfrac13(3x-4)+\tfrac43=x$.
\end{exemplo}

\textbf{Gráfico da inversa.} Se $(a,b)$ está no gráfico de $f$, então $(b,a)$ está no de
$f^{-1}$. Portanto o gráfico de $f^{-1}$ é a reflexão do gráfico de $f$ em relação à reta
$y=x$.

\figtikz{\begin{tikzpicture}
\begin{axis}[painel,width=6.2cm,height=6.2cm,xmin=-2.6,xmax=3,ymin=-2.6,ymax=3,xtick=\empty,ytick=\empty]
\addplot[cinza,dashed,domain=-2.6:3]{x} node[pos=0.97,below right,font=\scriptsize]{$y=x$};
\addplot[domain=-2:1.9,samples=100] ({x},{x^3/4+x/2}) node[pos=0.97,left,font=\scriptsize]{$f$};
\addplot[laranja,domain=-2:1.9,samples=100] ({x^3/4+x/2},{x}) node[pos=0.97,below,font=\scriptsize]{$f^{-1}$};
\fill[azul] (axis cs:1.2,1.032) circle (1.6pt) node[right,font=\scriptsize]{$(a,b)$};
\fill[laranja] (axis cs:1.032,1.2) circle (1.6pt) node[above left,font=\scriptsize]{$(b,a)$};
\draw[cinza,thin,dotted] (axis cs:1.2,1.032) -- (axis cs:1.032,1.2);
\end{axis}\end{tikzpicture}}{Se $(a,b)$ está no gráfico de $f$, então $(b,a)$ está no gráfico de $f^{-1}$: os gráficos são simétricos em relação à reta $y=x$.}

\textbf{Restrição de domínio.} Uma função que não é injetora pode ter inversa em um domínio
menor. Assim, $x^2$ restrita a $[0,\infty)$ tem inversa $\sqrt x$.

\begin{exemplo}{1.10}{Restringindo o domínio}
$f(x)=(x+1)^2$ não é injetora (é $x^2$ transladada 1 unidade para a esquerda). Em
$[-1,\infty)$ ela é injetora, com imagem $[0,\infty)$. Resolvendo $y=(x+1)^2$, vem
$x=-1\pm\sqrt y$; como $x\ge-1$, escolhemos o sinal $+$. Logo $f^{-1}(x)=-1+\sqrt x$, com
domínio $[0,\infty)$ e imagem $[-1,\infty)$.
\end{exemplo}

\subsection{Funções trigonométricas inversas}\index{funcoes trigonometricas inversas@funções trigonométricas inversas}

As funções trigonométricas não são injetoras, pois são periódicas. Restringindo o domínio,
obtemos as inversas.

\begin{definicao}
Para $-1\le x\le1$:
\[
\arcsen x=y \iff \sen y=x \text{ e } -\tfrac\pi2\le y\le\tfrac\pi2;\qquad
\arccos x=y \iff \cos y=x \text{ e } 0\le y\le\pi.
\]
Para todo $x$ real:
\[
\arctg x=y \iff \tg y=x \text{ e } -\tfrac\pi2<y<\tfrac\pi2;\qquad
\arccotg x=y \iff \cotg y=x \text{ e } 0<y<\pi.
\]
Para $|x|\ge1$, definem-se de modo análogo $\arccossec x$ (com
$-\tfrac\pi2\le y\le\tfrac\pi2$, $y\neq0$) e $\arcsec x$ (com $0\le y\le\pi$,
$y\neq\tfrac\pi2$). Em calculadoras e em livros em inglês aparecem as notações $\sin^{-1}$, $\cos^{-1}$, $\tan^{-1}$ (ou arcsin, arctan), com o mesmo significado.
\end{definicao}

\figtikz{\begin{tikzpicture}
\begin{groupplot}[group style={group size=3 by 1,horizontal sep=1.1cm},painel,width=5.2cm,height=5.2cm,
  xmin=-2,xmax=3.4,ymin=-2,ymax=3.4,
  xtick={-1,1},ytick={-1.5708,1.5708,3.1416},yticklabels={$-\frac\pi2$,$\frac\pi2$,$\pi$}]
\nextgroupplot[title={$y=\arcsen x$}]
\addplot[cinza,dotted,thin,domain=-2:3.4]{x};
\addplot[laranja!70,thin,domain=-1.5708:1.5708]{sin(deg(x))};
\addplot[domain=-1:1]{rad(asin(x))};
\nextgroupplot[title={$y=\arccos x$}]
\addplot[cinza,dotted,thin,domain=-2:3.4]{x};
\addplot[laranja!70,thin,domain=0:3.1416]{cos(deg(x))};
\addplot[domain=-1:1]{rad(acos(x))};
\nextgroupplot[title={$y=\arctg x$},xmin=-3.4,xtick={-2,2}]
\addplot[cinza,dotted,thin,domain=-3.4:3.4]{x};
\addplot[laranja!70,thin,domain=-1.25:1.25]{tan(deg(x))};
\addplot[domain=-3.4:3.4]{rad(atan(x))};
\draw[assint] (axis cs:-3.4,1.5708) -- (axis cs:3.4,1.5708);
\draw[assint] (axis cs:-3.4,-1.5708) -- (axis cs:3.4,-1.5708);
\end{groupplot}\end{tikzpicture}}{Cada inversa trigonométrica (em grafite) é a reflexão, em
relação à reta $y=x$ (pontilhada), da função trigonométrica restrita (em verde):
$\sen$ em $[-\frac\pi2,\frac\pi2]$, $\cos$ em $[0,\pi]$ e $\tg$ em
$(-\frac\pi2,\frac\pi2)$. As demais inversas são obtidas do mesmo modo.}

\begin{exemplo}{1.11}{Avaliando inversas trigonométricas}
(a) $\arcsen(-\sqrt3/2)$: o ângulo em $[-\pi/2,\pi/2]$ com seno $-\sqrt3/2$ é $-\pi/3$.

(b) $\arccos(\cos(5\pi/4))=\arccos(-\sqrt2/2)=3\pi/4$, e \emph{não} $5\pi/4$, pois
$\arccos$ toma valores em $[0,\pi]$.

(c) $\arcsen(\cos(2\pi/3))=\arcsen(-1/2)=-\pi/6$.
\end{exemplo}

% ---------------------------------------------------------------------
\section{Funções exponenciais e logarítmicas}

\subsection{Funções exponenciais}

Se uma população começa com $P_0$ indivíduos e cresce à taxa anual $r$, após $t$ anos ela
tem $P(t)=P_0(1+r)^t$ indivíduos. Esse é um exemplo de \textbf{função exponencial}\index{funcao@função!exponencial}: uma
função da forma $f(x)=b^x$, com $b>0$ e $b\neq1$. Seu domínio é $\R$ e sua imagem é
$(0,\infty)$. Se $b>1$, $b^x$ é crescente; se $0<b<1$, é decrescente. A exponencial
cresce muito mais rapidamente que qualquer potência: $2^x$ supera $x^2$ para $x$ grande.

\figtikz{\begin{tikzpicture}
\begin{axis}[painel,width=7cm,height=5cm,xmin=-3,xmax=3,ymin=-0.3,ymax=6.2,xtick={-1,1},ytick={1,2,4}]
\addplot[domain=-3:1.29]{4^x}; \addplot[domain=-3:2.58]{2^x};
\node[azul,font=\scriptsize,anchor=west] at (axis cs:1.35,5.5) {$4^x$};
\node[azul,font=\scriptsize,anchor=west] at (axis cs:2.2,4.2) {$2^x$};
\addplot[laranja,domain=-1.29:3]{(1/4)^x}; \addplot[laranja,domain=-2.58:3]{(1/2)^x};
\node[laranja,font=\scriptsize,anchor=east] at (axis cs:-1.35,5.5) {$(\frac14)^x$};
\node[laranja,font=\scriptsize,anchor=east] at (axis cs:-2.2,4.2) {$(\frac12)^x$};
\fill (axis cs:0,1) circle (1.5pt);
\end{axis}\end{tikzpicture}}{$b^x$ é crescente se $b>1$ (grafite) e decrescente se $0<b<1$ (verde). Todas passam por $(0,1)$.}

\begin{regra}{Regra: leis dos expoentes}\index{expoentes, leis dos}
Para $a>0$, $b>0$ e quaisquer $x$, $y$:
\[
b^xb^y=b^{x+y},\quad \frac{b^x}{b^y}=b^{x-y},\quad (b^x)^y=b^{xy},\quad
(ab)^x=a^xb^x,\quad \frac{a^x}{b^x}=\Big(\frac ab\Big)^x.
\]
\end{regra}

\subsection{O número $e$}

Se investimos $P$ reais a juros anuais $r$ capitalizados $n$ vezes por ano, após $t$ anos
temos $P(1+r/n)^{nt}$. Quando $n$ cresce (capitalização cada vez mais frequente), a
quantidade $(1+1/m)^m$ se aproxima de um número fixo,
\[
e\approx2{,}718282,
\]
e o montante tende a $Pe^{rt}$ (\emph{capitalização contínua}). O número $e$ é irracional.
A função $f(x)=e^x$ é a \textbf{exponencial natural}\index{exponencial natural}. Ela é a única exponencial cujo
gráfico tem reta tangente de inclinação exatamente 1 em $x=0$, fato que tornará as
derivadas e integrais de $e^x$ especialmente simples.

\subsection{Funções logarítmicas}

Como $b^x$ é injetora, ela tem inversa: o \textbf{logaritmo na base $b$},
\[
\log_b x=y \iff b^y=x,\qquad x>0 .
\]
Seu domínio é $(0,\infty)$ e sua imagem é $\R$. Escrevemos $\ln x=\log_e x$
(\textbf{logaritmo natural}\index{logaritmo natural}). Como inversas, $b^{\log_b x}=x$ para $x>0$ e
$\log_b(b^x)=x$ para todo $x$; em particular, $e^{\ln x}=x$ e $\ln(e^x)=x$.

\figtikz{\begin{tikzpicture}
\begin{axis}[painel,width=6.2cm,height=6.2cm,xmin=-3,xmax=4,ymin=-3,ymax=4,xtick={1},ytick={1}]
\addplot[cinza,dashed,domain=-3:4]{x} node[pos=0.97,below right,font=\scriptsize]{$y=x$};
\addplot[domain=-3:1.38]{exp(x)} node[pos=0.95,left,font=\scriptsize]{$e^x$};
\addplot[laranja,domain=0.05:4,samples=150]{ln(x)} node[pos=0.95,below,font=\scriptsize]{$\ln x$};
\end{axis}\end{tikzpicture}}{$y=e^x$ e $y=\ln x$ são inversas: seus gráficos
são simétricos em relação a $y=x$.}

\begin{regra}{Regra: propriedades dos logaritmos}\index{logaritmo!propriedades}
Se $a,b,c>0$, $b\neq1$ e $r\in\R$:
\[
\log_b(ac)=\log_b a+\log_b c,\qquad
\log_b\Big(\frac ac\Big)=\log_b a-\log_b c,\qquad
\log_b(a^r)=r\log_b a .
\]
\end{regra}

\begin{exemplo}{1.12}{Equações exponenciais e logarítmicas}
(a) $5^x=2$. Aplicando $\ln$: $x\ln5=\ln2$, logo $x=\ln2/\ln5$.

(b) $e^x+6e^{-x}=5$. Multiplicando por $e^x$: $(e^x)^2-5e^x+6=0$, isto é,
$(e^x-3)(e^x-2)=0$. Logo $x=\ln3$ ou $x=\ln2$.

(c) $\log_{10}\sqrt x+\log_{10}x=2$. O lado esquerdo é
$\log_{10}x^{3/2}=\tfrac32\log_{10}x$; logo $\log_{10}x=4/3$ e $x=10^{4/3}=10\sqrt[3]{10}$.
\end{exemplo}

\begin{regra}{Regra: mudança de base}\index{logaritmo!mudanca de base@mudança de base}
Para $a,b>0$ com $a,b\neq1$:
\[
a^x=b^{x\log_b a}\ \ (\text{em particular } a^x=e^{x\ln a}),\qquad
\log_a x=\frac{\log_b x}{\log_b a}\ \ (\text{em particular } \log_a x=\frac{\ln x}{\ln a}).
\]
\end{regra}

Graças a essas fórmulas, basta saber derivar e integrar $e^x$ e $\ln x$ para lidar com
qualquer base.

\begin{exemplo}{1.13}{A escala Richter}
Na escala Richter, de base 10, dois terremotos de magnitudes $R_1$ e $R_2$ e amplitudes
$A_1$ e $A_2$ satisfazem
\[
R_1-R_2=\log_{10}\Big(\frac{A_1}{A_2}\Big).
\]
Assim, um terremoto de magnitude 8 é 10 vezes mais intenso que um de magnitude 7, e
100 vezes mais intenso que um de magnitude 6. Comparemos o terremoto do Japão (2011,
magnitude 9) com o do Haiti (2010, magnitude 7,3).
\solucao
$9-7{,}3=\log_{10}(A_1/A_2)$, logo $A_1/A_2=10^{1{,}7}\approx50$. O terremoto do Japão
foi cerca de 50 vezes mais intenso.
\end{exemplo}

\subsection{Funções hiperbólicas}\index{funcoes hiperbolicas@funções hiperbólicas}

Certas combinações de $e^x$ e $e^{-x}$ aparecem com frequência. A forma de um cabo ou
corrente suspensa pelas extremidades, sob o próprio peso (a \emph{catenária}), é descrita
por uma delas.

\begin{definicao}
\[
\cosh x=\frac{e^x+e^{-x}}2,\qquad \senh x=\frac{e^x-e^{-x}}2,\qquad
\tgh x=\frac{\senh x}{\cosh x}=\frac{e^x-e^{-x}}{e^x+e^{-x}},
\]
\[
\cossech x=\frac1{\senh x},\qquad \sech x=\frac1{\cosh x},\qquad \cotgh x=\frac{\cosh x}{\senh x}.
\]
\end{definicao}

O nome vem da identidade $\cosh^2t-\senh^2t=1$: o ponto $(\cosh t,\senh t)$ percorre a
hipérbole $x^2-y^2=1$, assim como $(\cos t,\sen t)$ percorre o círculo $x^2+y^2=1$.

\figtikz{\begin{tikzpicture}
\begin{groupplot}[group style={group size=3 by 2,horizontal sep=1.1cm,vertical sep=1.3cm},painel,width=5.2cm,height=4.4cm,
  xmin=-3,xmax=3,xtick={-2,2}]
\nextgroupplot[title={$y=\cosh x$},ymin=-0.5,ymax=5,ytick={1,4}]
\addplot[domain=-3:3]{cosh(x)};
\addplot[cinza,dashed,thin,domain=-3:3]{exp(x)/2};
\addplot[cinza,dashed,thin,domain=-3:3]{exp(-x)/2};
\nextgroupplot[title={$y=\senh x$},ymin=-5,ymax=5,ytick={-4,4}]
\addplot[domain=-3:3]{sinh(x)};
\addplot[cinza,dashed,thin,domain=-3:3]{exp(x)/2};
\addplot[cinza,dashed,thin,domain=-3:3]{-exp(-x)/2};
\nextgroupplot[title={$y=\tgh x$},ymin=-2,ymax=2,ytick={-1,1}]
\addplot[domain=-3:3]{tanh(x)};
\draw[assint] (axis cs:-3,1) -- (axis cs:3,1);
\draw[assint] (axis cs:-3,-1) -- (axis cs:3,-1);
\nextgroupplot[title={$y=\cotgh x$},ymin=-4,ymax=4,ytick={-1,1}]
\addplot[domain=-3:-0.25]{1/tanh(x)};
\addplot[domain=0.25:3]{1/tanh(x)};
\draw[assint] (axis cs:-3,1) -- (axis cs:3,1);
\draw[assint] (axis cs:-3,-1) -- (axis cs:3,-1);
\nextgroupplot[title={$y=\operatorname{sech} x$},ymin=-0.5,ymax=1.5,ytick={1}]
\addplot[domain=-3:3]{1/cosh(x)};
\nextgroupplot[title={$y=\operatorname{csch} x$},ymin=-4,ymax=4,ytick={-2,2}]
\addplot[domain=-3:-0.25]{1/sinh(x)};
\addplot[domain=0.25:3]{1/sinh(x)};
\end{groupplot}\end{tikzpicture}}{As funções hiperbólicas. Em cinza tracejado,
$\pm\frac12e^{\pm x}$: para $|x|$ grande, $\cosh x$ e $\senh x$ se comportam como
$\frac12e^{|x|}$. Em verde, as assíntotas $y=\pm1$ de $\tgh$ e $\cotgh$.}

\begin{regra}{Regra: identidades hiperbólicas}\index{identidades hiperbolicas@identidades hiperbólicas}
\begin{tabular}{@{}ll@{}}
$\cosh(-x)=\cosh x$, \quad $\senh(-x)=-\senh x$ & $\cosh x+\senh x=e^x$, \quad $\cosh x-\senh x=e^{-x}$\\[2pt]
$\cosh^2x-\senh^2x=1$ & $1-\tgh^2x=\sech^2x$, \quad $\cotgh^2x-1=\cossech^2x$\\[2pt]
\multicolumn{2}{@{}l}{$\senh(x\pm y)=\senh x\cosh y\pm\cosh x\senh y$,\quad
$\cosh(x\pm y)=\cosh x\cosh y\pm\senh x\senh y$}
\end{tabular}
\end{regra}

\begin{exemplo}{1.14}{Avaliando funções hiperbólicas}
Se $\senh x=3/4$, então $\cosh^2x=1+(3/4)^2=25/16$; como $\cosh x\ge1$, temos
$\cosh x=5/4$. Daí $\tgh x=3/5$, $\cossech x=4/3$, $\sech x=4/5$ e $\cotgh x=5/3$.
\end{exemplo}

\textbf{Inversas hiperbólicas.} Como as hiperbólicas são combinações de exponenciais, suas
inversas se escrevem com logaritmos:
\[
\arcsenh x=\ln\!\big(x+\sqrt{x^2+1}\big),\qquad
\arccosh x=\ln\!\big(x+\sqrt{x^2-1}\big)\ (x\ge1),\]\[
\arctgh x=\frac12\ln\Big(\frac{1+x}{1-x}\Big)\ (|x|<1).
\]

% ---------------------------------------------------------------------
\section*{Resumo do capítulo}\markright{Resumo do capítulo}
\addcontentsline{toc}{section}{Resumo do capítulo}
\begin{itemize}
\item Uma função associa a cada entrada exatamente uma saída. Sem domínio explícito, o
  domínio é o conjunto dos $x$ em que a fórmula define um número real.
\item Teste da reta vertical: um gráfico de função é cortado no máximo uma vez por cada
  reta vertical. Uma função tem quantos zeros quiser, mas no máximo um intercepto $y$.
\item Para definir $g\circ f$, a imagem de $f$ deve estar contida no domínio de $g$; em
  geral $f\circ g\neq g\circ f$.
\item Funções pares são simétricas em relação ao eixo $y$; ímpares, em relação à origem.
\item A inclinação $m=\Delta y/\Delta x$ é a taxa de variação constante de uma função afim.
\item Translações, alongamentos, compressões e reflexões transformam gráficos conhecidos.
\item Ângulos em radianos: $\pi$ rad $=180^\circ$. Seno e cosseno têm período $2\pi$;
  tangente, período $\pi$.
\item Uma função tem inversa se, e somente se, é injetora; o gráfico de $f^{-1}$ é o reflexo
  do de $f$ em $y=x$. Restringindo o domínio, obtêm-se as inversas trigonométricas.
\item $b^x$ e $\log_b x$ são inversas; $\ln x=\log_e x$. Qualquer base se reduz à base $e$.
\item $\cosh x=(e^x+e^{-x})/2$, $\senh x=(e^x-e^{-x})/2$ e $\cosh^2x-\senh^2x=1$.
\end{itemize}

% Gerado por ferramentas/gerar.py; não edite à mão.
\section*{Exercícios selecionados}
\addcontentsline{toc}{section}{Exercícios selecionados}
Problemas escolhidos para desenvolver a resolução de problemas. As soluções completas,
com animações dos gráficos, estão no repositório do livro: \url{https://github.com/AleffAzevedo/calculo-1-pt/tree/main/exercicios/cap1}.

\begin{enumerate}[label=\textbf{\thechapter.\arabic*}, leftmargin=2.6em, itemsep=6pt]
\item \textbf{Composição e domínio.} Sejam $f(x)=\dfrac{3}{2x+1}$ e $g(x)=\dfrac{2}{x}$. Encontre $(f\circ g)(x)$ e
$(g\circ f)(x)$ e determine o domínio de cada composta.
\item \textbf{Custo e ponto de equilíbrio.} Uma confeitaria vende cupcakes em festivais. Montar a barraca custa R\$\,125,00 e cada
cupcake custa R\$\,0,75 para ser produzido.
(a) Escreva o custo total $C$ em função do número $x$ de cupcakes.
(b) Quanto custa produzir 160 cupcakes?
(c) Vendendo cada cupcake a R\$\,1,50, quantos ela precisa vender para começar a ter lucro?
\item \textbf{Identidade trigonométrica.} Verifique a identidade $\dfrac{1}{1-\sin\alpha}+\dfrac{1}{1+\sin\alpha}=2\sec^2\alpha$.
\item \textbf{Horas de luz do dia.} O número de horas de luz do dia em uma cidade é modelado por
$N(t)=12+3\sin\!\left[\frac{2\pi}{365}(t-79)\right]$, em que $t$ é o número de dias após
1º de janeiro.
(a) Encontre a amplitude e o período.
(b) Quantas horas de luz há no dia mais longo do ano? E no mais curto?
(c) Quantas horas de luz há 90 dias após 1º de janeiro?
\item \textbf{Custo de remoção de uma toxina.} O custo (em milhares de reais) para remover $p$ partes por bilhão (ppb) de uma toxina de um
lago é $C(p)=\dfrac{75p}{85-p}$, válido para $0\le p<85$.
(a) Calcule o custo para remover 25, 40 e 50 ppb.
(b) Encontre a função inversa $p=C^{-1}(C)$.
(c) Quanto da toxina se remove com R\$\,50 mil?
\item \textbf{Profundidade da maré.} A profundidade (em pés) da água em um cais é $D(t)=5\sin\!\left(\frac{\pi}{6}t-\frac{7\pi}{6}\right)+8$,
em que $t$ é o número de horas após a meia-noite. Determine o primeiro instante após a
meia-noite em que a profundidade é 11,75 pés.
\item \textbf{Expressões hiperbólicas.} Reescreva em termos de exponenciais e simplifique:
(a) $2\cosh(\ln x)$; (b) $\cosh 4x+\sinh 4x$; (c) $\cosh 2x-\sinh 2x$;
(d) $\ln(\cosh x+\sinh x)+\ln(\cosh x-\sinh x)$.
\item \textbf{Decaimento do iodo-131.} O iodo-131 decai segundo $Q(t)=Q_0e^{-0{,}08664t}$, com $t$ em dias. Quanto tempo leva
(arredondando para o dia mais próximo) para 95\% de uma amostra decair?
\item \textbf{Crescimento de bactérias.} Uma colônia de bactérias dobra a cada 12 horas e começa com 1000 bactérias.
(a) Usando $Q=Q_0e^{kt}$, determine $k$ (quatro casas decimais).
(b) Aproximadamente quanto tempo leva para haver 200\,000 bactérias?
\item \textbf{Escala Richter.} O terremoto de 1906 em São Francisco teve magnitude 8,3 na escala Richter. No mesmo
período, um terremoto de magnitude 4,9 no Japão causou poucos danos. Quantas vezes o de
São Francisco foi mais intenso?
\item \textbf{População sazonal.} A população de uma cidade litorânea, em milhares, é modelada por
$P(t)=82{,}5-67{,}5\cos\!\left(\frac\pi6t\right)$, com $t$ em meses ($t=0$ é 1º de janeiro).
Em que intervalos do ano a população é menor que 20 mil? E maior que 140 mil?
\end{enumerate}

\subsection*{Respostas}
{\small\begin{enumerate}[label=\textbf{\thechapter.\arabic*}, leftmargin=2.6em, itemsep=2pt]
\item $(f\circ g)(x)=\dfrac{3x}{x+4}$, com $x\neq0,-4$;\quad
$(g\circ f)(x)=\dfrac{2(2x+1)}{3}$, com $x\neq-\tfrac12$.
\item (a) $C(x)=0{,}75x+125$; (b) R\$\,245,00; (c) 167 cupcakes.
\item Some as frações: o denominador comum é $1-\sin^2\alpha=\cos^2\alpha$.
\item (a) Amplitude 3, período 365 dias; (b) 15 h e 9 h; (c) $\approx12{,}56$ h.
\item (a) $\approx$ R\$\,31,25 mil, R\$\,66,67 mil e R\$\,107,14 mil;
(b) $p=\dfrac{85C}{C+75}$; (c) 34 ppb.
\item $t=7+\frac6\pi\sin^{-1}(0{,}75)\approx8{,}62$ h, por volta de 8h37.
\item (a) $x+\frac1x$ ($x>0$); (b) $e^{4x}$; (c) $e^{-2x}$; (d) $0$.
\item $t=\dfrac{\ln 20}{0{,}08664}\approx35$ dias.
\item (a) $k=\frac{\ln2}{12}\approx0{,}0578$ por hora; (b) $\approx92$ horas.
\item $10^{3{,}4}\approx2512$ vezes.
\item $P<20$ para $t\in[0;\,0{,}74)\cup(11{,}26;\,12]$ (início de dez. a fim de jan.);
$P>140$ para $t\in(4{,}95;\,7{,}05)$ (fim de maio ao início de ago.).
\end{enumerate}}


% =====================================================================
\chapter{Limites}

O conceito de limite está na base de todo o Cálculo: a derivada e a integral são, ambas,
definidas por limites. Neste capítulo construímos a ideia de limite, primeiro de modo
intuitivo, com tabelas e gráficos, depois com regras de cálculo e, por fim, com uma
definição precisa. Também estudamos a continuidade, uma propriedade que depende
diretamente do limite.

% ---------------------------------------------------------------------
\section{Uma prévia do Cálculo}

O Cálculo nasceu de dois problemas geométricos: encontrar a reta tangente a uma curva
(o \emph{problema da tangente}, que leva ao Cálculo Diferencial) e calcular a área sob uma
curva (o \emph{problema da área}, que leva ao Cálculo Integral). Os dois envolvem uma
mesma ideia: aproximar uma quantidade cada vez melhor e observar para onde as aproximações
tendem.

\subsection{O problema da tangente}

Para uma reta, a taxa de variação é a inclinação. Para uma curva qualquer, aproximamos a
taxa de variação em um ponto $(a,f(a))$ por \emph{retas secantes}, que passam por
$(a,f(a))$ e por um ponto próximo $(x,f(x))$.

\begin{definicao}
A \textbf{reta secante}\index{reta!secante} ao gráfico de $f$ pelos pontos $(a,f(a))$ e $(x,f(x))$ é a reta
que passa por esses pontos. Sua inclinação é
\[
m_{\sec}=\frac{f(x)-f(a)}{x-a}.
\]
\end{definicao}

À medida que $x$ se aproxima de $a$, as secantes se aproximam de uma reta, a
\textbf{reta tangente}\index{reta!tangente} ao gráfico em $(a,f(a))$, e $m_{\sec}$ se aproxima da inclinação
da tangente, que é a \emph{taxa de variação instantânea} de $f$ em $a$.

\figtikz{\begin{tikzpicture}
\begin{axis}[painel,width=7cm,height=5.2cm,xmin=-0.3,xmax=3.6,ymin=0,ymax=3.6,xtick={1},xticklabels={$a$},ytick=\empty]
\addplot[domain=-0.3:3.4]{0.25*x^2+1} node[pos=0.98,left,font=\scriptsize]{$f$};
\foreach \b/\c in {3/laranja!90,2.2/laranja!60,1.5/laranja!35}{\edef\tmp{\noexpand\addplot[\c,thin,domain=-0.3:3.6]{1.25+0.25*(\b+1)*(x-1)};}\tmp}
\addplot[verde,thick,domain=-0.3:3.6]{1.25+0.5*(x-1)} node[pos=0.96,below,font=\scriptsize]{tangente};
\fill (axis cs:1,1.25) circle (1.6pt);
\end{axis}\end{tikzpicture}}{Quando $x$ se aproxima de $a$, as secantes (laranja) se aproximam da reta tangente (verde).}

\begin{exemplo}{2.1}{Inclinações de secantes}
Estime a inclinação da tangente a $f(x)=x^2$ em $x=1$ usando as secantes por $(1,1)$ e
(a) $(2,4)$; (b) $(3/2,\,9/4)$.
\solucao
(a) $m_{\sec}=\dfrac{4-1}{2-1}=3$. (b) $m_{\sec}=\dfrac{9/4-1}{3/2-1}=\dfrac52=2{,}5$. O
ponto de (b) está mais perto de $(1,1)$, então $2{,}5$ é uma estimativa melhor. Com pontos
cada vez mais próximos, veremos que a inclinação da tangente é exatamente 2.
\end{exemplo}

O mesmo raciocínio se aplica ao movimento. Se $s(t)$ é a posição de um objeto no instante
$t$, a \textbf{velocidade média}\index{velocidade!media@média} no intervalo $[a,t]$ é
\[
v_{\text{méd}}=\frac{s(t)-s(a)}{t-a},
\]
e a \textbf{velocidade instantânea}\index{velocidade!instantanea@instantânea} em $t=a$ é o valor do qual as velocidades médias se
aproximam quando $t$ se aproxima de $a$, se esse valor existir.

\subsection{O problema da área}

Para estimar a área sob o gráfico de uma função positiva entre $x=a$ e $x=b$, dividimos o
intervalo em partes e somamos as áreas de retângulos. Quanto mais retângulos, melhor a
aproximação; a área exata é o valor para o qual essas somas tendem. Esse é o ponto de
partida do Cálculo Integral (capítulo~5).

\figtikz{\begin{tikzpicture}
\begin{axis}[painel,width=7cm,height=5cm,xmin=0,xmax=5.6,ymin=0,ymax=2.6,xtick={1,5},xticklabels={$a$,$b$},ytick=\empty]
\foreach \a in {1,1.5,...,4.5}{\edef\tmp{\noexpand\draw[fill=azul!15,draw=azul] (axis cs:\a,0) rectangle (axis cs:{\a+0.5},{-0.15*(\a-3.2)^2+2.2});}\tmp}
\addplot[domain=0.2:5.5]{-0.15*(x-3.2)^2+2.2};
\end{axis}\end{tikzpicture}}{A área sob a curva aproximada por retângulos.}

% ---------------------------------------------------------------------
\section{O limite de uma função}

\subsection{Definição intuitiva}

\begin{definicao}
Seja $f$ definida em todos os pontos de um intervalo aberto contendo $a$, exceto
possivelmente em $a$, e seja $L$ um número real. Se os valores $f(x)$ se aproximam de $L$
quando $x$ ($\neq a$) se aproxima de $a$, dizemos que o \textbf{limite}\index{limite} de $f(x)$ quando
$x$ tende a $a$ é $L$, e escrevemos
\[
\lim_{x\to a}f(x)=L .
\]
Em outras palavras: à medida que $x$ fica perto de $a$, $f(x)$ fica perto de $L$ e
permanece perto.
\end{definicao}

\begin{nota}
O limite não depende do valor $f(a)$; $f$ nem precisa estar definida em $a$. Importa
apenas o comportamento de $f$ \emph{perto} de $a$.
\end{nota}

\begin{estrategia}{estimar um limite com uma tabela}\index{limite!estimado por tabela}
\begin{enumerate}[label=\arabic*.]
\item Monte duas colunas de valores de $x$ que se aproximam de $a$: uma pela esquerda
  ($a-0{,}1$, $a-0{,}01$, $a-0{,}001,\dots$) e outra pela direita ($a+0{,}1$,
  $a+0{,}01,\dots$), com os respectivos $f(x)$.
\item Se as duas colunas de $f(x)$ se aproximam de um mesmo valor $L$, a estimativa é
  $\lim_{x\to a}f(x)=L$.
\item Confirme com o gráfico de $f$ perto de $x=a$.
\end{enumerate}
\end{estrategia}

\begin{exemplo}{2.2}{Estimando $\lim_{x\to0}\frac{\sin x}{x}$}
Calculando $(\sin x)/x$ para $x=\pm0{,}1;\ \pm0{,}01;\ \pm0{,}001$, obtemos
$0{,}99833\ldots$, $0{,}99998\ldots$ e $0{,}9999998\ldots$ (a função é par, então os valores
pela esquerda e pela direita coincidem). Os valores se aproximam de 1, e o gráfico confirma:
\[
\lim_{x\to0}\frac{\sin x}{x}=1 .
\]
Provaremos esse resultado na seção~2.3.
\figtikz{\begin{tikzpicture}
\begin{axis}[painel,width=8cm,height=4.4cm,xmin=-8,xmax=8,ymin=-0.45,ymax=1.2,xtick={-6,-3,3,6},ytick={0.5,1}]
\addplot[domain=-8:-0.01,samples=150]{sin(deg(x))/x}; \addplot[domain=0.01:8,samples=150]{sin(deg(x))/x};
\draw[azul,fill=white] (axis cs:0,1) circle (2pt);
\end{axis}\end{tikzpicture}}{Gráfico de $f(x)=(\sin x)/x$ perto de $x=0$.}
\end{exemplo}

Dois limites básicos, intuitivamente evidentes, servem de ponto de partida:

\begin{teorema}[Dois limites importantes]{2.1}
Para todo número real $a$ e toda constante $c$:
\qquad $\displaystyle\lim_{x\to a}x=a$ \qquad e \qquad $\displaystyle\lim_{x\to a}c=c$.
\end{teorema}

\subsection{Quando o limite não existe}

Um limite pode não existir por diferentes motivos: a função pode se aproximar de valores
diferentes pela esquerda e pela direita, crescer sem limitação, ou oscilar.

\begin{exemplo}{2.3}{Um limite que não existe}
Considere $\lim_{x\to0}\sin(1/x)$. Para a sequência $x=\dfrac2\pi,\dfrac2{3\pi},
\dfrac2{5\pi},\dots$, que tende a 0, os valores de $\sin(1/x)$ são $1,-1,1,-1,\dots$. Eles não
se aproximam de nenhum número, logo o limite não existe. O gráfico oscila cada vez mais
rapidamente entre $-1$ e $1$ perto de $x=0$.
\figtikz{\begin{tikzpicture}
\begin{axis}[painel,width=8cm,height=4.4cm,xmin=-0.5,xmax=0.5,ymin=-1.3,ymax=1.3,xtick={-0.4,-0.2,0.2,0.4},ytick={-1,1},
  xticklabel style={/pgf/number format/use comma}]
\addplot[thin,domain=-0.5:-0.012,samples=1200]{sin(deg(1/x))}; \addplot[thin,domain=0.012:0.5,samples=1200]{sin(deg(1/x))};
\end{axis}\end{tikzpicture}}{$\sin(1/x)$ oscila entre $-1$ e $1$ quando
$x\to0$.}
\end{exemplo}

\subsection{Limites laterais}

\begin{definicao}
\textbf{Limite à esquerda:}\index{limite!lateral} se $f$ está definida em um intervalo $(c,a)$ e $f(x)$ se
aproxima de $L$ quando $x$ se aproxima de $a$ com $x<a$, escrevemos
$\displaystyle\lim_{x\to a^-}f(x)=L$.

\textbf{Limite à direita:} se $f$ está definida em um intervalo $(a,c)$ e $f(x)$ se
aproxima de $L$ quando $x$ se aproxima de $a$ com $x>a$, escrevemos
$\displaystyle\lim_{x\to a^+}f(x)=L$.
\end{definicao}

\begin{exemplo}{2.4}{Limites laterais}
Para $f(x)=\begin{cases}x+1, & x<2\\ x^2-4, & x\ge2,\end{cases}$ temos
$\lim_{x\to2^-}f(x)=2+1=3$ e $\lim_{x\to2^+}f(x)=2^2-4=0$. Como os limites laterais são
diferentes, $\lim_{x\to2}f(x)$ não existe.
\figtikz{\begin{tikzpicture}
\begin{axis}[painel,width=6.2cm,height=5cm,xmin=-3,xmax=4,ymin=-2.2,ymax=5.5,xtick={-1,2},ytick={3}]
\addplot[domain=-3:2]{x+1}; \addplot[laranja,domain=2:3.05]{x^2-4};
\draw[azul,fill=white] (axis cs:2,3) circle (2pt); \fill[laranja] (axis cs:2,0) circle (2pt);
\end{axis}\end{tikzpicture}}{Limites laterais diferentes em $x=2$.}
\end{exemplo}

\begin{teorema}[Limites laterais e bilaterais]{2.2}
Seja $f$ definida em um intervalo aberto contendo $a$, exceto possivelmente em $a$. Então
\[
\lim_{x\to a}f(x)=L \iff \lim_{x\to a^-}f(x)=L \ \text{ e } \ \lim_{x\to a^+}f(x)=L .
\]
\end{teorema}

\subsection{Limites infinitos}

Quando $f(x)$ cresce sem limitação à medida que $x$ se aproxima de $a$ (pela esquerda,
pela direita ou pelos dois lados), escrevemos, respectivamente,
\[
\lim_{x\to a^-}f(x)=+\infty,\qquad \lim_{x\to a^+}f(x)=+\infty,\qquad
\lim_{x\to a}f(x)=+\infty ;
\]
analogamente para $-\infty$, quando $f(x)$ decresce sem limitação. Atenção: dizer que o
limite é $\infty$ é uma forma de descrever \emph{como} o limite deixa de existir, pois
$\infty$ não é um número real. Por exemplo, para $f(x)=1/x$:
$\lim_{x\to0^-}\frac1x=-\infty$ e $\lim_{x\to0^+}\frac1x=+\infty$.

\figtikz{\begin{tikzpicture}
\begin{axis}[painel,width=6cm,height=5cm,xmin=-3,xmax=3,ymin=-5,ymax=5,xtick={-2,-1,1,2},ytick={-4,-2,2,4}]
\addplot[domain=-3:-0.2]{1/x}; \addplot[domain=0.2:3]{1/x};
\end{axis}\end{tikzpicture}}{O gráfico de $1/x$ perto de $x=0$.}

\begin{teorema}[Limites infinitos de potências inteiras]{2.3}\index{limite!infinito}
Se $n$ é um inteiro positivo par, $\displaystyle\lim_{x\to a}\frac{1}{(x-a)^n}=+\infty$.
Se $n$ é um inteiro positivo ímpar,
$\displaystyle\lim_{x\to a^+}\frac{1}{(x-a)^n}=+\infty$ e
$\displaystyle\lim_{x\to a^-}\frac{1}{(x-a)^n}=-\infty$.
\end{teorema}

\begin{definicao}
A reta $x=a$ é uma \textbf{assíntota vertical}\index{assintota@assíntota!vertical} de $f$ se pelo menos um dos limites
$\lim_{x\to a^-}f(x)$, $\lim_{x\to a^+}f(x)$ ou $\lim_{x\to a}f(x)$ for $+\infty$ ou
$-\infty$.
\end{definicao}

Por exemplo, $f(x)=\dfrac{1}{(x-3)^4}$ tem assíntota vertical $x=3$, pois
$\lim_{x\to3}f(x)=+\infty$ pelo Teorema~2.3.

% ---------------------------------------------------------------------
\section{As propriedades dos limites}

\subsection{Calculando limites com as propriedades}

Estimar limites com tabelas e gráficos é útil, mas trabalhoso e às vezes enganoso. As
propriedades a seguir permitem calcular limites diretamente.

\begin{teorema}[Propriedades dos limites]{2.4}
Sejam $f$ e $g$ definidas para $x\neq a$ em um intervalo aberto contendo $a$, com
$\lim_{x\to a}f(x)=L$ e $\lim_{x\to a}g(x)=M$ reais, e seja $c$ uma constante. Então:
\begin{itemize}
\item \textbf{Soma e diferença:} $\lim_{x\to a}\big(f(x)\pm g(x)\big)=L\pm M$.
\item \textbf{Múltiplo constante:} $\lim_{x\to a}c\,f(x)=cL$.
\item \textbf{Produto:} $\lim_{x\to a}f(x)\,g(x)=L\cdot M$.
\item \textbf{Quociente:} $\lim_{x\to a}\dfrac{f(x)}{g(x)}=\dfrac LM$, se $M\neq0$.
\item \textbf{Potência:} $\lim_{x\to a}\big(f(x)\big)^n=L^n$, para todo inteiro positivo $n$.
\item \textbf{Raiz:} $\lim_{x\to a}\sqrt[n]{f(x)}=\sqrt[n]{L}$, para todo $L$ se $n$ é ímpar,
  e para $L\ge0$ se $n$ é par e $f(x)\ge0$.
\end{itemize}
\end{teorema}

\begin{exemplo}{2.5}{Usando as propriedades}
\[
\lim_{x\to2}\frac{2x^2-3x+1}{x^3+4}
=\frac{2\lim x^2-3\lim x+\lim 1}{\lim x^3+\lim 4}
=\frac{2(4)-3(2)+1}{8+4}=\frac{3}{12}=\frac14 .
\]
\end{exemplo}

Aplicando as propriedades repetidamente, obtemos o resultado geral para polinômios e
funções racionais.

\begin{teorema}[Limites de polinômios e funções racionais]{2.5}
Se $p$ e $q$ são polinômios e $a$ é real, então
$\displaystyle\lim_{x\to a}p(x)=p(a)$ e
$\displaystyle\lim_{x\to a}\frac{p(x)}{q(x)}=\frac{p(a)}{q(a)}$ quando $q(a)\neq0$.
\end{teorema}

\subsection{A forma indeterminada $0/0$}

Quando $q(a)=0$, a propriedade do quociente não se aplica diretamente. Se também $p(a)=0$,
temos a \textbf{forma indeterminada $0/0$}: o limite pode existir ou não, e é preciso
manipular a expressão. A ideia central é que \emph{se $f(x)=g(x)$ para todo $x\neq a$
perto de $a$, então as duas funções têm o mesmo limite em $a$}. Por exemplo,
$\dfrac{x^2-1}{x-1}=x+1$ para $x\neq1$, logo
$\lim_{x\to1}\dfrac{x^2-1}{x-1}=\lim_{x\to1}(x+1)=2$.

\begin{estrategia}{limites da forma $0/0$}\index{forma indeterminada}
\begin{enumerate}[label=\arabic*.]
\item Verifique que as propriedades não se aplicam diretamente (forma $0/0$).
\item Encontre uma função igual a $f(x)/g(x)$ para $x\neq a$:
  (a) se forem polinômios, fatore e cancele os fatores comuns;
  (b) se houver diferença com raiz quadrada, multiplique pelo conjugado;
  (c) se houver fração composta, simplifique-a.
\item Aplique as propriedades dos limites.
\end{enumerate}
\end{estrategia}

\begin{exemplo}{2.6}{Fatorando e cancelando}
\[
\lim_{x\to3}\frac{x^2-3x}{2x^2-5x-3}=\lim_{x\to3}\frac{x(x-3)}{(x-3)(2x+1)}
=\lim_{x\to3}\frac{x}{2x+1}=\frac37 .
\]
\end{exemplo}

\begin{exemplo}{2.7}{Multiplicando pelo conjugado}
Calcule $\displaystyle\lim_{x\to-1}\frac{\sqrt{x+2}-1}{x+1}$.
\solucao
A expressão tem a forma $0/0$ em $x=-1$. Multiplicamos numerador e denominador pelo
conjugado $\sqrt{x+2}+1$:
\begin{align*}
\lim_{x\to-1}\frac{\sqrt{x+2}-1}{x+1}
&=\lim_{x\to-1}\frac{(\sqrt{x+2}-1)(\sqrt{x+2}+1)}{(x+1)(\sqrt{x+2}+1)}
 =\lim_{x\to-1}\frac{x+1}{(x+1)(\sqrt{x+2}+1)}\\
&=\lim_{x\to-1}\frac{1}{\sqrt{x+2}+1}=\frac12 .
\end{align*}
\end{exemplo}

\begin{exemplo}{2.8}{Simplificando uma fração composta}
\[
\lim_{x\to1}\frac{\frac{1}{x+1}-\frac12}{x-1}
=\lim_{x\to1}\frac{2-(x+1)}{2(x+1)(x-1)}
=\lim_{x\to1}\frac{-(x-1)}{2(x+1)(x-1)}
=\lim_{x\to1}\frac{-1}{2(x+1)}=-\frac14 .
\]
\end{exemplo}

\textbf{Forma $K/0$ com $K\neq0$.} Se o numerador tende a $K\neq0$ e o denominador tende a
0, o limite não é finito. Estudamos os limites laterais pelo sinal: por exemplo, em
$\lim_{x\to2}\frac{x-3}{x^2-2x}$, o numerador tende a $-1$ e o denominador $x(x-2)$ tende
a $0^-$ pela esquerda e a $0^+$ pela direita. Logo o limite é $+\infty$ pela esquerda e
$-\infty$ pela direita.

\subsection{O Teorema do Confronto}

\begin{teorema}[Teorema do Confronto]{2.6}
Sejam $f$, $g$ e $h$ definidas para $x\neq a$ em um intervalo aberto contendo $a$. Se
$f(x)\le g(x)\le h(x)$ nesse intervalo, para $x\neq a$, e
$\lim_{x\to a}f(x)=L=\lim_{x\to a}h(x)$, então $\lim_{x\to a}g(x)=L$.
\end{teorema}

\figtikz{\begin{tikzpicture}
\begin{axis}[painel,width=7cm,height=5cm,xmin=-1,xmax=3.2,ymin=-0.5,ymax=3.2,xtick={1},xticklabels={$a$},ytick=\empty]
\addplot[domain=-1:3]{2+0.5*(x-1)^2} node[pos=0.97,above left,font=\scriptsize]{$h$};
\addplot[verde,domain=-1:3]{2-0.3*(x-1)^2+0.15*(x-1)^2*sin(deg(4*x))} node[pos=0.97,left,font=\scriptsize]{$g$};
\addplot[laranja,domain=-0.6:2.6]{2-(x-1)^2} node[pos=0.97,left,font=\scriptsize]{$f$};
\fill (axis cs:1,2) circle (1.6pt);
\end{axis}\end{tikzpicture}}{O Teorema do Confronto: $g$ fica ``espremida''
entre $f$ e $h$.}

\begin{exemplo}{2.9}{Aplicando o Teorema do Confronto}
Como $-1\le\cos x\le1$, temos $-|x|\le x\cos x\le|x|$. Como $\lim_{x\to0}(-|x|)=0=
\lim_{x\to0}|x|$, concluímos que $\lim_{x\to0}x\cos x=0$.
\end{exemplo}

\textbf{Dois limites trigonométricos fundamentais.} Pelo círculo unitário, para
$0<\theta<\pi/2$, valem $0<\sin\theta<\theta<\tan\theta$ (o arco de comprimento $\theta$
fica entre o segmento vertical de comprimento $\sin\theta$ e o segmento tangente de
comprimento $\tan\theta$). Pelo confronto, isso dá $\lim_{\theta\to0}\sin\theta=0$ e,
daí, $\lim_{\theta\to0}\cos\theta=1$.

\figtikz{\begin{tikzpicture}[scale=2.3,font=\small]
\draw[->] (-0.2,0) -- (1.35,0) node[right]{$x$}; \draw[->] (0,-0.2) -- (0,1.15) node[above]{$y$};
\draw[thick] (0,0) ++(0:1) arc (0:95:1);
\fill[azul!15] (0,0) -- (1,0) arc (0:40:1) -- cycle;
\draw[thick,azul] (0,0) -- (40:1); \draw[laranja,thick] (40:1) -- ({cos(40)},0); \draw[laranja,thin] ({cos(40)},0.32) -- (0.45,0.62) node[above,font=\footnotesize]{$\sin\theta$};
\draw[thick] (0,0) -- (1,{tan(40)}); \draw[verde,thick] (1,0) -- (1,{tan(40)}) node[midway,right]{$\tan\theta$};
\draw[azul,very thick] (1,0) arc (0:40:1); \draw (0.22,0) arc (0:40:0.22); \node at (20:0.33) {$\theta$};
\node[below] at (1,0) {$1$};
\end{tikzpicture}}{Para $0<\theta<\pi/2$: $\sin\theta<\theta<\tan\theta$ (segmento laranja, arco azul e segmento verde).}

Dividindo $\sin\theta<\theta<\tan\theta$ por $\sin\theta>0$ e invertendo, obtemos
$\cos\theta<\dfrac{\sin\theta}{\theta}<1$. Essa desigualdade vale também para
$-\pi/2<\theta<0$, porque as três funções são pares. Pelo Teorema do Confronto:

\begin{teorema}[Limites trigonométricos especiais]{2.7}
\[
\lim_{\theta\to0}\frac{\sin\theta}{\theta}=1
\qquad\text{e}\qquad
\lim_{\theta\to0}\frac{1-\cos\theta}{\theta}=0 .
\]
\end{teorema}
O segundo limite sai do primeiro multiplicando pelo conjugado:
\[
\frac{1-\cos\theta}{\theta}=\frac{1-\cos^2\theta}{\theta(1+\cos\theta)}
=\frac{\sin\theta}{\theta}\cdot\frac{\sin\theta}{1+\cos\theta}\ \longrightarrow\ 1\cdot\frac02=0 .
\]
Esses dois limites serão essenciais para derivar seno e cosseno no capítulo~3.

% ---------------------------------------------------------------------
\section{Continuidade}\index{continuidade}

\subsection{Continuidade em um ponto}

Intuitivamente, uma função é contínua quando seu gráfico pode ser traçado sem tirar o lápis
do papel. A definição precisa exige três condições.

\begin{definicao}
Uma função $f$ é \textbf{contínua no ponto $a$} se, e somente se:
\begin{enumerate}[label=(\roman*)]
\item $f(a)$ está definida;
\item $\lim_{x\to a}f(x)$ existe;
\item $\lim_{x\to a}f(x)=f(a)$.
\end{enumerate}
Se alguma condição falha, $f$ é \textbf{descontínua} em $a$.
\end{definicao}

Para verificar a continuidade em $a$, testamos as condições na ordem: se uma falha,
paramos. Por exemplo, $f(x)=(x^2-4)/(x-2)$ não é contínua em $x=2$, pois $f(2)$ não está
definida. Já
\[
f(x)=\begin{cases}x^3, & x\neq1\\ 0, & x=1\end{cases}
\]
está definida em 1 e tem $\lim_{x\to1}f(x)=1$, mas $1\neq f(1)=0$: falha a condição (iii).

\begin{teorema}[Continuidade de polinômios e funções racionais]{2.8}
Polinômios e funções racionais são contínuos em todos os pontos de seus domínios.
\end{teorema}
Isso decorre do Teorema~2.5: para $q(a)\neq0$, $\lim_{x\to a}p(x)/q(x)=p(a)/q(a)$.

\subsection{Tipos de descontinuidade}

\begin{definicao}
Se $f$ é descontínua em $a$, então $f$ tem:
\begin{enumerate}[label=\arabic*.]
\item uma \textbf{descontinuidade removível}\index{descontinuidade} se $\lim_{x\to a}f(x)$ existe (é um número
  real);
\item uma \textbf{descontinuidade de salto} se os limites laterais existem, mas são
  diferentes;
\item uma \textbf{descontinuidade infinita} se pelo menos um limite lateral é $\pm\infty$.
\end{enumerate}
\end{definicao}

\figtikz{\begin{tikzpicture}
\begin{groupplot}[group style={group size=3 by 1,horizontal sep=1.1cm},painel,width=5.2cm,height=4.4cm,
  xmin=-0.5,xmax=4,ymin=-0.5,ymax=4,xtick={2},xticklabels={$a$},ytick=\empty]
\nextgroupplot[title={(a) removível}]
\addplot[domain=-0.5:4]{0.5*x+1};
\draw[azul,thick,fill=white] (axis cs:2,2) circle (2pt);
\fill[laranja] (axis cs:2,3.2) circle (2pt);
\nextgroupplot[title={(b) de salto}]
\addplot[domain=-0.5:2]{0.5*x+0.6};
\addplot[domain=2:4]{0.5*x+1.6};
\fill[azul] (axis cs:2,1.6) circle (2pt);
\draw[azul,thick,fill=white] (axis cs:2,2.6) circle (2pt);
\nextgroupplot[title={(c) infinita},ymin=-4,ymax=4]
\addplot[domain=-0.5:1.75]{1/(x-2)+0.5};
\addplot[domain=2.25:4]{1/(x-2)+0.5};
\draw[assint] (axis cs:2,-4) -- (axis cs:2,4);
\end{groupplot}\end{tikzpicture}}{Tipos de descontinuidade em $x=a$: (a) removível, com
o limite existindo, mas diferente de $f(a)$ ou com $f(a)$ indefinido; (b) de salto, com
limites laterais finitos e diferentes; (c) infinita, com limite lateral $\pm\infty$.}

\begin{exemplo}{2.10}{Classificando descontinuidades}
(a) $f(x)=\dfrac{x^2-4}{x-2}$ em $x=2$: $\lim_{x\to2}f(x)=\lim(x+2)=4$ existe;
descontinuidade \emph{removível}.

(b) $f(x)=\begin{cases}-x^2+4, & x\le3\\ 4x-8, & x>3\end{cases}$ em $x=3$: os limites
laterais são $-5$ e $4$; descontinuidade \emph{de salto}.

(c) $f(x)=\dfrac{x+2}{x+1}$ em $x=-1$: pela esquerda o limite é $-\infty$ e pela direita
$+\infty$; descontinuidade \emph{infinita}.
\end{exemplo}

\subsection{Continuidade em um intervalo}

\begin{definicao}[Continuidade lateral]
$f$ é \textbf{contínua à direita}\index{continuidade!lateral} em $a$ se $\lim_{x\to a^+}f(x)=f(a)$, e
\textbf{contínua à esquerda} em $a$ se $\lim_{x\to a^-}f(x)=f(a)$.
\end{definicao}

Uma função é contínua em um intervalo aberto se é contínua em cada ponto dele. É contínua
em $[a,b]$ se é contínua em $(a,b)$, contínua à direita em $a$ e contínua à esquerda em
$b$. Por exemplo, $f(x)=\sqrt{4-x^2}$ é contínua em $[-2,2]$.

\begin{teorema}[Função composta]{2.9}
Se $f$ é contínua em $L$ e $\lim_{x\to a}g(x)=L$, então
\[
\lim_{x\to a}f\big(g(x)\big)=f\Big(\lim_{x\to a}g(x)\Big)=f(L).
\]
\end{teorema}
Por exemplo, $\lim_{x\to\pi/2}\cos\!\big(x-\tfrac\pi2\big)=\cos0=1$. Em particular, a
composta de funções contínuas é contínua.

\begin{teorema}[Continuidade das funções trigonométricas]{2.10}
As funções trigonométricas são contínuas em todo o seu domínio.
\end{teorema}

Mais geralmente, exponenciais, logaritmos, raízes e todas as funções construídas a partir
delas e de polinômios por operações algébricas e composição são contínuas onde estão
definidas.

\subsection{O Teorema do Valor Intermediário}

\begin{teorema}[Teorema do Valor Intermediário]{2.11}
Seja $f$ contínua em um intervalo fechado e limitado $[a,b]$. Se $z$ é qualquer número
real entre $f(a)$ e $f(b)$, então existe $c\in[a,b]$ tal que $f(c)=z$.
\end{teorema}

\figtikz{\begin{tikzpicture}
\begin{axis}[painel,width=7cm,height=5.4cm,xmin=0,xmax=5.6,ymin=0,ymax=5,xtick={0.3,4.61,5},xticklabels={$a$,$c$,$b$},ytick={1.3645,3,4.333},yticklabels={$f(a)$,$z$,$f(b)$}]
\draw[dashed] (axis cs:0,1.3645) -| (axis cs:0.3,0); \draw[dashed] (axis cs:0,4.333) -| (axis cs:5,0);
\draw[laranja,dashed] (axis cs:0,3) -| (axis cs:4.61,0);
\addplot[domain=0.3:5]{x^3/6-x^2+1.5*x+1};
\fill[laranja] (axis cs:4.61,3) circle (1.6pt);
\end{axis}\end{tikzpicture}}{Existe $c\in[a,b]$ com $f(c)=z$.}

\begin{exemplo}{2.11}{Garantindo um zero}
Mostre que $f(x)=x-\cos x$ tem pelo menos um zero.
\solucao
$f$ é contínua em $\R$. Como $f(0)=-1<0$ e $f(\pi/2)=\pi/2>0$, o Teorema do Valor
Intermediário garante que existe $c\in[0,\pi/2]$ com $f(c)=0$.
\end{exemplo}

\begin{nota}
A continuidade é essencial. Para $f(x)=1/x$, temos $f(-1)=-1<0$ e $f(1)=1>0$, mas $f$ não
tem zero em $[-1,1]$, pois não é contínua em $x=0$.
\end{nota}

% ---------------------------------------------------------------------
\section{A definição precisa de limite}\index{limite!definicao precisa (-)@definição precisa ($\varepsilon$-$\delta$)}

\subsection{Quantificando a proximidade}

A definição intuitiva fala em $f(x)$ ``se aproximar'' de $L$. Para torná-la precisa,
medimos distâncias: $|f(x)-L|$ é a distância entre $f(x)$ e $L$, e $|x-a|$ é a distância
entre $x$ e $a$. Dizer que $f(x)$ fica tão perto de $L$ quanto quisermos significa: para
qualquer tolerância $\varepsilon>0$, conseguimos garantir $|f(x)-L|<\varepsilon$ tomando
$x$ suficientemente perto de $a$, a uma distância menor que algum $\delta>0$.

\begin{definicao}[Definição $\varepsilon$-$\delta$ de limite]
Seja $f$ definida para $x\neq a$ em um intervalo aberto contendo $a$, e seja $L$ real.
Então $\lim_{x\to a}f(x)=L$ se, para todo $\varepsilon>0$, existe $\delta>0$ tal que
\[
0<|x-a|<\delta\ \Longrightarrow\ |f(x)-L|<\varepsilon .
\]
\end{definicao}

\figtikz{\begin{tikzpicture}
\begin{axis}[painel,width=8cm,height=6.2cm,xmin=0,xmax=4.2,ymin=0,ymax=4.6,
  xtick={1.6,2,2.4},xticklabels={$a-\delta$,$a$,$a+\delta$},
  ytick={1.95,2.6,3.25},yticklabels={$L-\varepsilon$,$L$,$L+\varepsilon$},
  x tick label style={font=\scriptsize,rotate=0},axis on top]
\fill[laranja!12] (axis cs:0,1.95) rectangle (axis cs:4.2,3.25);
\fill[azul!12] (axis cs:1.6,0) rectangle (axis cs:2.4,4.6);
\draw[laranja,dashed,thin] (axis cs:0,1.95) -- (axis cs:4.2,1.95);
\draw[laranja,dashed,thin] (axis cs:0,3.25) -- (axis cs:4.2,3.25);
\draw[azul!60,dashed,thin] (axis cs:1.6,0) -- (axis cs:1.6,4.6);
\draw[azul!60,dashed,thin] (axis cs:2.4,0) -- (axis cs:2.4,4.6);
\addplot[domain=0.2:4]{0.15*x^3-0.9*x^2+2.1*x+0.6};
\draw[azul,thick,fill=white] (axis cs:2,2.6) circle (2pt);
\end{axis}\end{tikzpicture}}{A definição $\varepsilon$-$\delta$: para cada faixa
horizontal $(L-\varepsilon,L+\varepsilon)$ existe uma faixa vertical $(a-\delta,a+\delta)$
na qual o gráfico, exceto possivelmente em $x=a$, fica dentro da faixa horizontal.}

\begin{estrategia}{provar que $\lim_{x\to a}f(x)=L$}
\begin{enumerate}[label=\arabic*.]
\item Comece com: ``Seja $\varepsilon>0$.''
\item Encontre $\delta$ (geralmente trabalhando ``de trás para a frente'' a partir de
  $|f(x)-L|<\varepsilon$) e escreva: ``Escolha $\delta=\ldots$''.
\item Escreva: ``Suponha $0<|x-a|<\delta$.''
\item Usando essa hipótese, mostre que $|f(x)-L|<\varepsilon$.
\item Conclua: ``Portanto, $\lim_{x\to a}f(x)=L$.''
\end{enumerate}
\end{estrategia}

\begin{exemplo}{2.12}{Uma prova $\varepsilon$-$\delta$ com função afim}
Prove que $\lim_{x\to-1}(4x+1)=-3$.
\solucao
Rascunho: queremos $|(4x+1)-(-3)|=|4x+4|=4|x+1|<\varepsilon$, ou seja, $|x+1|<\varepsilon/4$.
Isso sugere $\delta=\varepsilon/4$.

Prova: seja $\varepsilon>0$. Escolha $\delta=\varepsilon/4$. Suponha $0<|x+1|<\delta$.
Então $|(4x+1)-(-3)|=4|x+1|<4\delta=\varepsilon$. Portanto,
$\lim_{x\to-1}(4x+1)=-3$.
\end{exemplo}

\begin{exemplo}{2.13}{Uma prova $\varepsilon$-$\delta$ com função quadrática}
Prove que $\lim_{x\to-1}(x^2-2x+3)=6$.
\solucao
Rascunho: $|(x^2-2x+3)-6|=|x+1|\cdot|x-3|$. O fator $|x-3|$ depende de $x$, então
primeiro o limitamos: se $|x+1|<1$, então $-1<x+1<1$, isto é, $-5<x-3<-3$, e logo
$|x-3|<5$. Basta então ter também $|x+1|<\varepsilon/5$.

Prova: seja $\varepsilon>0$. Escolha $\delta=\min\{1,\varepsilon/5\}$. Suponha
$0<|x+1|<\delta$. Então $|x+1|<1$, donde $|x-3|<5$, e $|x+1|<\varepsilon/5$. Logo
$|(x^2-2x+3)-6|=|x+1|\,|x-3|<\tfrac{\varepsilon}{5}\cdot5=\varepsilon$.
\end{exemplo}

\subsection{Provando as propriedades dos limites}

A definição precisa permite demonstrar as propriedades da seção~2.3. Como exemplo,
provemos a da soma. Uma ferramenta recorrente é a \textbf{desigualdade triangular}\index{desigualdade triangular}:
$|a+b|\le|a|+|b|$ para quaisquer reais $a$ e $b$.

\emph{Prova da propriedade da soma.} Suponha $\lim_{x\to a}f(x)=L$ e
$\lim_{x\to a}g(x)=M$. Seja $\varepsilon>0$. Existe $\delta_1>0$ tal que
$0<|x-a|<\delta_1$ implica $|f(x)-L|<\varepsilon/2$, e existe $\delta_2>0$ tal que
$0<|x-a|<\delta_2$ implica $|g(x)-M|<\varepsilon/2$. Escolha $\delta=\min\{\delta_1,\delta_2\}$.
Se $0<|x-a|<\delta$, então
\[
|(f(x)+g(x))-(L+M)|\le|f(x)-L|+|g(x)-M|<\frac\varepsilon2+\frac\varepsilon2=\varepsilon. \qquad\square
\]

Para mostrar que um limite \emph{não} existe, basta exibir um $\varepsilon>0$ para o qual
nenhum $\delta$ funciona. Por exemplo, se $f(x)=-1$ para $x<0$ e $f(x)=1$ para $x\ge0$, com
$\varepsilon=1/2$ não há $\delta$ que sirva para nenhum candidato $L$: perto de 0 há valores
de $f$ iguais a $-1$ e a $1$, que distam 2 entre si.

\subsection{Limites laterais e infinitos}

As outras formas de limite têm definições análogas:
\begin{itemize}
\item $\lim_{x\to a^+}f(x)=L$ se, para todo $\varepsilon>0$, existe $\delta>0$ tal que
  $0<x-a<\delta\Rightarrow|f(x)-L|<\varepsilon$;
\item $\lim_{x\to a^-}f(x)=L$ se, para todo $\varepsilon>0$, existe $\delta>0$ tal que
  $-\delta<x-a<0\Rightarrow|f(x)-L|<\varepsilon$;
\item $\lim_{x\to a}f(x)=+\infty$ se, para todo $M>0$, existe $\delta>0$ tal que
  $0<|x-a|<\delta\Rightarrow f(x)>M$ (e $-\infty$ com $f(x)<-M$).
\end{itemize}

% ---------------------------------------------------------------------
\section*{Resumo do capítulo}\markright{Resumo do capítulo}
\addcontentsline{toc}{section}{Resumo do capítulo}
\begin{itemize}
\item A inclinação da tangente e a velocidade instantânea são limites de inclinações de
  secantes e de velocidades médias; a área sob uma curva é o limite de somas de áreas de
  retângulos.
\item $\lim_{x\to a}f(x)=L$ quando $f(x)$ fica arbitrariamente perto de $L$ para $x$
  suficientemente perto de $a$ ($x\neq a$); o valor $f(a)$ é irrelevante.
\item O limite bilateral existe se, e somente se, os laterais existem e são iguais.
\item Um limite infinito descreve uma assíntota vertical; $\infty$ não é um número.
\item Para polinômios e racionais (com denominador não nulo), basta substituir.
  Na forma $0/0$, fatore, use o conjugado ou simplifique antes.
\item Teorema do Confronto; $\lim_{\theta\to0}\frac{\sin\theta}{\theta}=1$ e
  $\lim_{\theta\to0}\frac{1-\cos\theta}{\theta}=0$.
\item Contínua em $a$: $f(a)$ definida, o limite existe e é igual a $f(a)$.
  Descontinuidades: removível, de salto, infinita.
\item Teorema do Valor Intermediário: uma função contínua em $[a,b]$ assume todos os
  valores entre $f(a)$ e $f(b)$.
\item Definição precisa: para todo $\varepsilon>0$ existe $\delta>0$ tal que
  $0<|x-a|<\delta\Rightarrow|f(x)-L|<\varepsilon$.
\end{itemize}

% Gerado por ferramentas/gerar.py; não edite à mão.
\section*{Exercícios selecionados}
\addcontentsline{toc}{section}{Exercícios selecionados}
Problemas escolhidos para desenvolver a resolução de problemas. As soluções completas,
com animações dos gráficos, estão no repositório do livro: \url{https://github.com/AleffAzevedo/calculo-1-pt/tree/main/exercicios/cap2}.

\begin{enumerate}[label=\textbf{\thechapter.\arabic*}, leftmargin=2.6em, itemsep=6pt]
\item \textbf{Da secante à tangente.} Calcule $\displaystyle\lim_{h\to0}\frac{\frac{1}{a+h}-\frac1a}{h}$, em que $a\neq0$ é constante.
Interprete o resultado em termos de retas secantes ao gráfico de $f(x)=1/x$.
\item \textbf{O limite de $\sin(ax)/x$.} (a) Use tabelas com $x=\pm0{,}1;\ \pm0{,}01;\ \pm0{,}001$ para estimar
$\lim_{x\to0}\frac{\sin2x}{x}$ e $\lim_{x\to0}\frac{\sin3x}{x}$.
(b) Conjecture o valor de $\lim_{x\to0}\frac{\sin ax}{x}$ para $a>0$ e prove-o com o
Teorema~2.7.
\item \textbf{Conjugado com raiz.} Mostre que a substituição direta leva a uma forma indeterminada e calcule
$\displaystyle\lim_{t\to9}\frac{t-9}{\sqrt t-3}$.
\item \textbf{Fatorando cubos.} Calcule $\displaystyle\lim_{x\to1}\frac{x^3-1}{x^2-1}$.
\item \textbf{Limites laterais infinitos.} Calcule $\displaystyle\lim_{x\to1^-}\frac{2x^2+7x-4}{x^2+x-2}$ e
$\displaystyle\lim_{x\to1^+}\frac{2x^2+7x-4}{x^2+x-2}$.
\item \textbf{Teorema do Confronto.} Calcule $\displaystyle\lim_{\theta\to0}\theta^2\cos\!\left(\frac1\theta\right)$.
\item \textbf{Escolhendo $k$ para a continuidade (I).} Encontre $k$ para que $f(x)=\begin{cases}3x+2, & x<k\\ 2x-3, & k\le x\le8\end{cases}$
seja contínua em todo o seu domínio.
\item \textbf{Escolhendo $k$ para a continuidade (II).} Encontre $k$ para que $f(x)=\begin{cases}e^{kx}, & 0\le x<4\\ x+3, & 4\le x\le8\end{cases}$
seja contínua em $[0,8]$.
\item \textbf{TVI: duas partículas.} Uma partícula se move numa reta com posição $s(t)$ contínua, e $s(2)=5$, $s(5)=2$. Outra
partícula tem posição $h(t)=s(t)-t$. Explique por que existe um instante $t\in[2,5]$ em que
$h(t)=0$. O que isso significa para a primeira partícula?
\item \textbf{TVI: localizando uma raiz.} Use o TVI para decidir se $2^x=x^3$ tem solução em $[1{,}25;\,1{,}375]$ ou em
$[1{,}375;\,1{,}5]$. Justifique para cada intervalo.
\item \textbf{Prova $\varepsilon$-$\delta$.} Use a definição precisa de limite para provar que $\displaystyle\lim_{x\to2}(x^2+2x)=8$.
\item \textbf{Provando que um limite não existe.} Use a definição precisa para provar que $\displaystyle\lim_{x\to1}\frac{|x-1|}{x-1}$ não existe.
\end{enumerate}

\subsection*{Respostas}
{\small\begin{enumerate}[label=\textbf{\thechapter.\arabic*}, leftmargin=2.6em, itemsep=2pt]
\item $-\dfrac{1}{a^2}$: é a inclinação da reta tangente a $y=1/x$ em $x=a$.
\item (a) 2 e 3; (b) $\lim_{x\to0}\frac{\sin ax}{x}=a$.
\item $0/0$; o limite é 6.
\item $\dfrac32$.
\item $-\infty$ pela esquerda e $+\infty$ pela direita.
\item 0, pelo Teorema do Confronto.
\item $k=-5$.
\item $k=\dfrac{\ln7}{4}\approx0{,}4865$.
\item $h(2)=3>0$ e $h(5)=-3<0$; pelo TVI, $h(c)=0$ para algum $c\in[2,5]$, ou seja, $s(c)=c$.
\item Com $g(x)=2^x-x^3$: $g(1{,}25)>0$ e $g(1{,}375)<0$, logo há raiz no 1º intervalo; no 2º, $g$ não muda de sinal e o TVI nada garante.
\item Escolha $\delta=\min\{1,\varepsilon/7\}$.
\item Tome $\varepsilon=\tfrac12$: perto de 1 a função vale $-1$ e $1$, que não cabem juntos em nenhuma faixa de largura 1.
\end{enumerate}}


% =====================================================================
\chapter{Derivadas}

No capítulo anterior usamos limites para calcular a inclinação da reta tangente e a
velocidade instantânea. Essas duas ideias são casos de um único conceito: a
\emph{derivada}, que mede a taxa de variação instantânea de uma função. Neste capítulo
definimos a derivada, aprendemos regras para calculá-la sem recorrer ao limite a cada vez e
a aplicamos a movimento, crescimento populacional, custo e receita. Ao final teremos as
derivadas de todas as funções do capítulo~1.

% ---------------------------------------------------------------------
\section{Definindo a derivada}

\subsection{Retas tangentes}

Seja $f$ definida em um intervalo $I$ que contém $a$. Para $x\neq a$ em $I$, o
\textbf{quociente de diferenças}\index{quociente de diferencas@quociente de diferenças}
\[
Q=\frac{f(x)-f(a)}{x-a}
\qquad\text{ou, com } x=a+h,\qquad
Q=\frac{f(a+h)-f(a)}{h}
\]
é a inclinação da reta secante pelos pontos $(a,f(a))$ e $(x,f(x))$. Quando $x\to a$
(ou $h\to0$), as secantes se aproximam da reta tangente.

\figtikz{\begin{tikzpicture}
\begin{axis}[painel,width=8cm,height=6cm,xmin=-0.3,xmax=3.2,ymin=-1,ymax=6.5,xtick={1},xticklabels={$a$},ytick=\empty]
\addplot[domain=-0.3:2.5]{x^2};
\foreach \h/\o in {1.4/30,0.8/50,0.4/70}{
  \edef\tmp{\noexpand\addplot[laranja!\o!white,thin,domain=-0.3:3.2]{1+(2+\h)*(x-1)};
  \noexpand\fill[laranja!\o!white] (axis cs:{1+\h},{(1+\h)^2}) circle (1.6pt);}\tmp}
\addplot[laranja,thick,domain=-0.3:3.2]{1+2*(x-1)};
\fill[black] (axis cs:1,1) circle (2pt) node[below right,font=\footnotesize]{$(a,f(a))$};
\node[laranja,font=\footnotesize] at (axis cs:2.75,4.1) {tangente};
\end{axis}\end{tikzpicture}}{As secantes por $(a,f(a))$ e por pontos cada vez mais
próximos (laranja-claro) tendem à reta tangente (laranja).}

\begin{definicao}
Seja $f$ definida em um intervalo aberto contendo $a$. A \textbf{reta tangente} ao gráfico
de $f$ em $a$ é a reta que passa por $(a,f(a))$ com inclinação
\[
m_{\tan}=\lim_{x\to a}\frac{f(x)-f(a)}{x-a}=\lim_{h\to0}\frac{f(a+h)-f(a)}{h},
\]
desde que esse limite exista.
\end{definicao}

\begin{exemplo}{3.1}{Reta tangente}
Encontre a reta tangente a $f(x)=x^2$ em $x=3$ e a $g(x)=1/x$ em $x=2$.
\solucao
$m_{\tan}=\lim_{x\to3}\dfrac{x^2-9}{x-3}=\lim_{x\to3}(x+3)=6$; a tangente é
$y-9=6(x-3)$, ou $y=6x-9$.

Para $g$: $m_{\tan}=\lim_{h\to0}\dfrac{\frac1{2+h}-\frac12}{h}
=\lim_{h\to0}\dfrac{-h}{2(2+h)h}=-\dfrac14$; a tangente é $y-\tfrac12=-\tfrac14(x-2)$, ou
$y=-\tfrac14x+1$.
\end{exemplo}

\subsection{A derivada em um ponto}

A inclinação da tangente é tão importante que recebe um nome próprio.

\begin{definicao}
A \textbf{derivada}\index{derivada} de $f$ em $a$, denotada $f'(a)$, é
\[
f'(a)=\lim_{x\to a}\frac{f(x)-f(a)}{x-a}=\lim_{h\to0}\frac{f(a+h)-f(a)}{h},
\]
desde que o limite exista. Nesse caso, dizemos que $f$ é \textbf{derivável}\index{funcao@função!derivavel@derivável} (ou
diferenciável) em $a$. A derivada $f'(a)$ é a \textbf{taxa de variação instantânea}\index{taxa de variacao@taxa de variação} de $f$
em $a$.
\end{definicao}

\begin{exemplo}{3.2}{Calculando uma derivada pela definição}
Para $f(x)=3x^2-4x+1$, calcule $f'(2)$.
\solucao
\begin{align*}
f'(2)&=\lim_{h\to0}\frac{\big(3(2+h)^2-4(2+h)+1\big)-5}{h}
=\lim_{h\to0}\frac{12+12h+3h^2-8-4h+1-5}{h}\\
&=\lim_{h\to0}\frac{8h+3h^2}{h}=\lim_{h\to0}(8+3h)=8.
\end{align*}
\end{exemplo}

\subsection{Velocidade e taxas de variação}

Se $s(t)$ é a posição de um objeto que se move ao longo de uma reta, a velocidade média em
$[a,a+h]$ é $\frac{s(a+h)-s(a)}{h}$, e a \textbf{velocidade instantânea} em $t=a$ é o
limite dessas médias: $v(a)=s'(a)$. Mais geralmente, para qualquer grandeza $f$, a taxa
média de variação em um intervalo é um quociente de diferenças, e a taxa instantânea é a
derivada.

\begin{exemplo}{3.3}{Velocidade de uma bola em queda}
Uma bola solta do alto de um prédio de 64 pés tem altura $s(t)=-16t^2+64$ (em pés, com $t$
em segundos). Qual é sua velocidade em $t=1$?
\solucao
$s'(1)=\lim_{t\to1}\dfrac{-16t^2+64-48}{t-1}=\lim_{t\to1}\dfrac{-16(t^2-1)}{t-1}
=\lim_{t\to1}-16(t+1)=-32$ pés/s. O sinal negativo indica que a bola está descendo.
\end{exemplo}

% ---------------------------------------------------------------------
\section{A derivada como função}

\begin{definicao}
A \textbf{função derivada} $f'$ é a função cujo domínio é formado pelos $x$ para os quais
existe
\[
f'(x)=\lim_{h\to0}\frac{f(x+h)-f(x)}{h}.
\]
Também escrevemos $\dfrac{dy}{dx}$, $\dfrac{df}{dx}$, $\dfrac{d}{dx}\big(f(x)\big)$ ou $y'$
(\emph{notação de Leibniz} e \emph{de Lagrange}).
\end{definicao}

\begin{exemplo}{3.4}{Derivada da raiz quadrada}
Para $f(x)=\sqrt x$:
\[
f'(x)=\lim_{h\to0}\frac{\sqrt{x+h}-\sqrt x}{h}\cdot\frac{\sqrt{x+h}+\sqrt x}{\sqrt{x+h}+\sqrt x}
=\lim_{h\to0}\frac{h}{h(\sqrt{x+h}+\sqrt x)}=\frac{1}{2\sqrt x}.
\]
O domínio de $f'$ é $(0,\infty)$, menor que o de $f$: em $x=0$ a tangente é vertical.
\end{exemplo}

\subsection{O gráfico da derivada}

O valor de $f'(x)$ é a inclinação do gráfico de $f$ em $x$. Assim:
\begin{itemize}
\item onde $f$ é crescente, $f'>0$; onde $f$ é decrescente, $f'<0$;
\item onde o gráfico de $f$ tem tangente horizontal, $f'=0$.
\end{itemize}

\figtikz{\begin{tikzpicture}
\begin{groupplot}[group style={group size=1 by 2,vertical sep=1.1cm},painel,width=9cm,height=4.4cm,
  xmin=-2.4,xmax=2.4,xtick={-1,1},ytick={-2,2}]
\nextgroupplot[ymin=-3,ymax=3,title={$f(x)=x^3-3x$}]
\addplot[domain=-2.4:2.4]{x^3-3*x};
\draw[laranja,thick] (axis cs:-1.6,2) -- (axis cs:-0.4,2);
\draw[laranja,thick] (axis cs:0.4,-2) -- (axis cs:1.6,-2);
\draw[cinza,dotted] (axis cs:-1,-3) -- (axis cs:-1,3); \draw[cinza,dotted] (axis cs:1,-3) -- (axis cs:1,3);
\nextgroupplot[ymin=-4,ymax=8,title={$f'(x)=3x^2-3$}]
\addplot[domain=-2.4:2.4,laranja]{3*x^2-3};
\draw[cinza,dotted] (axis cs:-1,-4) -- (axis cs:-1,8); \draw[cinza,dotted] (axis cs:1,-4) -- (axis cs:1,8);
\end{groupplot}\end{tikzpicture}}{Onde $f$ cresce ($|x|>1$), $f'>0$; onde $f$ decresce
($|x|<1$), $f'<0$; nos pontos de tangente horizontal ($x=\pm1$), $f'=0$.}

\subsection{Derivabilidade e continuidade}

\begin{teorema}[Derivabilidade implica continuidade]{3.1}\index{derivabilidade e continuidade}
Se $f$ é derivável em $a$, então $f$ é contínua em $a$.
\end{teorema}
\emph{Prova.} Para $x\neq a$,
$f(x)-f(a)=\dfrac{f(x)-f(a)}{x-a}\,(x-a)\ \to\ f'(a)\cdot0=0$ quando $x\to a$.
Logo $\lim_{x\to a}f(x)=f(a)$. $\square$

A recíproca é falsa: uma função pode ser contínua em $a$ e não ser derivável ali. Há três
situações típicas:

\figtikz{\begin{tikzpicture}
\begin{groupplot}[group style={group size=3 by 1,horizontal sep=1.0cm},painel,width=5.3cm,height=4.4cm,
  xmin=-1.5,xmax=1.5,xtick=\empty,ytick=\empty]
\nextgroupplot[title={(a) bico: $|x|$},ymin=-0.4,ymax=1.5]
\addplot[domain=-1.5:1.5]{abs(x)};
\nextgroupplot[title={(b) tangente vertical: $\sqrt[3]{x}$},ymin=-1.3,ymax=1.3]
\addplot[domain=0:1.5,samples=200]{x^(1/3)}; \addplot[domain=-1.5:0,samples=200]{-(-x)^(1/3)};
\draw[laranja,dashed] (axis cs:0,-1.3) -- (axis cs:0,1.3);
\nextgroupplot[title={(c) oscilação: $x\sin\frac1x$},ymin=-0.5,ymax=0.5,xmin=-0.5,xmax=0.5]
\addplot[domain=-0.5:-0.002,samples=1500,thin]{x*sin(deg(1/x))};
\addplot[domain=0.002:0.5,samples=1500,thin]{x*sin(deg(1/x))};
\addplot[laranja,dashed,thin,domain=-0.5:0.5]{x}; \addplot[laranja,dashed,thin,domain=-0.5:0.5]{-x};
\end{groupplot}\end{tikzpicture}}{Funções contínuas em $0$, mas não deriváveis em $0$.
(a) Os limites laterais do quociente de diferenças são $-1$ e $1$. (b) O quociente tende a
$+\infty$. (c) O quociente $\sin(1/h)$ oscila e não tem limite.}

\begin{exemplo}{3.5}{Uma função por partes derivável}
Encontre $a$ e $b$ para que $f(x)=\begin{cases}ax+b,&x<3\\x^2,&x\ge3\end{cases}$ seja
derivável em todo $\R$.
\solucao
Contínua em 3: $3a+b=9$. Para a derivada existir em 3, os quocientes laterais devem ter o
mesmo limite. À direita, a inclinação de $x^2$ em 3 é 6; à esquerda, é $a$. Logo $a=6$ e
$b=-9$. A reta $y=6x-9$ é exatamente a tangente à parábola em $x=3$.
\end{exemplo}

\subsection{Derivadas de ordem superior}

A derivada de $f'$ é a \textbf{segunda derivada}\index{derivada!de ordem superior}, $f''$ ou $\dfrac{d^2y}{dx^2}$; depois vêm
$f'''$, $f^{(4)},\dots$ Se $s(t)$ é a posição, $v(t)=s'(t)$ é a velocidade e
$a(t)=v'(t)=s''(t)$ é a \textbf{aceleração}\index{aceleracao@aceleração}.

% ---------------------------------------------------------------------
\section{Regras de derivação}

Calcular derivadas pela definição é trabalhoso. As regras a seguir, provadas a partir da
definição, tornam o cálculo mecânico.

\begin{teorema}[Constante e potência]{3.2}
Se $c$ é constante, $\dfrac{d}{dx}(c)=0$. Para todo inteiro $n\ge1$,
$\dfrac{d}{dx}\,x^n=nx^{n-1}$.
\end{teorema}
\emph{Prova da regra da potência.} Pelo binômio de Newton,
$(x+h)^n=x^n+nx^{n-1}h+\binom n2x^{n-2}h^2+\dots+h^n$. Então
\[
\frac{(x+h)^n-x^n}{h}=nx^{n-1}+\binom n2x^{n-2}h+\dots+h^{n-1}\ \longrightarrow\ nx^{n-1}. \qquad\square
\]

\begin{teorema}[Soma, diferença e múltiplo constante]{3.3}
Se $f$ e $g$ são deriváveis e $k$ é constante:
\[
(f\pm g)'=f'\pm g',\qquad (kf)'=k\,f'.
\]
\end{teorema}

Com isso, derivamos qualquer polinômio termo a termo: por exemplo,
$\dfrac{d}{dx}\big(2x^5+7x^2-x+4\big)=10x^4+14x-1$.

\begin{teorema}[Regra do produto]{3.4}\index{regra do produto}
Se $f$ e $g$ são deriváveis, $\big(f\,g\big)'=f'\,g+f\,g'$.
\end{teorema}
\emph{Prova.} Somamos e subtraímos $f(x+h)g(x)$ no numerador:
\[
\frac{f(x+h)g(x+h)-f(x)g(x)}{h}
=f(x+h)\,\frac{g(x+h)-g(x)}{h}+g(x)\,\frac{f(x+h)-f(x)}{h}
\ \longrightarrow\ f(x)g'(x)+g(x)f'(x),
\]
usando que $f(x+h)\to f(x)$, pois $f$ é derivável, logo contínua. $\square$

\begin{nota}
A derivada do produto \emph{não} é o produto das derivadas: $(x\cdot x)'=2x$, e não
$1\cdot1$.
\end{nota}

\begin{teorema}[Regra do quociente]{3.5}\index{regra do quociente}
Se $f$ e $g$ são deriváveis e $g(x)\neq0$,
\[
\Big(\frac fg\Big)'=\frac{f'\,g-f\,g'}{g^2}.
\]
\end{teorema}

Aplicando a regra do quociente a $x^{-n}=1/x^n$, obtemos
$\dfrac{d}{dx}x^{-n}=\dfrac{0-nx^{n-1}}{x^{2n}}=-nx^{-n-1}$: a regra da potência vale também
para expoentes inteiros negativos.

\begin{exemplo}{3.6}{Produto e quociente}
(a) $\dfrac{d}{dx}\big[(x^2+2)(3x^3-5x)\big]=2x(3x^3-5x)+(x^2+2)(9x^2-5)=15x^4+3x^2-10$.

(b) $\dfrac{d}{dx}\Big[\dfrac{5x^2}{4x+3}\Big]=\dfrac{10x(4x+3)-5x^2\cdot4}{(4x+3)^2}
=\dfrac{20x^2+30x}{(4x+3)^2}$.

(c) $\dfrac{d}{dx}\Big[\dfrac{6}{x^2}\Big]=\dfrac{d}{dx}\big[6x^{-2}\big]=-12x^{-3}=-\dfrac{12}{x^3}$.
\end{exemplo}

\begin{exemplo}{3.7}{Tangentes horizontais}
Onde o gráfico de $f(x)=x^3-6x^2+9x+1$ tem tangente horizontal?
\solucao
$f'(x)=3x^2-12x+9=3(x-1)(x-3)=0$ em $x=1$ e $x=3$. Os pontos são $(1,5)$ e $(3,1)$.
\end{exemplo}

% ---------------------------------------------------------------------
\section{Derivadas como taxas de variação}

\textbf{Fórmula da variação aproximada.} Como $f'(a)\approx\frac{f(a+h)-f(a)}h$ para $h$
pequeno,
\[
f(a+h)\approx f(a)+f'(a)\,h .
\]
Por exemplo, se $f(3)=2$ e $f'(3)=5$, então $f(3{,}2)\approx2+5(0{,}2)=3$.

\subsection{Movimento retilíneo}

\begin{definicao}
Se $s(t)$ é a posição de um objeto no instante $t$: a \textbf{velocidade} é
$v(t)=s'(t)$; a \textbf{rapidez} (módulo da velocidade) é $|v(t)|$; a
\textbf{aceleração} é $a(t)=v'(t)=s''(t)$.
\end{definicao}

O objeto se move no sentido positivo quando $v>0$, no negativo quando $v<0$, e está
momentaneamente em repouso quando $v=0$. Ele acelera (a rapidez aumenta) quando $v$ e $a$
têm o mesmo sinal, e desacelera quando têm sinais opostos.

\begin{exemplo}{3.8}{Posição e velocidade}
Uma partícula tem posição $s(t)=t^3-9t^2+24t$ ($t\ge0$). Quando ela está em repouso? Em
que intervalos se move para a frente e para trás?
\solucao
$v(t)=3t^2-18t+24=3(t-2)(t-4)$. Repouso em $t=2$ e $t=4$. Pelo sinal de $v$: para a frente
em $[0,2)$ e $(4,\infty)$, para trás em $(2,4)$. A aceleração é $a(t)=6t-18$, nula em $t=3$.
\end{exemplo}

\subsection{Populações, custo e receita}

Se $P(t)$ é uma população, $P'(t)$ é a \textbf{taxa de crescimento populacional}. Na
economia, se $C(x)$ é o custo de produzir $x$ unidades e $R(x)$ é a receita com a venda de
$x$ unidades, o \textbf{custo marginal}\index{analise marginal@análise marginal} é $C'(x)$, a \textbf{receita marginal} é $R'(x)$
e o \textbf{lucro marginal} é $P'(x)=R'(x)-C'(x)$. Pela fórmula da variação aproximada, o
custo marginal $C'(x)$ estima o custo de produzir a unidade número $x+1$.

\begin{exemplo}{3.9}{Receita marginal}
Se $R(x)=4x-0{,}01x^2$ (em mil reais), então $R'(x)=4-0{,}02x$. Para $x=100$, $R'(100)=2$:
vender a 101ª unidade aumenta a receita em cerca de R\$\,2 mil. O valor exato é
$R(101)-R(100)=1{,}99$.
\end{exemplo}

% ---------------------------------------------------------------------
\section{Derivadas das funções trigonométricas}

\begin{teorema}[Derivadas de seno e cosseno]{3.6}
\[
\frac{d}{dx}(\sin x)=\cos x,\qquad \frac{d}{dx}(\cos x)=-\sin x .
\]
\end{teorema}
\emph{Prova para o seno.} Pela fórmula da soma,
\[
\frac{\sin(x+h)-\sin x}{h}=\sin x\cdot\frac{\cos h-1}{h}+\cos x\cdot\frac{\sin h}{h}
\ \longrightarrow\ \sin x\cdot0+\cos x\cdot1=\cos x,
\]
usando os dois limites trigonométricos do Teorema~2.7. A prova para o cosseno é
análoga. $\square$

\figtikz{\begin{tikzpicture}
\begin{axis}[painel,width=11cm,height=4.4cm,xmin=-6.6,xmax=6.6,ymin=-1.5,ymax=1.5,
  xtick={-6.2832,-3.1416,3.1416,6.2832},xticklabels={$-2\pi$,$-\pi$,$\pi$,$2\pi$},ytick={-1,1},
  legend style={font=\footnotesize,at={(1,1.05)},anchor=south east},legend columns=2]
\addplot[domain=-6.6:6.6]{sin(deg(x))}; \addlegendentry{$y=\sin x$}
\addplot[laranja,domain=-6.6:6.6]{cos(deg(x))}; \addlegendentry{$y=\cos x$ (derivada)}
\foreach \p in {-4.7124,-1.5708,1.5708,4.7124}{\edef\tmp{\noexpand\fill[black] (axis cs:\p,0) circle (1.6pt);}\tmp}
\end{axis}\end{tikzpicture}}{Onde $\sin x$ tem tangente horizontal (máximos e mínimos), sua
derivada $\cos x$ se anula (pontos pretos); onde $\sin x$ cresce mais rápido ($x=0$), $\cos x$
vale 1.}

Pela regra do quociente, $\dfrac{d}{dx}\tan x=\dfrac{\cos x\cos x-\sin x(-\sin x)}{\cos^2x}
=\dfrac1{\cos^2x}=\sec^2x$. Do mesmo modo:

\begin{teorema}[Demais funções trigonométricas]{3.7}
\[
(\tan x)'=\sec^2x,\qquad (\cot x)'=-\csc^2x,\qquad (\sec x)'=\sec x\tan x,\qquad
(\csc x)'=-\csc x\cot x .
\]
\end{teorema}

\textbf{Derivadas de ordem superior.} As derivadas de $\sin x$ se repetem com período 4:
$\sin x\to\cos x\to-\sin x\to-\cos x\to\sin x$. Assim,
$\dfrac{d^{74}}{dx^{74}}\sin x=\dfrac{d^2}{dx^2}\sin x=-\sin x$, pois $74=4\cdot18+2$.

\begin{exemplo}{3.10}{Uma massa presa a uma mola}
A posição de uma massa é $s(t)=2\sin t$ (cm). Então $v(t)=2\cos t$ e $a(t)=-2\sin t=-s(t)$:
a aceleração é proporcional ao deslocamento e tem sentido oposto. Em $t=\pi/2$, a massa está
em $s=2$, parada ($v=0$), com aceleração $-2$ cm/s$^2$, que a puxa de volta.
\end{exemplo}

% ---------------------------------------------------------------------
\section{A regra da cadeia}

Como derivar $h(x)=\sin(x^3)$ ou $(3x^2+1)^{20}$? São funções compostas, e sua derivada
segue uma regra simples: se $u=g(x)$ varia $g'(x)$ vezes mais rápido que $x$, e $y=f(u)$
varia $f'(u)$ vezes mais rápido que $u$, então $y$ varia $f'(u)\,g'(x)$ vezes mais rápido que
$x$.

\begin{regra}{Regra da cadeia}\index{regra da cadeia}
Se $g$ é derivável em $x$ e $f$ é derivável em $g(x)$, então $h=f\circ g$ é derivável em
$x$ e
\[
h'(x)=f'\big(g(x)\big)\,g'(x).
\]
Na notação de Leibniz, com $y=f(u)$ e $u=g(x)$:\qquad
$\dfrac{dy}{dx}=\dfrac{dy}{du}\cdot\dfrac{du}{dx}$.
\end{regra}

Em palavras: \emph{derive a função de fora, mantendo a de dentro, e multiplique pela
derivada da de dentro}. Dois casos frequentes:
\[
\frac{d}{dx}\big(g(x)\big)^n=n\big(g(x)\big)^{n-1}g'(x),\qquad
\frac{d}{dx}\sin\big(g(x)\big)=\cos\big(g(x)\big)\,g'(x),
\]
e analogamente $\frac{d}{dx}\cos u=-\sin u\,u'$, $\frac{d}{dx}\tan u=\sec^2u\,u'$, etc.

\begin{exemplo}{3.11}{Usando a regra da cadeia}
(a) $\dfrac{d}{dx}(3x^2+1)^3=3(3x^2+1)^2\cdot6x=18x(3x^2+1)^2$.

(b) $\dfrac{d}{dx}\sin^3x=3\sin^2x\cos x$.

(c) $\dfrac{d}{dx}\cos(5x^2)=-\sin(5x^2)\cdot10x$.

(d) $\dfrac{d}{dx}\Big[\dfrac{1}{(3x^2+1)^2}\Big]=\dfrac{d}{dx}(3x^2+1)^{-2}=-2(3x^2+1)^{-3}\cdot6x
=-\dfrac{12x}{(3x^2+1)^3}$.
\end{exemplo}

\textbf{Composição de três funções.} Para $k(x)=h\big(f(g(x))\big)$, aplicamos a regra duas
vezes: $k'(x)=h'\big(f(g(x))\big)\cdot f'\big(g(x)\big)\cdot g'(x)$. Por exemplo,
\[
\frac{d}{dx}\sin^4(7x-3)=4\sin^3(7x-3)\cdot\cos(7x-3)\cdot7 .
\]

\begin{exemplo}{3.12}{Cadeia e produto juntos}
$\dfrac{d}{dx}\big[x^2\cos(3x)\big]=2x\cos(3x)+x^2\cdot\big(-\sin(3x)\cdot3\big)
=2x\cos(3x)-3x^2\sin(3x)$.
\end{exemplo}

% ---------------------------------------------------------------------
\section{Derivadas de funções inversas}

Se $(a,b)$ está no gráfico de $f$, então $(b,a)$ está no gráfico de $f^{-1}$, que é a
reflexão em torno de $y=x$. A reflexão troca $\Delta x$ e $\Delta y$; logo, as inclinações
das tangentes em pontos correspondentes são \emph{recíprocas}.

\figtikz{\begin{tikzpicture}
\begin{axis}[painel,width=7cm,height=7cm,xmin=-2.2,xmax=4.2,ymin=-2.2,ymax=4.2,axis equal image,xtick=\empty,ytick=\empty]
\addplot[cinza,dotted,domain=-2.2:4.2]{x};
\addplot[domain=-2.2:1.45]{exp(x)}; \addplot[laranja,domain=0.12:4.2]{ln(x)};
\addplot[azul,dashed,thin,domain=-0.6:2.2]{exp(0.8)+exp(0.8)*(x-0.8)};
\addplot[laranja,dashed,thin,domain=0.6:4.2]{0.8+(x-exp(0.8))/exp(0.8)};
\fill[azul] (axis cs:0.8,{exp(0.8)}) circle (2pt) node[left,font=\footnotesize]{$(a,b)$};
\fill[laranja] (axis cs:{exp(0.8)},0.8) circle (2pt) node[below right,font=\footnotesize]{$(b,a)$};
\node[azul,font=\footnotesize] at (axis cs:-1.2,1.4) {$y=e^x$};
\node[laranja,font=\footnotesize] at (axis cs:3.4,0.7) {$y=\ln x$};
\end{axis}\end{tikzpicture}}{As tangentes a $f$ em $(a,b)$ e a $f^{-1}$ em $(b,a)$ são
reflexões uma da outra; suas inclinações são $m$ e $1/m$.}

\begin{teorema}[Derivada da função inversa]{3.8}
Seja $f$ inversível e derivável. Se $f'\big(f^{-1}(x)\big)\neq0$, então $f^{-1}$ é derivável em
$x$ e
\[
\big(f^{-1}\big)'(x)=\frac{1}{f'\big(f^{-1}(x)\big)} .
\]
\end{teorema}

\begin{exemplo}{3.13}{Aplicando o teorema da inversa}
Para $f(x)=x^3+x+1$, que é crescente, calcule $(f^{-1})'(3)$.
\solucao
$f(1)=3$, logo $f^{-1}(3)=1$. Como $f'(x)=3x^2+1$, temos $f'(1)=4$ e
$(f^{-1})'(3)=\dfrac1{f'(1)}=\dfrac14$. Não foi preciso encontrar a fórmula de $f^{-1}$.
\end{exemplo}

\textbf{Expoentes racionais.} Aplicando o teorema a $g(x)=x^{1/n}$, inversa de $x^n$,
obtemos $g'(x)=\frac1n x^{1/n-1}$. Combinando com a regra da cadeia, para todo racional $r$:
\[
\frac{d}{dx}x^r=r\,x^{r-1}.
\]
Por exemplo, $\dfrac{d}{dx}\sqrt[3]{x^2}=\dfrac{d}{dx}x^{2/3}=\dfrac23x^{-1/3}$.

\textbf{Inversas trigonométricas.} Para $y=\sin^{-1}x$, temos $\sin y=x$ com
$-\frac\pi2\le y\le\frac\pi2$. Pelo teorema, $\dfrac{dy}{dx}=\dfrac1{\cos y}$. Como
$\cos y\ge0$ nesse intervalo, $\cos y=\sqrt{1-\sin^2y}=\sqrt{1-x^2}$. Do mesmo modo:

\begin{teorema}[Derivadas das inversas trigonométricas]{3.9}
\[
\frac{d}{dx}\sin^{-1}x=\frac{1}{\sqrt{1-x^2}},\qquad
\frac{d}{dx}\cos^{-1}x=\frac{-1}{\sqrt{1-x^2}},\qquad
\frac{d}{dx}\tan^{-1}x=\frac{1}{1+x^2},
\]
\[
\frac{d}{dx}\cot^{-1}x=\frac{-1}{1+x^2},\qquad
\frac{d}{dx}\sec^{-1}x=\frac{1}{|x|\sqrt{x^2-1}},\qquad
\frac{d}{dx}\csc^{-1}x=\frac{-1}{|x|\sqrt{x^2-1}} .
\]
\end{teorema}

\begin{exemplo}{3.14}{Cadeia com inversa trigonométrica}
$\dfrac{d}{dx}\tan^{-1}(x^2)=\dfrac{2x}{1+x^4}$\qquad e\qquad
$\dfrac{d}{dx}\big[x\sin^{-1}x\big]=\sin^{-1}x+\dfrac{x}{\sqrt{1-x^2}}$.
\end{exemplo}

% ---------------------------------------------------------------------
\section{Derivação implícita}

Muitas curvas, como a circunferência $x^2+y^2=25$, não são gráficos de uma única função
$y=f(x)$. Ainda assim, perto de cada ponto (salvo exceções) a curva é o gráfico de alguma
função $y(x)$, e podemos derivar a equação \emph{como está}, tratando $y$ como função de $x$.

\begin{estrategia}{derivação implícita}\index{derivacao implicita@derivação implícita}
\begin{enumerate}[label=\arabic*.]
\item Derive os dois lados da equação em relação a $x$, lembrando que, pela regra da
  cadeia, $\frac{d}{dx}\big(g(y)\big)=g'(y)\,\frac{dy}{dx}$ (por exemplo,
  $\frac{d}{dx}y^2=2y\,y'$).
\item Passe os termos com $\frac{dy}{dx}$ para um lado e os demais para o outro.
\item Ponha $\frac{dy}{dx}$ em evidência e isole.
\end{enumerate}
\end{estrategia}

\begin{exemplo}{3.15}{Tangente a uma circunferência}
Na circunferência $x^2+y^2=25$: $2x+2y\,y'=0$, logo $y'=-\dfrac xy$. No ponto $(3,-4)$,
$y'=\frac34$, e a tangente é $y+4=\tfrac34(x-3)$. Note que a tangente é perpendicular ao raio,
cuja inclinação é $-\tfrac43$.
\end{exemplo}

\begin{exemplo}{3.16}{O fólio de Descartes}
Encontre a tangente à curva $x^3+y^3=6xy$ no ponto $(3,3)$.
\solucao
Derivando: $3x^2+3y^2y'=6y+6x\,y'$. Isolando,
\[
y'=\frac{2y-x^2}{y^2-2x}.
\]
Em $(3,3)$: $y'=\frac{6-9}{9-6}=-1$. A tangente é $y=-x+6$.
\figtikz{\begin{tikzpicture}
\begin{axis}[painel,width=6.2cm,height=6.2cm,xmin=-4,xmax=5,ymin=-4,ymax=5,axis equal image,xtick={3},ytick={3}]
\addplot[domain=-0.55:12,samples=400,variable=t] ({6*t/(1+t^3)},{6*t^2/(1+t^3)});
\addplot[domain=-50:-1.6,samples=300,variable=t] ({6*t/(1+t^3)},{6*t^2/(1+t^3)});
\addplot[laranja,domain=0.5:5.5]{-x+6};
\fill[laranja] (axis cs:3,3) circle (2pt);
\end{axis}\end{tikzpicture}}{O fólio de Descartes e a tangente em $(3,3)$.}
\end{exemplo}

% ---------------------------------------------------------------------
\section{Derivadas das funções exponenciais e logarítmicas}

Vimos no capítulo~1 que $e$ é a base cuja exponencial tem tangente de inclinação 1 em
$x=0$: $\lim_{h\to0}\frac{e^h-1}{h}=1$. Daí
\[
\frac{e^{x+h}-e^x}{h}=e^x\cdot\frac{e^h-1}{h}\ \longrightarrow\ e^x .
\]

\begin{teorema}[Derivada da exponencial natural]{3.10}
$\dfrac{d}{dx}e^x=e^x$, e, pela regra da cadeia, $\dfrac{d}{dx}e^{g(x)}=e^{g(x)}\,g'(x)$.
\end{teorema}

Como $\ln x$ é a inversa de $e^x$, o teorema da inversa dá
$\frac{d}{dx}\ln x=\frac{1}{e^{\ln x}}=\frac1x$.

\begin{teorema}[Derivada do logaritmo natural]{3.11}
Para $x>0$, $\dfrac{d}{dx}\ln x=\dfrac1x$, e $\dfrac{d}{dx}\ln g(x)=\dfrac{g'(x)}{g(x)}$
quando $g(x)>0$. Mais geralmente, $\dfrac{d}{dx}\ln|x|=\dfrac1x$ para $x\neq0$.
\end{teorema}

Para outras bases, escrevemos $b^x=e^{x\ln b}$ e $\log_bx=\frac{\ln x}{\ln b}$:

\begin{teorema}[Bases quaisquer]{3.12}
Para $b>0$, $b\neq1$:\qquad
$\dfrac{d}{dx}b^x=b^x\ln b$,\qquad $\dfrac{d}{dx}\log_bx=\dfrac{1}{x\ln b}$.
\end{teorema}

\begin{exemplo}{3.17}{Exponenciais e logaritmos}
(a) $\dfrac{d}{dx}\Big[\dfrac{e^{x^2}}{x}\Big]=\dfrac{2xe^{x^2}\cdot x-e^{x^2}}{x^2}
=\dfrac{e^{x^2}(2x^2-1)}{x^2}$.

(b) $\dfrac{d}{dx}\ln(x^3+3x-4)=\dfrac{3x^2+3}{x^3+3x-4}$.

(c) $\dfrac{d}{dx}\ln\Big(\dfrac{x^2\sin x}{2x+1}\Big)$: primeiro usamos as propriedades do
logaritmo, $2\ln x+\ln\sin x-\ln(2x+1)$; a derivada é $\dfrac2x+\cot x-\dfrac{2}{2x+1}$.

(d) $\dfrac{d}{dx}\,3^{x^2}=3^{x^2}\ln3\cdot2x$.
\end{exemplo}

\subsection{Derivação logarítmica}

Para funções da forma $f(x)^{g(x)}$, ou produtos e quocientes complicados, tomamos o
logaritmo dos dois lados antes de derivar.

\begin{estrategia}{derivação logarítmica}\index{derivacao logaritmica@derivação logarítmica}
\begin{enumerate}[label=\arabic*.]
\item Aplique $\ln$ aos dois lados de $y=h(x)$ e simplifique com as propriedades do
  logaritmo.
\item Derive implicitamente: o lado esquerdo vira $\frac{y'}{y}$.
\item Multiplique por $y$ e substitua $y=h(x)$.
\end{enumerate}
\end{estrategia}

\begin{exemplo}{3.18}{A derivada de $x^x$}
Seja $y=x^x$ ($x>0$). Então $\ln y=x\ln x$, e derivando,
$\dfrac{y'}{y}=\ln x+x\cdot\dfrac1x=\ln x+1$. Logo
\[
\frac{d}{dx}x^x=x^x(\ln x+1).
\]
Atenção: $x^x$ não é potência ($x^n$) nem exponencial ($b^x$), e nenhuma das duas regras
se aplica diretamente.
\end{exemplo}

Com a mesma técnica, para qualquer real $r$ e $x>0$, $y=x^r$ dá $\ln y=r\ln x$ e
$y'=y\cdot\frac rx=rx^{r-1}$: a regra da potência vale para \emph{todo} expoente real.

% ---------------------------------------------------------------------
\section*{Resumo do capítulo}\markright{Resumo do capítulo}
\addcontentsline{toc}{section}{Resumo do capítulo}
\begin{itemize}
\item $f'(a)=\lim_{h\to0}\frac{f(a+h)-f(a)}h$: inclinação da tangente e taxa de variação
  instantânea. Derivável implica contínua, mas não o contrário (bicos, tangentes verticais,
  oscilações).
\item Regras: $(x^r)'=rx^{r-1}$; $(f\pm g)'=f'\pm g'$; $(fg)'=f'g+fg'$;
  $(f/g)'=(f'g-fg')/g^2$; $(f\circ g)'=f'(g)\,g'$.
\item Trigonométricas: $(\sin x)'=\cos x$, $(\cos x)'=-\sin x$, $(\tan x)'=\sec^2x$,
  $(\sec x)'=\sec x\tan x$.
\item Inversas: $(f^{-1})'(x)=1/f'(f^{-1}(x))$; $(\sin^{-1}x)'=1/\sqrt{1-x^2}$,
  $(\tan^{-1}x)'=1/(1+x^2)$.
\item Exponenciais e logaritmos: $(e^x)'=e^x$, $(\ln x)'=1/x$, $(b^x)'=b^x\ln b$.
\item Movimento: $v=s'$, $a=v'=s''$. Economia: custo, receita e lucro marginais são
  derivadas. Variação aproximada: $f(a+h)\approx f(a)+f'(a)h$.
\item Derivação implícita: derive tratando $y$ como função de $x$ e isole $y'$.
  Derivação logarítmica para $f(x)^{g(x)}$ e produtos complicados.
\end{itemize}

% Gerado por ferramentas/gerar.py; não edite à mão.
\section*{Exercícios}\markright{Exercícios}
\addcontentsline{toc}{section}{Exercícios}
Os pontos antes de cada número indicam a dificuldade: \nivel{1} fácil, \nivel{2} médio, \nivel{3} difícil e \nivel{4} desafio. Os exercícios marcados com \estrela{} têm solução completa, com animação do gráfico, no repositório do livro: \url{https://github.com/AleffAzevedo/calculo-1-pt/blob/main/exercicios/cap3/README.md}. As respostas estão no fim do capítulo.

\begin{itemize}[leftmargin=5.2em,labelsep=0.6em,itemsep=4pt]
\item[\exnum{1}{3.1}] Use a definição de derivada para calcular $f'(x)$ para $f(x)=x^2-3x$.
\item[\exnum{1}{3.2}] Use a definição para calcular $f'(x)$ para $f(x)=\dfrac1{x+1}$.
\item[\exnum{1}{3.3}] Use a definição para calcular $f'(x)$ para $f(x)=\sqrt{2x+1}$. (Dica: conjugado.)
\item[\exnum{1}{3.4}] Encontre a reta tangente a $y=x^3$ em $x=2$.
\item[\exnum{1}{3.5}] Derive $y=4x^5-\dfrac3{x^2}+2\sqrt x-7$.
\item[\exnum{1}{3.6}] Derive $y=(x^2+1)(x^3-2x)$ de duas formas: pela regra do produto e expandindo antes.
\item[\exnum{1}{3.7}] Derive $y=\dfrac{2x-1}{x+3}$.
\item[\exnum{1}{3.8}] Derive $y=\dfrac{x^2}{x^2+1}$.
\item[\exnum{1}{3.9}] Derive $y=x^2\sen x$.
\item[\exnum{1}{3.10}] Derive $y=\dfrac{\tg x}{x}$.
\item[\exnum{1}{3.11}] Derive $y=(3x^2-5)^4$.
\item[\exnum{1}{3.12}] Derive $y=\sen(4x+1)$ e $y=\cos^2x$.
\item[\exnum{1}{3.13}] Derive $y=\sqrt{1+x^2}$.
\item[\exnum{1}{3.14}] Derive $y=e^{-3x}\cos x$.
\item[\exnum{1}{3.15}] Derive $y=\ln(x^2+4)$ e $y=\ln(\sen x)$.
\item[\exnum{1}{3.16}] Derive $y=x^3\,2^x$.
\item[\exnum{1}{3.17}] Derive $y=\arcsen(2x)$ e $y=\arctg(3x)$.
\item[\exnum{1}{3.18}] Calcule $y''$ para $y=x^4-2x^3+x$ e para $y=\sen(2x)$.
\item[\exnum{1}{3.19}] Para $s(t)=t^3-6t^2+9t$ (m), calcule a velocidade e a aceleração em $t=2$ s.
\item[\exnum{1}{3.20}] Use derivação implícita para encontrar $y'$ em $x^2+xy+y^2=7$.
\item[\exnum{2}{3.21}] \estrela\ \textbf{Retas tangente e normal.} Encontre as retas tangente e normal à parábola $y=x^2-4x+5$ no ponto de abscissa
$x=3$ e esboce as três curvas.
\item[\exnum{2}{3.22}] \estrela\ \textbf{Construindo o gráfico de $f'$.} Para $f(x)=x^3-3x^2$, determine onde $f'(x)>0$, $f'(x)<0$ e $f'(x)=0$, e esboce
o gráfico de $f'$ alinhado com o de $f$.
\item[\exnum{2}{3.23}] \estrela\ \textbf{Função por partes derivável.} Encontre $a$ e $b$ para que $f(x)=\begin{cases}x^2,&x\le1\\ax+b,&x>1\end{cases}$
seja derivável em todo $\mathbb R$.
\item[\exnum{2}{3.24}] Derive $y=(x^2+1)^3(2x-1)^2$ e fatore o resultado.
\item[\exnum{2}{3.25}] Derive $y=\sen^2(3x)$ e simplifique.
\item[\exnum{2}{3.26}] Derive $y=\cos\sqrt x$.
\item[\exnum{2}{3.27}] Derive $y=x\,e^{-x^2}$ e encontre onde $y'=0$.
\item[\exnum{2}{3.28}] Derive $y=\ln(\cos x)$ e $y=\ln|\sec x+\tg x|$.
\item[\exnum{2}{3.29}] Derive $y=\arctg\sqrt x$.
\item[\exnum{2}{3.30}] \estrela\ \textbf{O quociente de diferenças do seno.} Para $h=1$, $0{,}5$ e $0{,}1$, esboce $q_h(x)=\dfrac{\sen(x+h)-\sen x}{h}$ junto com
$\cos x$. O que acontece quando $h\to0$? Prove o resultado.
\item[\exnum{2}{3.31}] \estrela\ \textbf{Regra da cadeia e tangente.} Encontre a reta tangente a $y=\sen(x^2)$ em $x=\sqrt{\pi}$.
\item[\exnum{2}{3.32}] Calcule $y''$ para $y=xe^x$ e deduza uma fórmula para $y^{(n)}$.
\item[\exnum{2}{3.33}] Calcule $\dfrac{d^{50}}{dx^{50}}\cos x$ e $\dfrac{d^{99}}{dx^{99}}\sen x$.
\item[\exnum{2}{3.34}] Encontre os pontos em que $y=\dfrac{x}{x^2+1}$ tem tangente horizontal.
\item[\exnum{2}{3.35}] Use $f(a+h)\approx f(a)+f'(a)h$ para estimar $\sqrt{4{,}1}$ e $\sqrt[3]{8{,}3}$.
\item[\exnum{2}{3.36}] Sabendo que $g(1)=2$, $g'(1)=3$, $f(2)=5$ e $f'(2)=-1$, calcule $h'(1)$ para $h=f\circ g$ e para $k=g\cdot(f\circ g)$.
\item[\exnum{2}{3.37}] \estrela\ \textbf{Derivada de uma inversa sem a fórmula.} Seja $f(x)=x+e^x$. Mostre que $f$ é inversível e calcule $(f^{-1})'(1)$ e a reta
tangente ao gráfico de $f^{-1}$ no ponto $(1,0)$.
\item[\exnum{2}{3.38}] Use as propriedades do logaritmo antes de derivar: $y=\ln\dfrac{x^2+1}{x-1}$.
\item[\exnum{2}{3.39}] \estrela\ \textbf{Custo marginal.} O custo (em mil reais) de produzir $x$ unidades é $C(x)=0{,}01x^3-0{,}6x^2+15x+200$.
Calcule o custo marginal em $x=20$ e compare com o custo real de produzir a 21ª unidade.
\item[\exnum{2}{3.40}] Derive $y=3^{\sen x}$ e $y=\log_2(x^2+1)$.
\item[\exnum{3}{3.41}] \estrela\ \textbf{Partícula em movimento.} Uma partícula tem posição $s(t)=t^3-9t^2+24t$ (metros, $t\ge0$ em segundos).
(a) Quando ela está em repouso? (b) Quando se move para a frente e para trás?
(c) Qual é a distância total percorrida em $[0,5]$? (d) Quando está acelerando?
\item[\exnum{3}{3.42}] \estrela\ \textbf{Tangentes horizontais numa elipse.} Encontre os pontos da curva $x^2-xy+y^2=3$ onde a reta tangente é horizontal e onde é
vertical.
\item[\exnum{3}{3.43}] \estrela\ \textbf{O mínimo de $x^x$.} Use derivação logarítmica para encontrar $\dfrac{d}{dx}x^x$ ($x>0$) e determine o ponto onde
o gráfico de $y=x^x$ tem tangente horizontal.
\item[\exnum{3}{3.44}] \estrela\ \textbf{Tangentes por um ponto externo.} Encontre todas as retas tangentes à parábola $y=x^2$ que passam pelo ponto $(1,-3)$,
que não está na parábola.
\item[\exnum{3}{3.45}] Prove, pela definição, que $\dfrac{d}{dx}\cos x=-\sen x$.
\item[\exnum{3}{3.46}] Deduza $\dfrac{d}{dx}\arctg x=\dfrac1{1+x^2}$ a partir do teorema da derivada da inversa.
\item[\exnum{3}{3.47}] Derive $y=x^{\sen x}$ ($x>0$).
\item[\exnum{3}{3.48}] Use derivação logarítmica: $y=\dfrac{(x+1)^2(x-2)^3}{(x+3)^4}$.
\item[\exnum{3}{3.49}] Para a circunferência $x^2+y^2=r^2$, mostre que $y''=-\dfrac{r^2}{y^3}$.
\item[\exnum{3}{3.50}] Encontre a reta tangente à curva $x^2y+xy^2=6$ no ponto $(1,2)$.
\item[\exnum{3}{3.51}] Onde $f(x)=|x^2-4|$ é derivável? Calcule as derivadas laterais nos pontos problemáticos.
\item[\exnum{3}{3.52}] Mostre que $(\senh x)'=\cosh x$, $(\cosh x)'=\senh x$ e $(\tgh x)'=\operatorname{sech}^2x$.
\item[\exnum{3}{3.53}] Seja $f(x)=x^2\sen\frac1x$ para $x\neq0$ e $f(0)=0$. Calcule $f'(0)$ pela definição e mostre que $f'$ não é contínua em $0$.
\item[\exnum{3}{3.54}] A reta normal a $y=e^x$ em $x=0$ corta o eixo $x$ em que ponto?
\item[\exnum{3}{3.55}] Para que valor de $a$ a reta $y=ax$ é tangente à curva $y=e^x$? Em que ponto?
\item[\exnum{3}{3.56}] Derive $y=\arcsen x+\arccos x$ e conclua uma identidade.
\item[\exnum{3}{3.57}] Mostre que $y=e^{2x}+e^{-3x}$ satisfaz $y''+y'-6y=0$.
\item[\exnum{3}{3.58}] Encontre a reta tangente a $y=x^x$ em $x=e$.
\item[\exnum{3}{3.59}] O ângulo de um pêndulo é $\theta(t)=0{,}2\cos(2\pi t)$ rad. Qual é a velocidade angular máxima e quando ocorre?
\item[\exnum{3}{3.60}] Em que pontos a tangente a $y=x^3$ é paralela à reta $y=12x+1$? Escreva as tangentes.
\item[\exnum{4}{3.61}] Prove que $(fgh)'=f'gh+fg'h+fgh'$ e generalize para $n$ fatores.
\item[\exnum{4}{3.62}] Para que valores de $n$ a função $y=x^n$ satisfaz $x^2y''-3xy'+4y=0$?
\item[\exnum{4}{3.63}] Prove por indução que $\dfrac{d^n}{dx^n}(xe^x)=(x+n)e^x$.
\item[\exnum{4}{3.64}] Mostre que as curvas $x^2-y^2=a$ e $xy=b$ ($a,b\neq0$) se cortam em ângulo reto.
\item[\exnum{4}{3.65}] Encontre todas as retas que são tangentes simultaneamente a $y=x^2$ e a $y=-x^2+4x-4$.
\end{itemize}
\subsection*{Respostas}\markright{Respostas}
{\small\begin{multicols}{2}\raggedright
\begin{itemize}[leftmargin=2.7em,labelsep=0.5em,itemsep=1pt]
\item[\textbf{3.1}] $2x-3$.
\item[\textbf{3.2}] $-\dfrac1{(x+1)^2}$.
\item[\textbf{3.3}] $\dfrac1{\sqrt{2x+1}}$.
\item[\textbf{3.4}] $y=12x-16$.
\item[\textbf{3.5}] $20x^4+\dfrac6{x^3}+\dfrac1{\sqrt x}$.
\item[\textbf{3.6}] $5x^4-3x^2-2$.
\item[\textbf{3.7}] $\dfrac7{(x+3)^2}$.
\item[\textbf{3.8}] $\dfrac{2x}{(x^2+1)^2}$.
\item[\textbf{3.9}] $2x\sen x+x^2\cos x$.
\item[\textbf{3.10}] $\dfrac{x\sec^2x-\tg x}{x^2}$.
\item[\textbf{3.11}] $24x(3x^2-5)^3$.
\item[\textbf{3.12}] $4\cos(4x+1)$; $-2\sen x\cos x=-\sen2x$.
\item[\textbf{3.13}] $\dfrac{x}{\sqrt{1+x^2}}$.
\item[\textbf{3.14}] $-e^{-3x}(3\cos x+\sen x)$.
\item[\textbf{3.15}] $\dfrac{2x}{x^2+4}$; $\cotg x$.
\item[\textbf{3.16}] $2^x\big(3x^2+x^3\ln2\big)$.
\item[\textbf{3.17}] $\dfrac2{\sqrt{1-4x^2}}$; $\dfrac3{1+9x^2}$.
\item[\textbf{3.18}] $12x^2-12x$; $-4\sen2x$.
\item[\textbf{3.19}] $v(2)=-3$ m/s; $a(2)=0$.
\item[\textbf{3.20}] $y'=-\dfrac{2x+y}{x+2y}$.
\item[\textbf{3.21}] Tangente $y=2x-4$; normal $y=-\tfrac12x+\tfrac72$.
\item[\textbf{3.22}] $f'(x)=3x(x-2)$: positiva em $(-\infty,0)\cup(2,\infty)$, negativa em $(0,2)$, nula em $x=0$ e $x=2$.
\item[\textbf{3.23}] $a=2$, $b=-1$.
\item[\textbf{3.24}] $2(x^2+1)^2(2x-1)(8x^2-3x+2)$.
\item[\textbf{3.25}] $6\sen3x\cos3x=3\sen6x$.
\item[\textbf{3.26}] $-\dfrac{\sen\sqrt x}{2\sqrt x}$.
\item[\textbf{3.27}] $e^{-x^2}(1-2x^2)$; $x=\pm\tfrac1{\sqrt2}$.
\item[\textbf{3.28}] $-\tg x$; $\sec x$.
\item[\textbf{3.29}] $\dfrac1{2\sqrt x\,(1+x)}$.
\item[\textbf{3.30}] $q_h\to\cos x$; usa $\lim\frac{\sen h}h=1$ e $\lim\frac{\cos h-1}h=0$.
\item[\textbf{3.31}] $y=-2\sqrt\pi\,(x-\sqrt\pi)$.
\item[\textbf{3.32}] $y''=(x+2)e^x$; $y^{(n)}=(x+n)e^x$.
\item[\textbf{3.33}] $-\cos x$; $-\cos x$.
\item[\textbf{3.34}] $(1,\tfrac12)$ e $(-1,-\tfrac12)$.
\item[\textbf{3.35}] $2{,}025$; $2{,}025$.
\item[\textbf{3.36}] $h'(1)=-3$; $k'(1)=3\cdot5+2\cdot(-3)=9$.
\item[\textbf{3.37}] $(f^{-1})'(1)=\tfrac12$; tangente $y=\tfrac12(x-1)$.
\item[\textbf{3.38}] $\dfrac{2x}{x^2+1}-\dfrac1{x-1}$.
\item[\textbf{3.39}] $C'(20)=3$; o custo real é $C(21)-C(20)=3{,}01$.
\item[\textbf{3.40}] $3^{\sen x}\ln3\cos x$; $\dfrac{2x}{(x^2+1)\ln2}$.
\item[\textbf{3.41}] (a) $t=2$ e $t=4$; (b) frente em $[0,2)\cup(4,\infty)$, trás em $(2,4)$; (c) 28 m; (d) em $(2,3)$ e $(4,\infty)$.
\item[\textbf{3.42}] Horizontal em $(1,2)$ e $(-1,-2)$; vertical em $(2,1)$ e $(-2,-1)$.
\item[\textbf{3.43}] $y'=x^x(\ln x+1)$; tangente horizontal em $x=1/e$, com $y=e^{-1/e}\approx0{,}692$.
\item[\textbf{3.44}] $y=6x-9$ (tangente em $x=3$) e $y=-2x-1$ (tangente em $x=-1$).
\item[\textbf{3.45}] Use $\cos(x+h)=\cos x\cos h-\sen x\sen h$ e os limites do Teorema 2.7.
\item[\textbf{3.46}] $y=\arctg x\Rightarrow y'=\dfrac1{\sec^2y}=\dfrac1{1+\tg^2y}$.
\item[\textbf{3.47}] $x^{\sen x}\Big(\cos x\ln x+\dfrac{\sen x}x\Big)$.
\item[\textbf{3.48}] $y\Big(\dfrac2{x+1}+\dfrac3{x-2}-\dfrac4{x+3}\Big)$.
\item[\textbf{3.49}] $y'=-\tfrac xy$; derive de novo e use $x^2+y^2=r^2$.
\item[\textbf{3.50}] $y'=-\tfrac85$; $y=-\tfrac85x+\tfrac{18}5$.
\item[\textbf{3.51}] Em todo $x\neq\pm2$; em $x=2$: $-4$ e $4$; em $x=-2$: $-4$ e $4$.
\item[\textbf{3.52}] Derive as definições com exponenciais; use $\cosh^2-\senh^2=1$.
\item[\textbf{3.53}] $f'(0)=0$; $f'(x)=2x\sen\tfrac1x-\cos\tfrac1x$ não tem limite em 0.
\item[\textbf{3.54}] Normal $y=-x+1$; corta em $(1,0)$.
\item[\textbf{3.55}] $a=e$, no ponto $(1,e)$.
\item[\textbf{3.56}] $y'=0$, logo $\arcsen x+\arccos x=\tfrac\pi2$.
\item[\textbf{3.57}] Substitua $y'$ e $y''$; os termos se cancelam.
\item[\textbf{3.58}] $y=e^e+2e^e(x-e)$.
\item[\textbf{3.59}] $0{,}4\pi\approx1{,}26$ rad/s, quando $\theta=0$ ($t=\tfrac14+\tfrac n2$).
\item[\textbf{3.60}] $x=\pm2$: $y=12x-16$ e $y=12x+16$.
\item[\textbf{3.61}] Aplique a regra do produto a $(fg)\cdot h$; por indução, soma de $n$ termos com uma derivada em cada.
\item[\textbf{3.62}] $n=2$ (raiz dupla de $n^2-4n+4=0$).
\item[\textbf{3.63}] Base $n=1$; passo: derive $(x+n)e^x$ e obtenha $(x+n+1)e^x$.
\item[\textbf{3.64}] Inclinações $\tfrac xy$ e $-\tfrac yx$: o produto é $-1$.
\item[\textbf{3.65}] $y=0$ e $y=4x-4$.
\end{itemize}\end{multicols}}


% =====================================================================
\chapter{Aplicações das derivadas}

Com as regras de derivação em mãos, passamos às suas consequências. A derivada responde a
perguntas muito diferentes: com que rapidez uma grandeza varia quando outra varia (taxas
relacionadas), qual é o maior ou o menor valor de uma função (otimização), qual é a forma de
um gráfico, como calcular limites difíceis (regra de L'Hôpital) e como aproximar raízes de
equações (método de Newton). Terminamos com o caminho inverso, as antiderivadas, que abrem
as portas para o capítulo de integração.

% ---------------------------------------------------------------------
\section{Taxas relacionadas}

Quando duas grandezas estão ligadas por uma equação e ambas variam com o tempo, suas taxas
de variação também ficam ligadas. Derivando a equação em relação a $t$ (implicitamente),
obtemos uma relação entre as taxas.

\begin{estrategia}{taxas relacionadas}
\begin{enumerate}[label=\arabic*.]
\item Faça um desenho e dê nome às variáveis, indicando quais dependem do tempo.
\item Anote as taxas conhecidas e a taxa pedida, como derivadas em relação a $t$.
\item Escreva uma equação que relacione as variáveis (geometria, física etc.).
\item Derive a equação em relação a $t$, usando a regra da cadeia.
\item Só então substitua os valores do instante pedido e isole a taxa desejada.
\end{enumerate}
\end{estrategia}

\begin{nota}
Substitua os valores numéricos \emph{depois} de derivar. Se você substituir antes,
grandezas que variam viram constantes e suas derivadas somem indevidamente.
\end{nota}

\begin{exemplo}{4.1}{Enchendo um balão}
Um balão esférico é inflado a $3$ cm$^3$/s. Com que rapidez o raio cresce quando $r=8$ cm?
\solucao
$V=\frac43\pi r^3\Rightarrow\frac{dV}{dt}=4\pi r^2\frac{dr}{dt}$. Com $\frac{dV}{dt}=3$ e
$r=8$: $\frac{dr}{dt}=\frac{3}{256\pi}\approx0{,}0037$ cm/s. Quanto maior o balão, mais devagar
o raio cresce para a mesma vazão.
\end{exemplo}

\begin{exemplo}{4.2}{A escada que escorrega}
Uma escada de 10 m está apoiada numa parede. A base se afasta da parede a 1 m/s. Com que
rapidez o topo desce quando a base está a 6 m da parede?
\solucao
\figtikz{\begin{tikzpicture}[scale=0.42]
\fill[cinza!25] (-0.4,0) rectangle (0,9); \draw[thick] (0,0) -- (0,9); \draw[thick] (0,0) -- (9,0);
\draw[azul,line width=2pt] (0,8) -- (6,0);
\draw[<->,laranja] (0,-0.7) -- node[below,font=\small]{$x$} (6,-0.7);
\draw[<->,laranja] (-1,0) -- node[left,font=\small]{$y$} (-1,8);
\node[azul,font=\small] at (3.8,4.6) {10 m};
\draw[-{Stealth},thick] (6.3,0.4) -- (8,0.4); \draw[-{Stealth},thick] (0.4,8.3) -- (0.4,6.8);
\end{tikzpicture}}{A escada: $x^2+y^2=100$.}
Seja $x$ a distância da base à parede e $y$ a altura do topo: $x^2+y^2=100$. Derivando,
$2x\frac{dx}{dt}+2y\frac{dy}{dt}=0$. Quando $x=6$, $y=8$; com $\frac{dx}{dt}=1$:
$\frac{dy}{dt}=-\frac{6\cdot1}{8}=-\frac34$ m/s. O topo desce a $0{,}75$ m/s.
\end{exemplo}

\begin{exemplo}{4.3}{Enchendo um tanque cônico}
Um tanque em forma de cone invertido tem 4 m de altura e raio do topo de 2 m. Entra água a
$2$ m$^3$/min. Com que rapidez o nível sobe quando a água tem 3 m de profundidade?
\solucao
Por semelhança de triângulos, $r=\frac h2$. Então $V=\frac13\pi r^2h=\frac{\pi h^3}{12}$ e
$\frac{dV}{dt}=\frac{\pi h^2}{4}\frac{dh}{dt}$. Com $h=3$: $2=\frac{9\pi}{4}\frac{dh}{dt}$, logo
$\frac{dh}{dt}=\frac{8}{9\pi}\approx0{,}28$ m/min.
\end{exemplo}

% ---------------------------------------------------------------------
\section{Aproximações lineares e diferenciais}

Perto do ponto de tangência, o gráfico de uma função derivável é quase indistinguível de
sua reta tangente. A \textbf{aproximação linear} (ou linearização) de $f$ em $a$ é
\[
L(x)=f(a)+f'(a)(x-a),\qquad f(x)\approx L(x)\ \text{ para } x \text{ perto de } a .
\]

\begin{exemplo}{4.4}{Linearizações úteis}
(a) Para $f(x)=\sqrt x$ em $a=9$: $L(x)=3+\frac16(x-9)$, logo
$\sqrt{9{,}1}\approx3+\frac{0{,}1}6\approx3{,}0167$ (o valor exato é $3{,}01662\ldots$).

(b) Para $f(x)=\sin x$ em $a=0$: $L(x)=x$. É a aproximação $\sin x\approx x$ para ângulos
pequenos, usada, por exemplo, no estudo do pêndulo.

(c) Para $f(x)=(1+x)^n$ em $a=0$: $(1+x)^n\approx1+nx$. Assim, $\sqrt[3]{1{,}03}\approx1{,}01$.
\end{exemplo}

\textbf{Diferenciais.} Se $y=f(x)$, a \textbf{diferencial} $dy=f'(x)\,dx$ aproxima a variação
real $\Delta y=f(x+dx)-f(x)$. Na prática, as diferenciais estimam como um erro de medida se
propaga.

\begin{exemplo}{4.5}{Propagação de erro}
A aresta de um cubo mede $5$ cm, com erro possível de $\pm0{,}1$ cm. Estime o erro no volume.
\solucao
$V=x^3$, $dV=3x^2\,dx=3(25)(0{,}1)=7{,}5$ cm$^3$. O volume é $125\pm7{,}5$ cm$^3$. O erro
relativo é $\frac{dV}V=\frac{7{,}5}{125}=6\%$, o triplo do erro relativo da aresta ($2\%$),
como sugere $\frac{dV}{V}=3\frac{dx}{x}$.
\end{exemplo}

% ---------------------------------------------------------------------
\section{Máximos e mínimos}

\begin{definicao}
$f$ tem um \textbf{máximo absoluto} em $c$ se $f(c)\ge f(x)$ para todo $x$ do domínio, e um
\textbf{mínimo absoluto} se $f(c)\le f(x)$ para todo $x$. Tem um \textbf{máximo local}
(ou relativo) em $c$ se $f(c)\ge f(x)$ para todo $x$ em algum intervalo aberto contendo $c$;
analogamente para \textbf{mínimo local}. Esses valores são os \textbf{extremos} de $f$.
\end{definicao}

\figtikz{\begin{tikzpicture}
\begin{axis}[painel,width=10cm,height=5cm,xmin=-0.2,xmax=6.3,ymin=-1.5,ymax=3.8,xtick=\empty,ytick=\empty]
\addplot[domain=0.3:6,samples=200]{0.18*(x-0.3)*(x-2.4)*(x-4.6)*(x-5.6)+1.2};
\fill[laranja] (axis cs:0.3,1.2) circle (2pt);
\fill[laranja] (axis cs:6,{0.18*(5.7)*(3.6)*(1.4)*(0.4)+1.2}) circle (2pt) node[left,font=\scriptsize]{máx. absoluto};
\node[font=\scriptsize] at (axis cs:1.35,3.1) {máx. local};
\node[font=\scriptsize] at (axis cs:3.55,-1.15) {mín. absoluto (e local)};
\node[font=\scriptsize] at (axis cs:5.15,1.9) {máx. local};
\end{axis}\end{tikzpicture}}{Extremos locais e absolutos de uma função contínua em um
intervalo fechado. Os extremos absolutos podem ocorrer nas extremidades.}

\begin{teorema}[Teorema do Valor Extremo]{4.1}
Se $f$ é contínua em um intervalo fechado e limitado $[a,b]$, então $f$ tem máximo e mínimo
absolutos em $[a,b]$.
\end{teorema}

As duas hipóteses importam: $f(x)=x$ em $(0,1)$ não tem máximo nem mínimo (o intervalo é
aberto), e $f(x)=1/x$ em $[-1,1]\setminus\{0\}$ também não (falta continuidade).

\begin{teorema}[Teorema de Fermat]{4.2}
Se $f$ tem um extremo local em $c$ e $f$ é derivável em $c$, então $f'(c)=0$.
\end{teorema}

Um \textbf{ponto crítico} de $f$ é um ponto $c$ do domínio em que $f'(c)=0$ ou $f'(c)$ não
existe. Pelo teorema de Fermat, \emph{extremos locais só ocorrem em pontos críticos}. A
recíproca é falsa: $f(x)=x^3$ tem $f'(0)=0$ e nenhum extremo em $0$.

\begin{estrategia}{extremos absolutos em $[a,b]$ (método do intervalo fechado)}
\begin{enumerate}[label=\arabic*.]
\item Calcule $f$ nas extremidades $a$ e $b$.
\item Encontre os pontos críticos em $(a,b)$ e calcule $f$ neles.
\item O maior desses valores é o máximo absoluto; o menor, o mínimo absoluto.
\end{enumerate}
\end{estrategia}

\begin{exemplo}{4.6}{Extremos absolutos}
Encontre os extremos absolutos de $f(x)=x^2-3x^{2/3}$ em $[0,2]$.
\solucao
$f'(x)=2x-2x^{-1/3}=\dfrac{2(x^{4/3}-1)}{x^{1/3}}$. Pontos críticos: $f'(x)=0$ em $x=1$, e
$f'$ não existe em $x=0$ (que é também uma extremidade). Valores: $f(0)=0$, $f(1)=-2$ e
$f(2)=4-3\sqrt[3]4\approx-0{,}76$. Máximo absoluto $0$ em $x=0$; mínimo absoluto $-2$ em $x=1$.
\end{exemplo}

% ---------------------------------------------------------------------
\section{O Teorema do Valor Médio}

\begin{teorema}[Teorema de Rolle]{4.3}
Se $f$ é contínua em $[a,b]$, derivável em $(a,b)$ e $f(a)=f(b)$, então existe $c\in(a,b)$ com
$f'(c)=0$.
\end{teorema}
\emph{Ideia da prova.} Pelo Teorema do Valor Extremo, $f$ tem máximo e mínimo em $[a,b]$. Se
ambos ocorrem nas extremidades, $f$ é constante e $f'\equiv0$. Caso contrário, um deles
ocorre em um ponto interior $c$, e pelo teorema de Fermat $f'(c)=0$. $\square$

\begin{teorema}[Teorema do Valor Médio]{4.4}
Se $f$ é contínua em $[a,b]$ e derivável em $(a,b)$, então existe $c\in(a,b)$ tal que
\[
f'(c)=\frac{f(b)-f(a)}{b-a}.
\]
\end{teorema}

\figtikz{\begin{tikzpicture}
\begin{axis}[painel,width=8cm,height=5.4cm,xmin=0,xmax=5,ymin=0,ymax=4.5,xtick={0.6,4.4},xticklabels={$a$,$b$},ytick=\empty]
\addplot[domain=0.3:4.7]{0.5+0.9*x-0.08*x^2+0.6*sin(deg(1.4*x))};
\pgfmathsetmacro{\fa}{0.5+0.9*0.6-0.08*0.36+0.6*sin(deg(0.84))}
\pgfmathsetmacro{\fb}{0.5+0.9*4.4-0.08*19.36+0.6*sin(deg(6.16))}
\pgfmathsetmacro{\m}{(\fb-\fa)/3.8}
\addplot[laranja,domain=0.6:4.4]{\fa+\m*(x-0.6)};
\fill[laranja] (axis cs:0.6,\fa) circle (2pt); \fill[laranja] (axis cs:4.4,\fb) circle (2pt);
\addplot[verde,dashed,domain=1.6:4.6]{0.5+0.9*3.17-0.08*3.17^2+0.6*sin(deg(1.4*3.17))+\m*(x-3.17)};
\fill[verde] (axis cs:3.17,{0.5+0.9*3.17-0.08*3.17^2+0.6*sin(deg(1.4*3.17))}) circle (2pt) node[below,font=\scriptsize]{$c$};
\end{axis}\end{tikzpicture}}{Teorema do Valor Médio: existe um ponto $c$ em que a tangente
(verde) é paralela à secante entre as extremidades (laranja).}

Em termos físicos: se um carro percorre 200 km em 2 horas, em algum instante sua velocidade
foi exatamente 100 km/h.

\begin{exemplo}{4.7}{Verificando o TVM}
Para $f(x)=\sqrt x$ em $[0,9]$: $f$ é contínua em $[0,9]$ e derivável em $(0,9)$. A
inclinação média é $\frac{3-0}{9}=\frac13$, e $f'(c)=\frac1{2\sqrt c}=\frac13$ dá
$c=\frac94\in(0,9)$.
\end{exemplo}

O TVM é a ponte entre o sinal da derivada e o comportamento da função:

\begin{teorema}[Consequências do TVM]{4.5}
Seja $f$ derivável em um intervalo $I$.
\begin{enumerate}[label=(\roman*)]
\item Se $f'(x)=0$ para todo $x\in I$, então $f$ é constante em $I$.
\item Se $f'(x)=g'(x)$ para todo $x\in I$, então $f(x)=g(x)+C$ para alguma constante $C$.
\item Se $f'>0$ em $I$, $f$ é crescente em $I$; se $f'<0$, $f$ é decrescente.
\end{enumerate}
\end{teorema}
\emph{Prova de (iii).} Se $x_1<x_2$ em $I$, o TVM dá $c$ com
$f(x_2)-f(x_1)=f'(c)(x_2-x_1)>0$. $\square$

% ---------------------------------------------------------------------
\section{Derivadas e a forma do gráfico}

\begin{teorema}[Teste da primeira derivada]{4.6}
Seja $c$ um ponto crítico de $f$, contínua em um intervalo aberto contendo $c$.
\begin{itemize}
\item Se $f'$ muda de positiva para negativa em $c$, $f$ tem um máximo local em $c$.
\item Se $f'$ muda de negativa para positiva em $c$, $f$ tem um mínimo local em $c$.
\item Se $f'$ não muda de sinal em $c$, não há extremo local em $c$.
\end{itemize}
\end{teorema}

\textbf{Concavidade.} Se $f'$ é crescente em $I$, o gráfico de $f$ se curva para cima, e
dizemos que $f$ é \textbf{côncava para cima} em $I$. Se $f'$ é decrescente, $f$ é
\textbf{côncava para baixo}. Um \textbf{ponto de inflexão} é um ponto do gráfico em que a
concavidade muda.

\begin{teorema}[Teste da concavidade]{4.7}
Se $f''>0$ em $I$, $f$ é côncava para cima em $I$; se $f''<0$, côncava para baixo.
\end{teorema}

\begin{teorema}[Teste da segunda derivada]{4.8}
Suponha $f'(c)=0$ e $f''$ contínua perto de $c$. Se $f''(c)>0$, $f$ tem mínimo local em $c$;
se $f''(c)<0$, máximo local. Se $f''(c)=0$, o teste é inconclusivo.
\end{teorema}

\begin{exemplo}{4.8}{Análise completa de um polinômio}
Estude $f(x)=x^3-3x^2-9x+1$.
\solucao
$f'(x)=3x^2-6x-9=3(x-3)(x+1)$: crescente em $(-\infty,-1)$ e $(3,\infty)$, decrescente em
$(-1,3)$. Máximo local $f(-1)=6$; mínimo local $f(3)=-26$.
$f''(x)=6x-6$: côncava para baixo em $(-\infty,1)$ e para cima em $(1,\infty)$; inflexão em
$(1,-10)$. Pelo teste da segunda derivada, $f''(-1)=-12<0$ confirma o máximo e $f''(3)=12>0$,
o mínimo.
\figtikz{\begin{tikzpicture}
\begin{axis}[painel,width=8cm,height=5.4cm,xmin=-3,xmax=5.5,ymin=-32,ymax=14,xtick={-1,1,3},ytick={-26,-10,6}]
\addplot[domain=-3:1,samples=120]{x^3-3*x^2-9*x+1};
\addplot[laranja,domain=1:5.3,samples=120]{x^3-3*x^2-9*x+1};
\fill[black] (axis cs:-1,6) circle (2pt); \fill[black] (axis cs:3,-26) circle (2pt);
\fill[verde] (axis cs:1,-10) circle (2.2pt);
\node[font=\scriptsize,azul] at (axis cs:-2.1,-14) {côncava p/ baixo};
\node[font=\scriptsize,laranja] at (axis cs:4.2,-10) {côncava p/ cima};
\end{axis}\end{tikzpicture}}{$f(x)=x^3-3x^2-9x+1$: extremos locais (pretos) e ponto de
inflexão (verde). A cor indica a concavidade.}
\end{exemplo}

% ---------------------------------------------------------------------
\section{Limites no infinito e assíntotas}

Escrevemos $\lim_{x\to\infty}f(x)=L$ se $f(x)$ fica arbitrariamente perto de $L$ para $x$
suficientemente grande; analogamente para $x\to-\infty$. Nesse caso, a reta $y=L$ é uma
\textbf{assíntota horizontal}. As propriedades dos limites continuam valendo, e o limite
básico é $\lim_{x\to\pm\infty}\frac1{x^n}=0$ para $n>0$.

\textbf{Funções racionais.} Dividimos numerador e denominador pela maior potência de $x$ do
denominador. Por exemplo,
\[
\lim_{x\to\infty}\frac{3x^2+5}{2x^2-x}=\lim_{x\to\infty}\frac{3+5/x^2}{2-1/x}=\frac32 .
\]
Em geral, para $\frac{p(x)}{q(x)}$: se $\operatorname{grau}p<\operatorname{grau}q$, o limite
é $0$; se os graus são iguais, é a razão dos coeficientes dominantes; se
$\operatorname{grau}p>\operatorname{grau}q$, o limite é $\pm\infty$. Quando
$\operatorname{grau}p=\operatorname{grau}q+1$, a divisão de polinômios dá uma
\textbf{assíntota oblíqua}: $\frac{p(x)}{q(x)}=mx+b+\frac{r(x)}{q(x)}$, com
$\frac{r}{q}\to0$.

\begin{nota}
Cuidado com raízes quando $x\to-\infty$: $\sqrt{x^2}=|x|=-x$. Assim,
$\lim_{x\to\infty}\frac{\sqrt{x^2+1}}{x}=1$, mas $\lim_{x\to-\infty}\frac{\sqrt{x^2+1}}{x}=-1$.
\end{nota}

Comportamentos transcendentes a lembrar: $e^x\to\infty$ e $e^{-x}\to0$ quando $x\to\infty$;
$\ln x\to\infty$ (devagar); $\tan^{-1}x\to\pm\frac\pi2$ quando $x\to\pm\infty$.

\begin{estrategia}{esboço do gráfico de uma função}
\begin{enumerate}[label=\arabic*.]
\item Domínio, interceptos e simetria.
\item Assíntotas verticais (limites infinitos) e horizontais/oblíquas (limites no infinito).
\item $f'$: intervalos de crescimento e extremos locais.
\item $f''$: concavidade e pontos de inflexão.
\item Junte tudo num esboço.
\end{enumerate}
\end{estrategia}

\begin{exemplo}{4.9}{Esboço com assíntota oblíqua}
Esboce $f(x)=\dfrac{x^2}{x-1}$.
\solucao
\emph{Domínio:} $x\neq1$; intercepto $(0,0)$. \emph{Assíntotas:} vertical $x=1$
($-\infty$ pela esquerda, $+\infty$ pela direita); como $\frac{x^2}{x-1}=x+1+\frac1{x-1}$, a
reta $y=x+1$ é assíntota oblíqua. \emph{Primeira derivada:}
$f'(x)=\frac{x(x-2)}{(x-1)^2}$: máximo local $f(0)=0$, mínimo local $f(2)=4$.
\emph{Segunda derivada:} $f''(x)=\frac2{(x-1)^3}$: côncava para baixo em $x<1$, para cima em
$x>1$, sem inflexão (que exigiria $x=1$, fora do domínio).
\figtikz{\begin{tikzpicture}
\begin{axis}[painel,width=7.5cm,height=6cm,xmin=-4,xmax=6,ymin=-8,ymax=11,xtick={1,2},ytick={4}]
\addplot[domain=-4:0.85,samples=200]{x^2/(x-1)}; \addplot[domain=1.12:6,samples=200]{x^2/(x-1)};
\addplot[laranja,dashed,domain=-4:6]{x+1}; \draw[laranja,dashed] (axis cs:1,-8) -- (axis cs:1,11);
\fill[black] (axis cs:0,0) circle (2pt); \fill[black] (axis cs:2,4) circle (2pt);
\end{axis}\end{tikzpicture}}{$f(x)=x^2/(x-1)$, com assíntotas $x=1$ e $y=x+1$.}
\end{exemplo}

% ---------------------------------------------------------------------
\section{Otimização}

Muitos problemas práticos pedem o maior ou o menor valor de uma grandeza: área máxima, custo
mínimo, tempo mínimo. A estratégia combina modelagem com o método do intervalo fechado ou com
o teste da primeira derivada.

\begin{estrategia}{problemas de otimização}
\begin{enumerate}[label=\arabic*.]
\item Faça um desenho e nomeie as variáveis.
\item Escreva a grandeza a otimizar em função das variáveis.
\item Use as restrições do problema para reduzir a uma função de \emph{uma} variável, e
  determine o intervalo de valores possíveis.
\item Encontre os pontos críticos e compare os valores (extremidades incluídas), ou use o
  teste da primeira ou da segunda derivada.
\item Responda à pergunta, com unidades.
\end{enumerate}
\end{estrategia}

\begin{exemplo}{4.10}{A caixa de volume máximo}
De uma folha de $24\times36$ cm cortam-se quadrados de lado $x$ nos cantos e dobram-se as
abas para formar uma caixa sem tampa. Qual $x$ maximiza o volume?
\solucao
$V(x)=x(36-2x)(24-2x)=4x^3-120x^2+864x$, com $0\le x\le12$. Então
$V'(x)=12x^2-240x+864=12(x^2-20x+72)=0$ dá $x=10\pm\sqrt{28}$; só $x=10-2\sqrt7\approx4{,}71$
está em $[0,12]$. Como $V(0)=V(12)=0$, o máximo é $V(4{,}71)\approx1825$ cm$^3$.
\end{exemplo}

\begin{exemplo}{4.11}{A lata mais econômica}
Uma lata cilíndrica fechada deve ter volume $V$. Que proporção entre altura e raio minimiza a
área de material?
\solucao
$V=\pi r^2h\Rightarrow h=\frac V{\pi r^2}$. A área é
$A(r)=2\pi r^2+2\pi rh=2\pi r^2+\frac{2V}r$, para $r>0$. $A'(r)=4\pi r-\frac{2V}{r^2}=0$ dá
$r^3=\frac V{2\pi}$. Como $A''(r)=4\pi+\frac{4V}{r^3}>0$, é mínimo. Nesse caso
$h=\frac{V}{\pi r^2}=\frac{2\pi r^3}{\pi r^2}=2r$: a lata ótima tem \emph{altura igual ao
diâmetro}.
\end{exemplo}

% ---------------------------------------------------------------------
\section{A regra de L'Hôpital}

\begin{teorema}[Regra de L'Hôpital]{4.9}
Sejam $f$ e $g$ deriváveis perto de $a$ (exceto possivelmente em $a$), com $g'(x)\neq0$
perto de $a$. Se $\lim_{x\to a}f(x)=\lim_{x\to a}g(x)=0$, ou se ambos são $\pm\infty$, então
\[
\lim_{x\to a}\frac{f(x)}{g(x)}=\lim_{x\to a}\frac{f'(x)}{g'(x)},
\]
desde que o limite da direita exista (ou seja $\pm\infty$). O mesmo vale para $x\to\pm\infty$ e
para limites laterais.
\end{teorema}

\begin{nota}
Deriva-se numerador e denominador \emph{separadamente}; não se usa a regra do quociente. E
a regra só vale para as formas $\frac00$ e $\frac\infty\infty$: por exemplo,
$\lim_{x\to0}\frac{x+1}{x+2}=\frac12$, enquanto ``aplicar L'Hôpital'' daria, erradamente, 1.
\end{nota}

\begin{exemplo}{4.12}{As formas $\frac00$ e $\frac\infty\infty$}
(a) $\displaystyle\lim_{x\to0}\frac{1-\cos x}{x^2}=\lim_{x\to0}\frac{\sin x}{2x}=\frac12$.\qquad
(b) $\displaystyle\lim_{x\to\infty}\frac{e^x}{x^2}=\lim_{x\to\infty}\frac{e^x}{2x}
=\lim_{x\to\infty}\frac{e^x}{2}=\infty$.
\end{exemplo}

\textbf{Outras formas indeterminadas.} As formas $0\cdot\infty$ e $\infty-\infty$ são
reescritas como quocientes; as formas $1^\infty$, $0^0$ e $\infty^0$ são tratadas com
logaritmo: se $y=f(x)^{g(x)}$, calculamos $\lim\ln y=\lim g(x)\ln f(x)$ e depois
exponenciamos.

\begin{exemplo}{4.13}{Formas indeterminadas diversas}
(a) $0\cdot\infty$: $\displaystyle\lim_{x\to0^+}x\ln x=\lim_{x\to0^+}\frac{\ln x}{1/x}
=\lim_{x\to0^+}\frac{1/x}{-1/x^2}=\lim_{x\to0^+}(-x)=0$.

(b) $\infty-\infty$: $\displaystyle\lim_{x\to0^+}\Big(\frac1x-\frac1{\sin x}\Big)
=\lim_{x\to0^+}\frac{\sin x-x}{x\sin x}=\lim_{x\to0^+}\frac{\cos x-1}{\sin x+x\cos x}
=\lim_{x\to0^+}\frac{-\sin x}{2\cos x-x\sin x}=0$.

(c) $1^\infty$: para $y=\big(1+\frac1x\big)^x$, $\ln y=x\ln\big(1+\frac1x\big)
=\frac{\ln(1+1/x)}{1/x}\to1$ (por L'Hôpital). Logo $\lim_{x\to\infty}\big(1+\frac1x\big)^x=e$.

(d) $0^0$: para $y=x^x$, $\ln y=x\ln x\to0$ por (a). Logo $\lim_{x\to0^+}x^x=e^0=1$.
\end{exemplo}

\textbf{Taxas de crescimento.} Aplicando L'Hôpital repetidamente: para todo $p>0$,
$\frac{\ln x}{x^p}\to0$ e $\frac{x^p}{e^x}\to0$ quando $x\to\infty$. Ou seja, no infinito,
\emph{o logaritmo cresce mais devagar que qualquer potência, e toda potência cresce mais
devagar que a exponencial}.

% ---------------------------------------------------------------------
\section{O método de Newton}

Para aproximar uma raiz de $f(x)=0$, partimos de um palpite $x_0$ e trocamos a curva pela
tangente em $(x_0,f(x_0))$. O ponto em que a tangente corta o eixo $x$ é o novo palpite:
\[
x_{n+1}=x_n-\frac{f(x_n)}{f'(x_n)} .
\]

\figtikz{\begin{tikzpicture}
\begin{axis}[painel,width=8cm,height=5.2cm,xmin=0.8,xmax=3.2,ymin=-1.5,ymax=7.5,xtick={3,1.8333,1.4621},xticklabels={$x_0$,$x_1$,$x_2$},ytick=\empty,
  x tick label style={font=\scriptsize}]
\addplot[domain=0.8:3.2]{x^2-2};
\addplot[laranja,domain=1.7:3.2]{7+6*(x-3)};
\addplot[laranja!70,domain=1.35:2.3]{1.3611+3.6667*(x-1.8333)};
\draw[cinza,dotted] (axis cs:3,0) -- (axis cs:3,7); \draw[cinza,dotted] (axis cs:1.8333,0) -- (axis cs:1.8333,1.3611);
\fill[verde] (axis cs:1.4142,0) circle (2pt) node[below left,font=\scriptsize]{$\sqrt2$};
\end{axis}\end{tikzpicture}}{Método de Newton para $f(x)=x^2-2$: cada tangente leva a um
palpite mais próximo da raiz $\sqrt2$.}

\begin{exemplo}{4.14}{Calculando $\sqrt2$}
Para $f(x)=x^2-2$, a iteração é $x_{n+1}=x_n-\frac{x_n^2-2}{2x_n}=\frac12\Big(x_n+\frac2{x_n}\Big)$.
Com $x_0=1$: $x_1=1{,}5$; $x_2\approx1{,}416667$; $x_3\approx1{,}414216$; $x_4\approx1{,}4142136$.
O número de casas corretas praticamente dobra a cada passo.
\end{exemplo}

\textbf{Quando falha.} Se $f'(x_n)=0$, a tangente é horizontal e o método para. Ele também
pode oscilar: para $f(x)=x^3-2x+2$ com $x_0=0$, obtemos $x_1=1$, $x_2=0$, $x_3=1,\dots$, sem
convergir. Por isso convém escolher $x_0$ perto da raiz, por exemplo com o auxílio de um
gráfico ou do Teorema do Valor Intermediário.

% ---------------------------------------------------------------------
\section{Antiderivadas}

\begin{definicao}
$F$ é uma \textbf{antiderivada} (ou primitiva) de $f$ em um intervalo $I$ se $F'(x)=f(x)$ para
todo $x\in I$.
\end{definicao}

Pela consequência (ii) do TVM, se $F$ é uma antiderivada de $f$, então \emph{todas} as
antiderivadas são da forma $F(x)+C$. Escrevemos a \textbf{integral indefinida}
\[
\int f(x)\,dx=F(x)+C .
\]

\figtikz{\begin{tikzpicture}
\begin{axis}[painel,width=6.5cm,height=5cm,xmin=-2.3,xmax=2.3,ymin=-2.5,ymax=5,xtick=\empty,ytick=\empty]
\foreach \c/\o in {-2/40,-1/55,0/100,1/55,2/40}{\edef\tmp{\noexpand\addplot[azul!\o!white,domain=-2.3:2.3]{x^2+\c};}\tmp}
\end{axis}\end{tikzpicture}}{A família $x^2+C$: todas as antiderivadas de $2x$ diferem por
uma translação vertical.}

Cada regra de derivação, lida ao contrário, dá uma regra de antiderivação:
\[
\int x^n\,dx=\frac{x^{n+1}}{n+1}+C\ (n\neq-1),\qquad
\int\frac1x\,dx=\ln|x|+C,\qquad \int e^x\,dx=e^x+C,
\]
\[
\int\cos x\,dx=\sin x+C,\quad \int\sin x\,dx=-\cos x+C,\quad \int\sec^2x\,dx=\tan x+C,\quad
\int\frac{dx}{1+x^2}=\tan^{-1}x+C .
\]
Além disso, $\int\big(kf\pm g\big)\,dx=k\int f\,dx\pm\int g\,dx$. Uma antiderivada sempre pode
ser conferida derivando.

\begin{exemplo}{4.15}{Calculando integrais indefinidas}
$\displaystyle\int\Big(5x^3-\frac{7}{x^2}+4\sqrt x\Big)dx
=\frac{5x^4}4+\frac7x+\frac{8}{3}x^{3/2}+C$
\quad e\quad
$\displaystyle\int\frac{x^2+1}{x}\,dx=\int\Big(x+\frac1x\Big)dx=\frac{x^2}2+\ln|x|+C$.
\end{exemplo}

\textbf{Problemas de valor inicial.} Uma condição como $F(a)=b$ fixa a constante $C$.

\begin{exemplo}{4.16}{Um carro freando}
Um carro a 30 m/s freia com desaceleração constante de 6 m/s$^2$. Que distância percorre até
parar?
\solucao
$a(t)=-6\Rightarrow v(t)=-6t+C_1$, e $v(0)=30$ dá $v(t)=30-6t$. Então
$s(t)=30t-3t^2+C_2$, com $s(0)=0$. O carro para quando $v=0$, em $t=5$ s, tendo percorrido
$s(5)=150-75=75$ m.
\end{exemplo}

% ---------------------------------------------------------------------
\section*{Resumo do capítulo}
\addcontentsline{toc}{section}{Resumo do capítulo}
\begin{itemize}
\item Taxas relacionadas: derive a equação que liga as variáveis em relação a $t$; substitua
  os valores só depois.
\item Linearização $L(x)=f(a)+f'(a)(x-a)$; diferencial $dy=f'(x)\,dx$ estima variações e
  erros.
\item Extremos locais ocorrem em pontos críticos ($f'=0$ ou $\nexists$). Em $[a,b]$ com $f$
  contínua, compare $f$ nos pontos críticos e nas extremidades.
\item TVM: $f'(c)=\frac{f(b)-f(a)}{b-a}$. Consequências: $f'=0\Rightarrow f$ constante; sinal
  de $f'$ dá crescimento.
\item Testes da 1ª e da 2ª derivada; $f''$ dá a concavidade; inflexão onde ela muda.
\item Assíntotas horizontais e oblíquas via limites no infinito; roteiro de esboço.
\item Otimização: modele, reduza a uma variável, encontre e compare pontos críticos.
\item L'Hôpital para $\frac00$ e $\frac\infty\infty$; as demais formas se reduzem a essas.
  $\ln x\ll x^p\ll e^x$.
\item Newton: $x_{n+1}=x_n-f(x_n)/f'(x_n)$.
\item Antiderivadas: $\int f\,dx=F+C$; problemas de valor inicial fixam $C$.
\end{itemize}

% Gerado por ferramentas/gerar.py; não edite à mão.
\section*{Exercícios}
\addcontentsline{toc}{section}{Exercícios}
Os pontos antes de cada número indicam a dificuldade: \nivel{1} fácil, \nivel{2} médio, \nivel{3} difícil e \nivel{4} desafio. Os exercícios marcados com \estrela{} têm solução completa, com animação do gráfico, no repositório do livro: \url{https://github.com/AleffAzevedo/calculo-1-pt/blob/main/exercicios/cap4/README.md}. As respostas estão no fim do capítulo.

\begin{itemize}[leftmargin=5.2em,labelsep=0.6em,itemsep=4pt]
\item[\exnum{1}{4.1}] Encontre a linearização de $f(x)=\sqrt[3]x$ em $a=8$ e use-a para estimar $\sqrt[3]{8{,}2}$.
\item[\exnum{1}{4.2}] Para $y=x^2-3x$, compare $dy$ e $\Delta y$ quando $x=2$ e $dx=0{,}1$.
\item[\exnum{1}{4.3}] Encontre os pontos críticos de $f(x)=x^3-12x+1$.
\item[\exnum{1}{4.4}] Encontre os pontos críticos de $f(x)=x^{2/3}(x-5)$ (inclua os pontos em que $f'$ não existe).
\item[\exnum{1}{4.5}] Encontre os extremos absolutos de $f(x)=x^2-4x+3$ em $[0,3]$.
\item[\exnum{1}{4.6}] Encontre os extremos absolutos de $f(x)=\sin x+\cos x$ em $[0,\pi]$.
\item[\exnum{1}{4.7}] Verifique as hipóteses do Teorema de Rolle para $f(x)=x^2-4x$ em $[0,4]$ e encontre $c$.
\item[\exnum{1}{4.8}] Encontre o $c$ do TVM para $f(x)=x^2$ em $[1,3]$ e interprete.
\item[\exnum{1}{4.9}] Determine os intervalos de crescimento e de decrescimento de $f(x)=x^3-3x^2$.
\item[\exnum{1}{4.10}] Estude a concavidade de $f(x)=x^4-4x^3$ e encontre os pontos de inflexão.
\item[\exnum{1}{4.11}] Use o teste da segunda derivada para classificar os pontos críticos de $f(x)=2x^3-3x^2-12x$.
\item[\exnum{1}{4.12}] Calcule $\displaystyle\lim_{x\to\infty}\frac{3x^2-x}{5x^2+2}$.
\item[\exnum{1}{4.13}] Calcule $\displaystyle\lim_{x\to-\infty}\frac{2x+1}{\sqrt{x^2+3}}$.
\item[\exnum{1}{4.14}] Encontre as assíntotas verticais e horizontais de $f(x)=\dfrac{2x^2}{x^2-9}$.
\item[\exnum{1}{4.15}] Encontre a assíntota oblíqua de $f(x)=\dfrac{x^2+1}{x+1}$.
\item[\exnum{1}{4.16}] Use L'Hôpital para calcular $\displaystyle\lim_{x\to1}\frac{\ln x}{x-1}$ e $\displaystyle\lim_{x\to0}\frac{e^x-1}{x}$.
\item[\exnum{1}{4.17}] Calcule $\displaystyle\lim_{x\to\infty}x^2e^{-x}$.
\item[\exnum{1}{4.18}] Faça duas iterações do método de Newton para $f(x)=x^2-5$ com $x_0=2$.
\item[\exnum{1}{4.19}] Calcule $\int(3x^2-4x+5)\,dx$ e $\int(\sin x+e^x)\,dx$.
\item[\exnum{1}{4.20}] Resolva o problema de valor inicial $f'(x)=6x^2-2$, $f(1)=3$.
\item[\exnum{2}{4.21}] \estrela\ \textbf{A escada escorregando.} Uma escada de 5 m está apoiada numa parede vertical. A base se afasta da parede a
$0{,}5$ m/s. (a) Com que rapidez o topo desce quando a base está a 3 m da parede?
(b) Em que posição as duas extremidades se movem com a mesma rapidez?
\item[\exnum{2}{4.22}] \estrela\ \textbf{A sombra no poste.} Uma pessoa de $1{,}8$ m se afasta de um poste de luz de 6 m de altura a $1{,}5$ m/s.
(a) Com que rapidez sua sombra aumenta? (b) Com que rapidez a ponta da sombra se move?
\item[\exnum{2}{4.23}] \estrela\ \textbf{Avião e radar.} Um avião voa horizontalmente a 4 km de altitude, a 600 km/h, numa trajetória que passa
diretamente acima de uma estação de radar. Com que rapidez a distância avião–radar aumenta
quando essa distância é 5 km (e o avião já passou da estação)?
\item[\exnum{2}{4.24}] Uma mancha de óleo circular se espalha com o raio crescendo a 2 m/s. Com que rapidez a área cresce quando $r=10$ m?
\item[\exnum{2}{4.25}] A aresta de um cubo cresce a $0{,}5$ cm/s. Com que rapidez o volume cresce quando a aresta mede 4 cm?
\item[\exnum{2}{4.26}] O raio de uma esfera mede $10$ cm com erro de $\pm0{,}2$ cm. Estime o erro absoluto e o percentual no volume.
\item[\exnum{2}{4.27}] Use uma aproximação linear para estimar $\sin31^\circ$.
\item[\exnum{2}{4.28}] \estrela\ \textbf{Método do intervalo fechado.} Encontre os valores máximo e mínimo absolutos de $f(x)=x\,e^{-x}$ em $[0,4]$.
\item[\exnum{2}{4.29}] \estrela\ \textbf{Teorema do Valor Médio.} Verifique as hipóteses do TVM para $f(x)=x^3-x$ em $[0,2]$ e encontre todos os $c$ garantidos
pelo teorema.
\item[\exnum{2}{4.30}] Mostre que $x^3+x-1=0$ tem exatamente uma raiz real.
\item[\exnum{2}{4.31}] Mostre que $\tan^{-1}x+\tan^{-1}\tfrac1x=\tfrac\pi2$ para $x>0$.
\item[\exnum{2}{4.32}] Encontre os extremos locais de $f(x)=x^2e^{-x}$.
\item[\exnum{2}{4.33}] Encontre os extremos locais de $f(x)=x-2\sin x$ em $(0,2\pi)$.
\item[\exnum{2}{4.34}] Calcule $\displaystyle\lim_{x\to\infty}\big(\sqrt{x^2+x}-x\big)$.
\item[\exnum{2}{4.35}] Calcule $\displaystyle\lim_{x\to0}\frac{x-\sin x}{x^3}$.
\item[\exnum{2}{4.36}] Calcule $\displaystyle\lim_{x\to\infty}x^{1/x}$.
\item[\exnum{2}{4.37}] Calcule $\displaystyle\lim_{x\to\pi/2^-}(\sec x-\tan x)$.
\item[\exnum{2}{4.38}] \estrela\ \textbf{O limite $(1+a/x)^x$.} Mostre, com L'Hôpital, que $\displaystyle\lim_{x\to\infty}\Big(1+\frac ax\Big)^x=e^a$ para
qualquer constante $a$. Use o resultado para calcular o montante de R\$\,1000 a 6\% ao ano
com capitalização contínua por 10 anos.
\item[\exnum{2}{4.39}] Calcule $\displaystyle\int\frac{x^3-2x}{x^2}\,dx$.
\item[\exnum{2}{4.40}] Calcule $\displaystyle\int\Big(\frac2{1+x^2}-\frac3{\sqrt{1-x^2}}\Big)dx$.
\item[\exnum{3}{4.41}] \estrela\ \textbf{Esboço completo de um gráfico.} Faça o estudo completo e esboce $f(x)=\dfrac{x}{x^2+1}$: domínio, simetria, assíntotas,
crescimento, extremos, concavidade e inflexões.
\item[\exnum{3}{4.42}] \estrela\ \textbf{Cercado junto ao rio.} Um fazendeiro tem 600 m de cerca para construir um cercado retangular junto a um rio reto
(sem cerca do lado do rio), dividido em dois currais por uma cerca perpendicular ao rio. Que
dimensões maximizam a área total?
\item[\exnum{3}{4.43}] \estrela\ \textbf{Ponto mais próximo.} Encontre os pontos da parábola $y=x^2$ mais próximos do ponto $(0,3)$ e a distância mínima.
\item[\exnum{3}{4.44}] \estrela\ \textbf{O cabo mais barato.} Uma usina fica numa margem de um rio de 1 km de largura. Uma fábrica fica na outra
margem, 3 km rio abaixo. O cabo custa R\$\,5 mil/km sob a água e R\$\,4 mil/km em terra.
Qual é o trajeto mais barato (atravessar em diagonal até um ponto e seguir pela margem)?
\item[\exnum{3}{4.45}] \estrela\ \textbf{Método de Newton para $\cos x=x$.} Use o método de Newton com $x_0=1$ para aproximar a solução de $\cos x=x$ com 9 casas
decimais. Quantas iterações são necessárias?
\item[\exnum{3}{4.46}] \estrela\ \textbf{Lançamento vertical.} Uma bola é lançada para cima a partir de 2 m de altura com velocidade de 15 m/s
($g=9{,}8$ m/s²). Partindo de $a(t)=-9{,}8$, encontre $v(t)$ e $s(t)$, a altura máxima e o
instante em que a bola atinge o chão.
\item[\exnum{3}{4.47}] Uma calha de 10 m tem seção em triângulo isósceles invertido, com 1 m de largura no topo e $0{,}5$ m de altura. Entra água a $0{,}4$ m³/min. Com que rapidez o nível sobe quando $h=0{,}2$ m?
\item[\exnum{3}{4.48}] Dois carros partem do mesmo ponto: um para o norte a 60 km/h e outro para o leste a 80 km/h. Com que rapidez a distância entre eles aumenta?
\item[\exnum{3}{4.49}] Um balão sobe verticalmente a 2 m/s e é observado de um ponto no solo a 50 m da base. Com que rapidez o ângulo de elevação aumenta quando o balão está a 50 m de altura?
\item[\exnum{3}{4.50}] Encontre o retângulo de maior área inscrito num semicírculo de raio $R$ (com um lado sobre o diâmetro).
\item[\exnum{3}{4.51}] Uma caixa sem tampa de base quadrada deve ter volume $32\,000$ cm³. Que dimensões minimizam o material?
\item[\exnum{3}{4.52}] Encontre o retângulo de maior área com base no eixo $x$ e vértices superiores na parábola $y=12-x^2$.
\item[\exnum{3}{4.53}] A demanda de um produto é $q=1000-20p$ ($p$ em reais). Que preço maximiza a receita $R=pq$?
\item[\exnum{3}{4.54}] Faça o estudo completo de $f(x)=x^{2/3}(x-5)$ (há uma cúspide) e esboce o gráfico.
\item[\exnum{3}{4.55}] Faça o estudo completo de $f(x)=e^{-x^2}$ e esboce o gráfico.
\item[\exnum{3}{4.56}] Calcule $\displaystyle\lim_{x\to0}\frac{e^x-1-x}{x^2}$ e $\displaystyle\lim_{x\to\infty}x\Big(\frac\pi2-\tan^{-1}x\Big)$.
\item[\exnum{3}{4.57}] Use o TVM para provar que $|\sin a-\sin b|\le|a-b|$ para todos $a,b$.
\item[\exnum{3}{4.58}] Prove que $e^x\ge1+x$ para todo $x$.
\item[\exnum{3}{4.59}] Resolva $f''(x)=6x+2$, $f'(0)=1$, $f(0)=-1$.
\item[\exnum{3}{4.60}] Prove que a equação $\cos x=x$ tem exatamente uma solução.
\item[\exnum{4}{4.61}] Prove que a única função derivável com $f'=f$ e $f(0)=1$ é $f(x)=e^x$. (Dica: derive $g(x)=f(x)e^{-x}$.)
\item[\exnum{4}{4.62}] A luz vai de $A$ (no ar, velocidade $v_1$) a $B$ (na água, velocidade $v_2$) pelo caminho de tempo mínimo. Deduza a lei de Snell $\dfrac{\sin\theta_1}{v_1}=\dfrac{\sin\theta_2}{v_2}$.
\item[\exnum{4}{4.63}] Para $a>0$, mostre que o método de Newton para $x^2=a$ dá $x_{n+1}=\tfrac12\big(x_n+\tfrac a{x_n}\big)$ e que $x_n\ge\sqrt a$ para todo $n\ge1$, qualquer que seja $x_0>0$.
\item[\exnum{4}{4.64}] Calcule $\displaystyle\lim_{x\to0}\Big(\frac1{x^2}-\frac1{\sin^2x}\Big)$.
\item[\exnum{4}{4.65}] Qual é o comprimento da maior barra que pode ser carregada horizontalmente pela esquina entre dois corredores de larguras 8 m e 27 m?
\end{itemize}
\subsection*{Respostas}
{\small\begin{multicols}{2}\raggedright
\begin{itemize}[leftmargin=2.7em,labelsep=0.5em,itemsep=1pt]
\item[\textbf{4.1}] $L(x)=2+\tfrac{x-8}{12}$; $\approx2{,}0167$.
\item[\textbf{4.2}] $dy=0{,}1$; $\Delta y=0{,}11$.
\item[\textbf{4.3}] $x=\pm2$.
\item[\textbf{4.4}] $x=2$ e $x=0$.
\item[\textbf{4.5}] Máximo $3$ em $x=0$; mínimo $-1$ em $x=2$.
\item[\textbf{4.6}] Máximo $\sqrt2$ em $x=\tfrac\pi4$; mínimo $-1$ em $x=\pi$.
\item[\textbf{4.7}] $c=2$.
\item[\textbf{4.8}] $c=2$: tangente paralela à secante.
\item[\textbf{4.9}] Cresce em $(-\infty,0)\cup(2,\infty)$; decresce em $(0,2)$.
\item[\textbf{4.10}] Para cima em $(-\infty,0)\cup(2,\infty)$, para baixo em $(0,2)$; inflexões $(0,0)$ e $(2,-16)$.
\item[\textbf{4.11}] Máx. local $f(-1)=7$; mín. local $f(2)=-20$.
\item[\textbf{4.12}] $\tfrac35$.
\item[\textbf{4.13}] $-2$.
\item[\textbf{4.14}] Verticais $x=\pm3$; horizontal $y=2$.
\item[\textbf{4.15}] $y=x-1$.
\item[\textbf{4.16}] $1$ e $1$.
\item[\textbf{4.17}] $0$.
\item[\textbf{4.18}] $x_1=2{,}25$; $x_2\approx2{,}236111$.
\item[\textbf{4.19}] $x^3-2x^2+5x+C$; $-\cos x+e^x+C$.
\item[\textbf{4.20}] $f(x)=2x^3-2x+3$.
\item[\textbf{4.21}] (a) $0{,}375$ m/s; (b) quando $x=y=\tfrac{5}{\sqrt2}\approx3{,}54$ m.
\item[\textbf{4.22}] (a) $\tfrac{9}{14}\approx0{,}64$ m/s; (b) $\tfrac{15}{7}\approx2{,}14$ m/s.
\item[\textbf{4.23}] 360 km/h.
\item[\textbf{4.24}] $40\pi\approx125{,}7$ m²/s.
\item[\textbf{4.25}] $24$ cm³/s.
\item[\textbf{4.26}] $dV=80\pi\approx251$ cm³; $6\%$.
\item[\textbf{4.27}] $\tfrac12+\tfrac{\sqrt3}2\cdot\tfrac\pi{180}\approx0{,}5151$.
\item[\textbf{4.28}] Máximo $\tfrac1e\approx0{,}368$ em $x=1$; mínimo $0$ em $x=0$.
\item[\textbf{4.29}] $c=\tfrac{2}{\sqrt3}\approx1{,}155$.
\item[\textbf{4.30}] Existe pelo TVI; é única porque $f'=3x^2+1>0$ (pelo Rolle, duas raízes exigiriam $f'=0$).
\item[\textbf{4.31}] A derivada da soma é $0$, então ela é constante; em $x=1$ vale $\tfrac\pi2$.
\item[\textbf{4.32}] Mín. local $f(0)=0$; máx. local $f(2)=\tfrac4{e^2}\approx0{,}541$.
\item[\textbf{4.33}] Mín. $f(\tfrac\pi3)=\tfrac\pi3-\sqrt3$; máx. $f(\tfrac{5\pi}3)=\tfrac{5\pi}3+\sqrt3$.
\item[\textbf{4.34}] $\tfrac12$.
\item[\textbf{4.35}] $\tfrac16$.
\item[\textbf{4.36}] $1$.
\item[\textbf{4.37}] $0$.
\item[\textbf{4.38}] $e^a$; $1000e^{0{,}6}\approx$ R\$\,1822{,}12.
\item[\textbf{4.39}] $\tfrac{x^2}2-2\ln|x|+C$.
\item[\textbf{4.40}] $2\tan^{-1}x-3\sin^{-1}x+C$.
\item[\textbf{4.41}] Ímpar; assíntota $y=0$; máx. $(1,\tfrac12)$, mín. $(-1,-\tfrac12)$; inflexões em $x=0$ e $x=\pm\sqrt3$.
\item[\textbf{4.42}] 100 m (perpendicular ao rio) por 300 m (paralelo); área 30\,000 m².
\item[\textbf{4.43}] $\big(\pm\sqrt{5/2},\,\tfrac52\big)$; distância $\tfrac{\sqrt{11}}{2}\approx1{,}658$.
\item[\textbf{4.44}] Atravessar até o ponto a $\tfrac43$ km rio abaixo; custo mínimo R\$\,15 mil.
\item[\textbf{4.45}] $x\approx0{,}739085133$; 3 iterações (a 4ª só confirma).
\item[\textbf{4.46}] $s(t)=-4{,}9t^2+15t+2$; altura máx. $\approx13{,}48$ m em $t\approx1{,}53$ s; chão em $t\approx3{,}19$ s.
\item[\textbf{4.47}] $0{,}1$ m/min.
\item[\textbf{4.48}] $100$ km/h (constante).
\item[\textbf{4.49}] $0{,}02$ rad/s.
\item[\textbf{4.50}] Lados $\sqrt2R$ e $\tfrac R{\sqrt2}$; área $R^2$.
\item[\textbf{4.51}] Base $40\times40$ cm, altura $20$ cm; área $4800$ cm².
\item[\textbf{4.52}] Base 4, altura 8; área 32.
\item[\textbf{4.53}] $p=25$; receita R\$\,12\,500.
\item[\textbf{4.54}] Máx. local $(0,0)$ (cúspide); mín. local $(2,-3\sqrt[3]4)\approx(2;-4{,}76)$; inflexão $(-1,-6)$.
\item[\textbf{4.55}] Par; máx. $(0,1)$; inflexões em $x=\pm\tfrac1{\sqrt2}$; assíntota $y=0$.
\item[\textbf{4.56}] $\tfrac12$ e $1$.
\item[\textbf{4.57}] $\sin a-\sin b=\cos c\,(a-b)$ e $|\cos c|\le1$.
\item[\textbf{4.58}] $g(x)=e^x-1-x$ tem mínimo absoluto $g(0)=0$.
\item[\textbf{4.59}] $f(x)=x^3+x^2+x-1$.
\item[\textbf{4.60}] Existência pelo TVI; unicidade porque $g(x)=\cos x-x$ tem $g'=-\sin x-1\le0$, zerando só em pontos isolados.
\item[\textbf{4.61}] $g'=0$, logo $g$ é constante $=g(0)=1$.
\item[\textbf{4.62}] Escreva o tempo em função do ponto de cruzamento e anule a derivada.
\item[\textbf{4.63}] Desigualdade entre as médias aritmética e geométrica.
\item[\textbf{4.64}] $-\tfrac13$.
\item[\textbf{4.65}] $13\sqrt{13}\approx46{,}9$ m (minimizando $\tfrac8{\sin\theta}+\tfrac{27}{\cos\theta}$).
\end{itemize}\end{multicols}}


% =====================================================================
\chapter{Integração}

No capítulo~2 vimos que o problema da área sob uma curva leva a um limite de somas de
retângulos. Neste capítulo esse limite ganha nome, \emph{integral definida}, e descobrimos o
fato central do Cálculo: integrar e derivar são operações inversas (Teorema Fundamental do
Cálculo). Com ele, calcular áreas, deslocamentos e variações totais se reduz a encontrar
antiderivadas, e a técnica de substituição amplia muito a lista de funções que sabemos
integrar.

% ---------------------------------------------------------------------
\section{Aproximando áreas}

\subsection{Notação de somatório}\index{somatorio@somatório}

Escrevemos $\displaystyle\sum_{i=1}^n a_i=a_1+a_2+\dots+a_n$. Valem
$\sum(a_i\pm b_i)=\sum a_i\pm\sum b_i$ e $\sum c\,a_i=c\sum a_i$, além das fórmulas
\[
\sum_{i=1}^n1=n,\qquad \sum_{i=1}^n i=\frac{n(n+1)}2,\qquad
\sum_{i=1}^n i^2=\frac{n(n+1)(2n+1)}6,\qquad \sum_{i=1}^n i^3=\frac{n^2(n+1)^2}4 .
\]

\subsection{Somas de Riemann}

Para aproximar a área sob o gráfico de $f\ge0$ em $[a,b]$, dividimos o intervalo em $n$
subintervalos de largura $\Delta x=\frac{b-a}n$, com extremos $x_i=a+i\,\Delta x$, e somamos
áreas de retângulos. A altura de cada retângulo é o valor de $f$ em algum ponto $x_i^*$ do
subintervalo:
\[
\sum_{i=1}^n f(x_i^*)\,\Delta x \qquad(\textbf{soma de Riemann}\index{soma de riemann@soma de Riemann}).
\]
Usando a extremidade esquerda ($x_{i-1}$), a direita ($x_i$) ou o ponto médio, obtemos as
aproximações $L_n$, $R_n$ e $M_n$.

\figtikz{\begin{tikzpicture}
\begin{groupplot}[group style={group size=3 by 1,horizontal sep=1.0cm},painel,width=5.3cm,height=4.6cm,
  xmin=-0.1,xmax=2.2,ymin=0,ymax=4.6,xtick={0,0.5,1,1.5,2},ytick={1,4},xticklabel style={font=\tiny}]
\nextgroupplot[title={$L_4=1{,}75$}]
\foreach \a in {0,0.5,1,1.5}{\edef\tmp{\noexpand\draw[fill=laranja!25,draw=laranja] (axis cs:\a,0) rectangle (axis cs:{\a+0.5},{\a*\a});}\tmp}
\addplot[domain=0:2.1]{x^2};
\nextgroupplot[title={$R_4=3{,}75$}]
\foreach \a in {0,0.5,1,1.5}{\edef\tmp{\noexpand\draw[fill=laranja!25,draw=laranja] (axis cs:\a,0) rectangle (axis cs:{\a+0.5},{(\a+0.5)^2});}\tmp}
\addplot[domain=0:2.1]{x^2};
\nextgroupplot[title={$M_4=2{,}625$}]
\foreach \a in {0,0.5,1,1.5}{\edef\tmp{\noexpand\draw[fill=laranja!25,draw=laranja] (axis cs:\a,0) rectangle (axis cs:{\a+0.5},{(\a+0.25)^2});}\tmp}
\addplot[domain=0:2.1]{x^2};
\end{groupplot}\end{tikzpicture}}{Somas de Riemann com $n=4$ para $f(x)=x^2$ em $[0,2]$:
extremidade esquerda, direita e ponto médio. A área exata é $\tfrac83\approx2{,}667$.}

\begin{exemplo}{5.1}{Aproximações e o limite}
Para $f(x)=x^2$ em $[0,2]$ com $n=4$ ($\Delta x=0{,}5$):
$L_4=0{,}5(0+0{,}25+1+2{,}25)=1{,}75$ e $R_4=0{,}5(0{,}25+1+2{,}25+4)=3{,}75$. Como $f$ é
crescente, $L_n$ subestima e $R_n$ superestima a área.

Para $n$ qualquer, $x_i=\frac{2i}n$ e
\[
R_n=\sum_{i=1}^n\Big(\frac{2i}n\Big)^2\frac2n=\frac8{n^3}\sum_{i=1}^ni^2
=\frac8{n^3}\cdot\frac{n(n+1)(2n+1)}6\ \xrightarrow{n\to\infty}\ \frac{16}{6}=\frac83 .
\]
\end{exemplo}

% ---------------------------------------------------------------------
\section{A integral definida}

\begin{definicao}
Se $f$ está definida em $[a,b]$ e o limite
\[
\int_a^b f(x)\,dx=\lim_{n\to\infty}\sum_{i=1}^n f(x_i^*)\,\Delta x
\]
existe e é o mesmo para quaisquer escolhas dos pontos $x_i^*$, dizemos que $f$ é
\textbf{integrável}\index{funcao@função!integravel@integrável} em $[a,b]$ e o limite é a \textbf{integral definida}\index{integral!definida} de $f$ de $a$ a $b$.
Os números $a$ e $b$ são os \emph{limites de integração}, e $f$ é o \emph{integrando}.
\end{definicao}

Toda função contínua em $[a,b]$ é integrável (e também as que têm um número finito de
descontinuidades de salto). A variável de integração é ``muda'':
$\int_a^bf(x)\,dx=\int_a^bf(t)\,dt$.

\textbf{Área com sinal.}\index{area com sinal@área com sinal} Se $f$ muda de sinal, os retângulos abaixo do eixo contribuem
negativamente: $\int_a^bf(x)\,dx$ é a \emph{área líquida} (área acima do eixo menos área
abaixo). A \emph{área total} é $\int_a^b|f(x)|\,dx$.

\figtikz{\begin{tikzpicture}
\begin{axis}[painel,width=9cm,height=4.4cm,xmin=-0.2,xmax=6.6,ymin=-1.3,ymax=1.3,
  xtick={3.1416,6.2832},xticklabels={$\pi$,$2\pi$},ytick={-1,1}]
\addplot[fill=azul!25,draw=none,domain=0:3.1416] {sin(deg(x))} \closedcycle;
\addplot[fill=laranja!30,draw=none,domain=3.1416:4.7124] {sin(deg(x))} \closedcycle;
\addplot[domain=-0.2:6.6]{sin(deg(x))};
\node[azul,font=\small] at (axis cs:1.57,0.4) {$+2$}; \node[laranja,font=\small] at (axis cs:3.95,-0.4) {$-1$};
\end{axis}\end{tikzpicture}}{$\int_0^{3\pi/2}\sen x\,dx=2-1=1$ (área líquida), enquanto a área
total é $2+1=3$.}

\begin{teorema}[Propriedades da integral definida]{5.1}
Para $f$ e $g$ integráveis e $k$ constante:
\begin{enumerate}[label=(\roman*)]
\item $\int_a^a f=0$ \ e \ $\int_b^a f=-\int_a^b f$;
\item $\int_a^b(f\pm g)=\int_a^bf\pm\int_a^bg$ \ e \ $\int_a^bkf=k\int_a^bf$;
\item $\int_a^bf=\int_a^cf+\int_c^bf$;
\item se $f\le g$ em $[a,b]$, então $\int_a^bf\le\int_a^bg$; se $m\le f\le M$, então
  $m(b-a)\le\int_a^bf\le M(b-a)$.
\end{enumerate}
\end{teorema}

\textbf{Valor médio.} O \textbf{valor médio}\index{valor medio de uma funcao@valor médio de uma função} de $f$ em $[a,b]$ é
$f_{\text{méd}}=\dfrac1{b-a}\displaystyle\int_a^bf(x)\,dx$: a altura do retângulo de base
$[a,b]$ com a mesma área líquida que a região sob $f$.

% ---------------------------------------------------------------------
\section{O Teorema Fundamental do Cálculo}

\begin{teorema}[Teorema do Valor Médio para integrais]{5.2}\index{teorema@Teorema!do valor medio para integrais@do Valor Médio para integrais}
Se $f$ é contínua em $[a,b]$, existe $c\in[a,b]$ tal que $f(c)=\dfrac1{b-a}\displaystyle\int_a^bf(x)\,dx$.
\end{teorema}

Considere a \emph{função área} $F(x)=\int_a^xf(t)\,dt$: a área acumulada de $a$ até $x$. Para
$h$ pequeno, $F(x+h)-F(x)=\int_x^{x+h}f(t)\,dt\approx f(x)\,h$ (uma fatia de largura $h$ e
altura $\approx f(x)$). Isso sugere que $F'(x)=f(x)$.

\figtikz{\begin{tikzpicture}
\begin{axis}[painel,width=8cm,height=4.8cm,xmin=0,xmax=5,ymin=0,ymax=3.6,xtick={0.5,2.6,3.0},xticklabels={$a$,$x$,$x\!+\!h$},ytick=\empty]
\addplot[fill=azul!20,draw=none,domain=0.5:2.6] {1+0.6*x-0.08*x^2+0.3*sin(deg(2*x))} \closedcycle;
\addplot[fill=laranja!40,draw=none,domain=2.6:3.0] {1+0.6*x-0.08*x^2+0.3*sin(deg(2*x))} \closedcycle;
\addplot[domain=0:5]{1+0.6*x-0.08*x^2+0.3*sin(deg(2*x))};
\node[azul,font=\small] at (axis cs:1.55,0.8) {$F(x)$};
\end{axis}\end{tikzpicture}}{A função área $F(x)=\int_a^xf(t)\,dt$. Ao passar de $x$ para
$x+h$, a área cresce cerca de $f(x)\cdot h$ (faixa verde).}

\begin{teorema}[Teorema Fundamental do Cálculo, parte 1]{5.3}\index{teorema@Teorema!fundamental do calculo@Fundamental do Cálculo}
Se $f$ é contínua em $[a,b]$, então $F(x)=\int_a^xf(t)\,dt$ é derivável em $(a,b)$ e
\[
F'(x)=\frac{d}{dx}\int_a^xf(t)\,dt=f(x).
\]
\end{teorema}
\emph{Prova.} Pelo TVM para integrais, $\frac{F(x+h)-F(x)}h=\frac1h\int_x^{x+h}f(t)\,dt=f(c_h)$
para algum $c_h$ entre $x$ e $x+h$. Quando $h\to0$, $c_h\to x$ e, por continuidade,
$f(c_h)\to f(x)$. $\square$

Combinada com a regra da cadeia:
$\dfrac{d}{dx}\displaystyle\int_a^{g(x)}f(t)\,dt=f\big(g(x)\big)\,g'(x)$. Por exemplo,
$\dfrac{d}{dx}\displaystyle\int_1^{x^2}\cos t\,dt=2x\cos(x^2)$.

\begin{teorema}[Teorema Fundamental do Cálculo, parte 2]{5.4}
Se $f$ é contínua em $[a,b]$ e $F$ é qualquer antiderivada de $f$, então
\[
\int_a^bf(x)\,dx=F(b)-F(a)=\Big[F(x)\Big]_a^b .
\]
\end{teorema}
\emph{Prova.} $G(x)=\int_a^xf(t)\,dt$ também é antiderivada de $f$ (parte 1), logo
$G=F+C$. Então $\int_a^bf=G(b)-G(a)=F(b)-F(a)$. $\square$

\begin{exemplo}{5.2}{Usando o TFC}
(a) $\int_0^2x^2\,dx=\Big[\frac{x^3}3\Big]_0^2=\frac83$, confirmando o Exemplo~5.1 sem limites.

(b) $\int_0^\pi\sen x\,dx=\big[-\cos x\big]_0^\pi=1+1=2$.

(c) $\displaystyle\int_1^4\Big(\sqrt x+\frac1{x^2}\Big)dx=\Big[\tfrac23x^{3/2}-\tfrac1x\Big]_1^4
=\Big(\tfrac{16}3-\tfrac14\Big)-\Big(\tfrac23-1\Big)=\tfrac{65}{12} .$
\end{exemplo}

% ---------------------------------------------------------------------
\section{Fórmulas de integração e o Teorema da Variação Total}

Lendo a parte 2 do TFC com $F'$ no lugar de $f$: a integral de uma taxa de variação é a
variação total.

\begin{teorema}[Teorema da Variação Total]{5.5}\index{teorema@Teorema!da variacao total@da Variação Total}
$F(b)=F(a)+\displaystyle\int_a^bF'(x)\,dx$, ou seja, $\displaystyle\int_a^bF'(x)\,dx=F(b)-F(a)$.
\end{teorema}

Se $v(t)$ é a velocidade de um objeto, $\int_a^bv(t)\,dt$ é o \textbf{deslocamento}\index{deslocamento} (variação
de posição), enquanto $\int_a^b|v(t)|\,dt$ é a \textbf{distância percorrida}\index{distancia percorrida@distância percorrida}. Se $r(t)$ é a
vazão de água num tanque, $\int_a^br(t)\,dt$ é a variação do volume.

\begin{exemplo}{5.3}{Deslocamento \texorpdfstring{$\times$}{x} distância}
Uma partícula tem velocidade $v(t)=t^2-4t+3=(t-1)(t-3)$ m/s, $0\le t\le4$.
\solucao
Deslocamento: $\int_0^4v\,dt=\big[\frac{t^3}3-2t^2+3t\big]_0^4=\frac43$ m.
Para a distância, separamos onde $v$ muda de sinal ($t=1$ e $t=3$):
$\int_0^1v=\frac43$, $\int_1^3v=-\frac43$, $\int_3^4v=\frac43$. Distância
$=\frac43+\frac43+\frac43=4$ m.
\end{exemplo}

\textbf{Funções pares e ímpares.}\index{simetria em integrais} Se $f$ é par, $\int_{-a}^af=2\int_0^af$; se $f$ é ímpar,
$\int_{-a}^af=0$ (as áreas se cancelam). Por exemplo, $\int_{-2}^2x^3\cos x\,dx=0$ sem nenhuma
conta.

% ---------------------------------------------------------------------
\section{Substituição}

A regra da cadeia diz que $\frac{d}{dx}F(g(x))=F'(g(x))\,g'(x)$. Lida ao contrário:
\[
\int f\big(g(x)\big)\,g'(x)\,dx=\int f(u)\,du,\qquad u=g(x),\ du=g'(x)\,dx .
\]

\begin{estrategia}{integração por substituição}\index{substituicao@substituição}
\begin{enumerate}[label=\arabic*.]
\item Escolha $u=g(x)$, em geral a ``função de dentro'', cuja derivada aparece (a menos de
  uma constante) no integrando.
\item Calcule $du=g'(x)\,dx$ e reescreva \emph{toda} a integral em termos de $u$.
\item Integre em $u$ e volte para $x$. Em integrais definidas, troque também os limites:
  $x=a\mapsto u=g(a)$ e $x=b\mapsto u=g(b)$.
\end{enumerate}
\end{estrategia}

\begin{exemplo}{5.4}{Substituições}
(a) $\int x^2(x^3+5)^9\,dx$: com $u=x^3+5$, $du=3x^2dx$, fica
$\frac13\int u^9\,du=\frac{u^{10}}{30}+C=\frac{(x^3+5)^{10}}{30}+C$.

(b) $\int\sen(5x)\,dx=-\frac15\cos(5x)+C$ \quad ($u=5x$).

(c) $\int x\sqrt{x-1}\,dx$: com $u=x-1$, temos $x=u+1$ e
$\int(u+1)u^{1/2}\,du=\frac25u^{5/2}+\frac23u^{3/2}+C=\frac25(x-1)^{5/2}+\frac23(x-1)^{3/2}+C$.

(d) $\int_0^{\pi/2}\sen^2x\cos x\,dx$: com $u=\sen x$, os limites viram $0$ e $1$:
$\int_0^1u^2\,du=\frac13$.
\end{exemplo}

% ---------------------------------------------------------------------
\section{Integrais com exponenciais e logaritmos}

Das derivadas do capítulo~3:
\[
\int e^u\,du=e^u+C,\qquad \int a^u\,du=\frac{a^u}{\ln a}+C,\qquad \int\frac{du}u=\ln|u|+C .
\]

\begin{exemplo}{5.5}{Exponenciais e logaritmos}
(a) $\int xe^{-x^2}dx=-\frac12\int e^u\,du=-\frac12e^{-x^2}+C$ \quad ($u=-x^2$).

(b) $\int\frac{3x^2}{x^3+1}\,dx=\ln|x^3+1|+C$: o numerador é a derivada do denominador.

(c) $\int\tg x\,dx=\int\frac{\sen x}{\cos x}\,dx=-\ln|\cos x|+C=\ln|\sec x|+C$ \quad
($u=\cos x$).

(d) $\int_1^e\frac{\ln x}x\,dx=\int_0^1u\,du=\frac12$ \quad ($u=\ln x$).
\end{exemplo}

% ---------------------------------------------------------------------
\section{Integrais que resultam em inversas trigonométricas}

Das derivadas das inversas trigonométricas, para $a>0$:
\[
\int\frac{du}{\sqrt{a^2-u^2}}=\arcsen\frac ua+C,\qquad
\int\frac{du}{a^2+u^2}=\frac1a\arctg\frac ua+C,\qquad
\int\frac{du}{u\sqrt{u^2-a^2}}=\frac1a\arcsec\frac{|u|}a+C .
\]

\begin{exemplo}{5.6}{Inversas trigonométricas}
(a) $\displaystyle\int\frac{dx}{\sqrt{4-9x^2}}$: com $u=3x$, $\frac13\int\frac{du}{\sqrt{4-u^2}}
=\frac13\arcsen\frac{3x}2+C$.

(b) $\displaystyle\int_0^1\frac{dx}{1+x^2}=\big[\arctg x\big]_0^1=\frac\pi4$. Daí,
$\pi=\int_0^1\frac{4}{1+x^2}\,dx$.

(c) $\displaystyle\int\frac{dx}{x^2+4x+13}$: completando o quadrado, $x^2+4x+13=(x+2)^2+9$, e
com $u=x+2$: $\frac13\arctg\frac{x+2}3+C$.
\end{exemplo}

% ---------------------------------------------------------------------
\section*{Resumo do capítulo}\markright{Resumo do capítulo}
\addcontentsline{toc}{section}{Resumo do capítulo}
\begin{itemize}
\item Somas de Riemann $\sum f(x_i^*)\Delta x$ aproximam a área; $\int_a^bf$ é o seu limite e
  mede a área com sinal.
\item Propriedades: linearidade, aditividade em intervalos, inversão de limites, comparação;
  valor médio $\frac1{b-a}\int_a^bf$.
\item TFC 1: $\frac d{dx}\int_a^xf(t)\,dt=f(x)$ (com a cadeia, $f(g(x))g'(x)$).
  TFC 2: $\int_a^bf=F(b)-F(a)$.
\item Variação total: $\int_a^bF'=F(b)-F(a)$. Deslocamento $=\int v$; distância
  $=\int|v|$.
\item Simetria: $\int_{-a}^a$(ímpar) $=0$; $\int_{-a}^a$(par) $=2\int_0^a$.
\item Substituição $u=g(x)$, $du=g'(x)dx$; nas definidas, troque os limites.
\item $\int e^u=e^u$, $\int\frac{du}u=\ln|u|$, $\int\frac{du}{\sqrt{a^2-u^2}}=\arcsen\frac ua$,
  $\int\frac{du}{a^2+u^2}=\frac1a\arctg\frac ua$.
\end{itemize}

% Gerado por ferramentas/gerar.py; não edite à mão.
\section*{Exercícios}\markright{Exercícios}
\addcontentsline{toc}{section}{Exercícios}
Os pontos antes de cada número indicam a dificuldade: \nivel{1} fácil, \nivel{2} médio, \nivel{3} difícil e \nivel{4} desafio. Os exercícios marcados com \estrela{} têm solução completa, com animação do gráfico, no repositório do livro: \url{https://github.com/AleffAzevedo/calculo-1-pt/blob/main/exercicios/cap5/README.md}. As respostas estão no fim do capítulo.

\begin{itemize}[leftmargin=5.2em,labelsep=0.6em,itemsep=4pt]
\item[\exnum{1}{5.1}] Calcule $\displaystyle\sum_{i=1}^{20}(3i-2)$ e $\displaystyle\sum_{i=1}^{10}i^2$ usando as fórmulas de somatório.
\item[\exnum{1}{5.2}] Calcule $L_4$ e $R_4$ para $f(x)=x^2+1$ em $[0,2]$. Qual subestima a área? Por quê?
\item[\exnum{1}{5.3}] Calcule a soma do ponto médio $M_4$ para $f(x)=\dfrac1x$ em $[1,3]$.
\item[\exnum{1}{5.4}] Interprete $\int_0^3\sqrt{9-x^2}\,dx$ como área e calcule-a sem antiderivadas.
\item[\exnum{1}{5.5}] Use áreas de figuras geométricas para calcular $\int_{-2}^{4}|x-1|\,dx$.
\item[\exnum{1}{5.6}] Sabendo que $\int_0^5f=7$, $\int_0^3f=4$ e $\int_0^5g=-2$, calcule $\int_3^5f$, $\int_5^0(2f-3g)$ e $\int_3^3g$.
\item[\exnum{1}{5.7}] Calcule o valor médio de $f(x)=x^2$ em $[0,3]$ e encontre $c$ com $f(c)=f_{\text{méd}}$.
\item[\exnum{1}{5.8}] Calcule $\dfrac{d}{dx}\displaystyle\int_2^x\frac{\sin t}{t}\,dt$ e $\dfrac{d}{dx}\displaystyle\int_0^{x^3}\sqrt{1+t^2}\,dt$.
\item[\exnum{1}{5.9}] Calcule $\displaystyle\int_1^2\Big(3x^2-\frac2{x^3}\Big)dx$.
\item[\exnum{1}{5.10}] Calcule $\displaystyle\int_0^{\pi/3}\sec^2x\,dx$ e $\displaystyle\int_0^{\pi}(\cos x+2)\,dx$.
\item[\exnum{1}{5.11}] Calcule $\displaystyle\int_1^4\frac{x-1}{\sqrt x}\,dx$.
\item[\exnum{1}{5.12}] Calcule $\displaystyle\int_0^1(e^x-2^x)\,dx$.
\item[\exnum{1}{5.13}] Calcule $\displaystyle\int(2x+1)^7\,dx$ e $\displaystyle\int\cos(3x-1)\,dx$.
\item[\exnum{1}{5.14}] Calcule $\displaystyle\int x\sqrt{x^2+1}\,dx$.
\item[\exnum{1}{5.15}] Calcule $\displaystyle\int\frac{x}{x^2+4}\,dx$ e $\displaystyle\int\frac{dx}{2x+5}$.
\item[\exnum{1}{5.16}] Calcule $\displaystyle\int e^{5x}\,dx$ e $\displaystyle\int x^2e^{x^3}\,dx$.
\item[\exnum{1}{5.17}] Calcule $\displaystyle\int\frac{dx}{\sqrt{9-x^2}}$ e $\displaystyle\int\frac{dx}{16+x^2}$.
\item[\exnum{1}{5.18}] Calcule $\displaystyle\int_0^{\pi/2}\sin x\cos^3x\,dx$.
\item[\exnum{1}{5.19}] Calcule $\displaystyle\int_{-1}^{1}(x^5-3x^3+x^2)\,dx$ usando simetria.
\item[\exnum{1}{5.20}] Uma partícula tem $v(t)=2t-4$ m/s em $[0,5]$. Calcule o deslocamento e a distância percorrida.
\item[\exnum{2}{5.21}] \estrela\ \textbf{Somas de Riemann convergindo.} Para $f(x)=4-x^2$ em $[0,2]$, calcule $L_4$ e $R_4$, escreva $R_n$ em função de $n$ e
encontre a área exata como $\lim_{n\to\infty}R_n$.
\item[\exnum{2}{5.22}] \estrela\ \textbf{Integral pela definição.} Use a definição (limite de somas de Riemann com extremidades direitas) para calcular
$\int_0^1x^3\,dx$.
\item[\exnum{2}{5.23}] Escreva $\displaystyle\lim_{n\to\infty}\sum_{i=1}^n\Big(1+\frac{2i}n\Big)^2\frac2n$ como integral definida e calcule.
\item[\exnum{2}{5.24}] Mostre, sem calcular, que $2\le\displaystyle\int_0^2\sqrt{1+x^3}\,dx\le6$.
\item[\exnum{2}{5.25}] \estrela\ \textbf{A função área.} Seja $g(x)=\displaystyle\int_0^x(t-1)(t-3)\,dt$ para $0\le x\le4$. Use o TFC para encontrar
$g'(x)$, os intervalos em que $g$ cresce e decresce, e os extremos locais de $g$. Calcule
$g(1)$, $g(3)$ e $g(4)$.
\item[\exnum{2}{5.26}] \estrela\ \textbf{TFC com limites variáveis.} Calcule $\dfrac{d}{dx}\displaystyle\int_x^{x^2}e^{t^2}\,dt$.
\item[\exnum{2}{5.27}] Calcule $\dfrac{d}{dx}\displaystyle\int_{\cos x}^{1}\frac{dt}{1+t^2}$.
\item[\exnum{2}{5.28}] Calcule $\displaystyle\int_0^3|x^2-4|\,dx$.
\item[\exnum{2}{5.29}] \estrela\ \textbf{Deslocamento e distância.} Uma partícula tem velocidade $v(t)=3t^2-12t+9$ m/s para $0\le t\le4$. Calcule o
deslocamento e a distância percorrida. Se $s(0)=2$, onde ela está em $t=4$?
\item[\exnum{2}{5.30}] \estrela\ \textbf{Valor médio de uma função.} Calcule o valor médio de $f(x)=\sin x$ em $[0,\pi]$ e encontre todos os $c$ garantidos pelo
TVM para integrais.
\item[\exnum{2}{5.31}] Calcule $\displaystyle\int x\sqrt{x+2}\,dx$.
\item[\exnum{2}{5.32}] Calcule $\displaystyle\int\frac{\cos x}{\sin^2x}\,dx$ e $\displaystyle\int\sec^2x\,\tan^3x\,dx$.
\item[\exnum{2}{5.33}] \estrela\ \textbf{Substituição em integral definida.} Calcule $\displaystyle\int_0^2x\sqrt{4-x^2}\,dx$ trocando os limites de integração, e interprete
a substituição geometricamente.
\item[\exnum{2}{5.34}] Calcule $\displaystyle\int_0^{1}\frac{e^x}{1+e^x}\,dx$.
\item[\exnum{2}{5.35}] Calcule $\displaystyle\int\frac{(\ln x)^2}{x}\,dx$ e $\displaystyle\int\frac{dx}{x\ln x}$.
\item[\exnum{2}{5.36}] Calcule $\displaystyle\int\frac{e^{x}}{1+e^{2x}}\,dx$.
\item[\exnum{2}{5.37}] Calcule $\displaystyle\int_0^{1/2}\frac{dx}{\sqrt{1-x^2}}$ e $\displaystyle\int_{0}^{\sqrt3}\frac{dx}{1+x^2}$.
\item[\exnum{2}{5.38}] Calcule $\displaystyle\int\frac{dx}{x^2-6x+13}$ completando o quadrado.
\item[\exnum{2}{5.39}] \estrela\ \textbf{Simetria salva contas.} Calcule $\displaystyle\int_{-\pi}^{\pi}\big(x^3\cos x+2\cos x+x\big)\,dx$ usando simetria.
\item[\exnum{2}{5.40}] Calcule $\displaystyle\int\frac{x+3}{x^2+1}\,dx$ (separe em duas integrais).
\item[\exnum{3}{5.41}] \estrela\ \textbf{Calculando $\pi$ com uma integral.} Mostre que $\displaystyle\int_0^1\frac{4}{1+x^2}\,dx=\pi$ e use a regra do ponto médio com $n=4$ para
estimar $\pi$.
\item[\exnum{3}{5.42}] \estrela\ \textbf{Vazão e volume.} Água entra num tanque à taxa $r(t)=200-4t$ L/min, $0\le t\le50$, que inicialmente tem
1000 L. Quanto há no tanque em $t=50$? Em que instante o volume cresce mais rapidamente?
Esboce $V(t)$.
\item[\exnum{3}{5.43}] \estrela\ \textbf{Integral da tangente e da secante.} Mostre que $\int\tan x\,dx=\ln|\sec x|+C$ e que $\int\sec x\,dx=\ln|\sec x+\tan x|+C$.
Calcule $\int_0^{\pi/4}\sec x\,dx$.
\item[\exnum{3}{5.44}] Use somas de Riemann (extremidades direitas) para calcular $\int_1^3(2x+1)\,dx$ pela definição.
\item[\exnum{3}{5.45}] Encontre $f$ contínua e a constante $a$ tais que $\displaystyle6+\int_a^x\frac{f(t)}{t^2}\,dt=2\sqrt x$ para $x>0$.
\item[\exnum{3}{5.46}] Encontre os extremos locais e os pontos de inflexão de $F(x)=\displaystyle\int_0^x\frac{t-2}{1+t^2}\,dt$, $x\ge0$.
\item[\exnum{3}{5.47}] Calcule $\displaystyle\lim_{x\to0}\frac1{x^3}\int_0^x\sin(t^2)\,dt$.
\item[\exnum{3}{5.48}] Calcule $\displaystyle\int_0^{\pi}|\cos x|\,dx$ e $\displaystyle\int_0^{2\pi}|\sin x-\cos x|\,dx$.
\item[\exnum{3}{5.49}] A taxa de vendas de um produto é $r(t)=100\,t\,e^{-t^2/10}$ unidades/semana. Quantas unidades são vendidas nas 4 primeiras semanas? E no longo prazo?
\item[\exnum{3}{5.50}] Calcule $\displaystyle\int\frac{x^3}{\sqrt{1-x^2}}\,dx$ (dica: $u=1-x^2$).
\item[\exnum{3}{5.51}] Calcule $\displaystyle\int\frac{dx}{\sqrt{e^{2x}-1}}$ (dica: $u=e^x$).
\item[\exnum{3}{5.52}] Calcule $\displaystyle\int\frac{\sin x}{1+\cos^2x}\,dx$ e $\displaystyle\int\frac{dx}{x\sqrt{x^2-4}}$.
\item[\exnum{3}{5.53}] Calcule $\displaystyle\int\frac{dx}{\sqrt{-x^2+4x}}$ completando o quadrado.
\item[\exnum{3}{5.54}] Mostre que, se $f$ é contínua e periódica de período $T$, então $\int_a^{a+T}f$ não depende de $a$.
\item[\exnum{3}{5.55}] Um carro freia com $a(t)=-8$ m/s² a partir de $30$ m/s. Use integrais para achar a distância até parar.
\item[\exnum{3}{5.56}] O valor eficaz (RMS) de $f$ em $[a,b]$ é $\sqrt{\tfrac1{b-a}\int_a^bf^2}$. Calcule o RMS de $f(t)=\sin t$ em $[0,2\pi]$. (Dica: $\sin^2t=\tfrac{1-\cos2t}2$.)
\item[\exnum{3}{5.57}] Seja $g(x)=\int_0^x f(t)\,dt$, com $f$ crescente e $f(0)=0$. Mostre que $g$ é côncava para cima e que $g(x)\ge0$ para $x\ge0$.
\item[\exnum{3}{5.58}] Calcule $\displaystyle\lim_{n\to\infty}\frac1n\sum_{i=1}^n\sqrt{\frac in}$ interpretando como integral.
\item[\exnum{3}{5.59}] Calcule $\displaystyle\lim_{n\to\infty}\sum_{i=1}^n\frac{n}{n^2+i^2}$.
\item[\exnum{3}{5.60}] Calcule $\displaystyle\int_0^{\ln2}e^x\sqrt{e^x-1}\,dx$ e $\displaystyle\int_1^{e^2}\frac{dx}{x\sqrt{\ln x}}$.
\item[\exnum{4}{5.61}] Prove que $\displaystyle\int_0^{\pi/2}\frac{\sin^nx}{\sin^nx+\cos^nx}\,dx=\frac\pi4$ para todo $n>0$. (Dica: $u=\tfrac\pi2-x$.)
\item[\exnum{4}{5.62}] Mostre que $\ln x=\int_1^x\frac{dt}{t}$ satisfaz $\ln(ab)=\ln a+\ln b$ usando apenas propriedades da integral e substituição.
\item[\exnum{4}{5.63}] Calcule $\displaystyle\int_0^1\frac{\ln(1+x)}{1+x^2}\,dx$. (Dica: $x=\tan\theta$ e depois $\theta\mapsto\tfrac\pi4-\theta$.)
\item[\exnum{4}{5.64}] Encontre todas as funções contínuas $f$ com $\displaystyle\int_0^xf(t)\,dt=x^2+\int_x^1t^2f(t)\,dt$ para todo $x$.
\item[\exnum{4}{5.65}] Prove a desigualdade $\ln(1+n)<1+\tfrac12+\dots+\tfrac1n<1+\ln n$ comparando somas com $\int\frac{dx}x$.
\end{itemize}
\subsection*{Respostas}\markright{Respostas}
{\small\begin{multicols}{2}\raggedright
\begin{itemize}[leftmargin=2.7em,labelsep=0.5em,itemsep=1pt]
\item[\textbf{5.1}] $590$ e $385$.
\item[\textbf{5.2}] $L_4=3{,}75$; $R_4=5{,}75$; $L_4$ subestima, pois $f$ é crescente.
\item[\textbf{5.3}] $M_4=\tfrac{2}{5}+\tfrac27+\tfrac29+\tfrac2{11}\approx1{,}0898$ (exato: $\ln3\approx1{,}0986$).
\item[\textbf{5.4}] Um quarto do círculo de raio 3: $\tfrac{9\pi}4$.
\item[\textbf{5.5}] $\tfrac92+\tfrac92=9$.
\item[\textbf{5.6}] $3$; $-20$; $0$.
\item[\textbf{5.7}] $3$; $c=\sqrt3$.
\item[\textbf{5.8}] $\tfrac{\sin x}x$; $3x^2\sqrt{1+x^6}$.
\item[\textbf{5.9}] $\tfrac{25}4$.
\item[\textbf{5.10}] $\sqrt3$; $2\pi$.
\item[\textbf{5.11}] $\tfrac83$.
\item[\textbf{5.12}] $e-1-\tfrac1{\ln2}\approx0{,}276$.
\item[\textbf{5.13}] $\tfrac{(2x+1)^8}{16}+C$; $\tfrac13\sin(3x-1)+C$.
\item[\textbf{5.14}] $\tfrac13(x^2+1)^{3/2}+C$.
\item[\textbf{5.15}] $\tfrac12\ln(x^2+4)+C$; $\tfrac12\ln|2x+5|+C$.
\item[\textbf{5.16}] $\tfrac15e^{5x}+C$; $\tfrac13e^{x^3}+C$.
\item[\textbf{5.17}] $\sin^{-1}\tfrac x3+C$; $\tfrac14\tan^{-1}\tfrac x4+C$.
\item[\textbf{5.18}] $\tfrac14$.
\item[\textbf{5.19}] $\tfrac23$.
\item[\textbf{5.20}] $5$ m e $13$ m.
\item[\textbf{5.21}] $L_4=6{,}25$; $R_4=4{,}25$; $R_n=8-\tfrac{4(n+1)(2n+1)}{3n^2}$; área $\tfrac{16}3$.
\item[\textbf{5.22}] $\tfrac14$.
\item[\textbf{5.23}] $\int_1^3x^2\,dx=\tfrac{26}3$.
\item[\textbf{5.24}] $1\le\sqrt{1+x^3}\le3$ em $[0,2]$.
\item[\textbf{5.25}] $g'(x)=(x-1)(x-3)$; máx. local $g(1)=\tfrac43$; mín. local $g(3)=0$; $g(4)=\tfrac43$.
\item[\textbf{5.26}] $2x\,e^{x^4}-e^{x^2}$.
\item[\textbf{5.27}] $\dfrac{\sin x}{1+\cos^2x}$.
\item[\textbf{5.28}] $\tfrac{23}3$.
\item[\textbf{5.29}] Deslocamento 4 m; distância 12 m; $s(4)=6$.
\item[\textbf{5.30}] $f_{\text{méd}}=\tfrac2\pi\approx0{,}637$; $c=\sin^{-1}\tfrac2\pi\approx0{,}690$ ou $c=\pi-0{,}690\approx2{,}451$.
\item[\textbf{5.31}] $\tfrac25(x+2)^{5/2}-\tfrac43(x+2)^{3/2}+C$.
\item[\textbf{5.32}] $-\csc x+C$; $\tfrac14\tan^4x+C$.
\item[\textbf{5.33}] $\tfrac83$.
\item[\textbf{5.34}] $\ln\tfrac{1+e}2\approx0{,}620$.
\item[\textbf{5.35}] $\tfrac13(\ln x)^3+C$; $\ln|\ln x|+C$.
\item[\textbf{5.36}] $\tan^{-1}(e^x)+C$.
\item[\textbf{5.37}] $\tfrac\pi6$; $\tfrac\pi3$.
\item[\textbf{5.38}] $\tfrac12\tan^{-1}\tfrac{x-3}2+C$.
\item[\textbf{5.39}] $0$.
\item[\textbf{5.40}] $\tfrac12\ln(x^2+1)+3\tan^{-1}x+C$.
\item[\textbf{5.41}] $4\tan^{-1}1=\pi$; $M_4\approx3{,}14680$.
\item[\textbf{5.42}] $6000$ L; em $t=0$; $V(t)=1000+200t-2t^2$, parábola crescente até o vértice em $t=50$.
\item[\textbf{5.43}] $\ln(\sqrt2+1)\approx0{,}8814$.
\item[\textbf{5.44}] $10$.
\item[\textbf{5.45}] $f(x)=x^{3/2}$, $a=9$.
\item[\textbf{5.46}] Mín. local em $x=2$; inflexão em $x=2+\sqrt5$.
\item[\textbf{5.47}] $\tfrac13$ (L'Hôpital + TFC).
\item[\textbf{5.48}] $2$; $4\sqrt2$.
\item[\textbf{5.49}] $500(1-e^{-1{,}6})\approx399$; tende a $500$.
\item[\textbf{5.50}] $-\sqrt{1-x^2}+\tfrac13(1-x^2)^{3/2}+C$.
\item[\textbf{5.51}] $\sec^{-1}(e^x)+C$.
\item[\textbf{5.52}] $-\tan^{-1}(\cos x)+C$; $\tfrac12\sec^{-1}\tfrac{|x|}2+C$.
\item[\textbf{5.53}] $\sin^{-1}\tfrac{x-2}2+C$.
\item[\textbf{5.54}] Derive $G(a)=\int_a^{a+T}f$: $G'(a)=f(a+T)-f(a)=0$.
\item[\textbf{5.55}] $56{,}25$ m.
\item[\textbf{5.56}] $\tfrac1{\sqrt2}\approx0{,}707$.
\item[\textbf{5.57}] $g''=f'\ge0$; $g$ acumula valores $f(t)\ge0$.
\item[\textbf{5.58}] $\int_0^1\sqrt x\,dx=\tfrac23$.
\item[\textbf{5.59}] $\int_0^1\tfrac{dx}{1+x^2}=\tfrac\pi4$.
\item[\textbf{5.60}] $\tfrac23$; $2\sqrt2$.
\item[\textbf{5.61}] A integral $I$ satisfaz $I=\int_0^{\pi/2}\tfrac{\cos^n}{\sin^n+\cos^n}$; some as duas: $2I=\tfrac\pi2$.
\item[\textbf{5.62}] $\int_a^{ab}\tfrac{dt}t=\int_1^b\tfrac{du}u$ com $t=au$.
\item[\textbf{5.63}] $\tfrac\pi8\ln2$.
\item[\textbf{5.64}] Nenhuma: derivando, $f=\tfrac{2x}{1+x^2}$, mas em $x=0$ seria preciso $\int_0^1t^2f=0$, e essa integral vale $1-\ln2$.
\item[\textbf{5.65}] Retângulos acima e abaixo de $y=1/x$ em $[1,n+1]$ e $[1,n]$.
\end{itemize}\end{multicols}}


% =====================================================================
\chapter{Aplicações da integração}

A ideia do capítulo~5 (dividir em pedaços pequenos, aproximar cada pedaço e somar) vale
muito além da área sob uma curva. Fatiando uma região, um sólido, um fio ou um volume de
líquido, a integral calcula áreas entre curvas, volumes, comprimentos, massas, trabalho,
forças de pressão e centros de massa. O roteiro é sempre o mesmo: escrever a contribuição
de uma fatia típica de espessura $\Delta x$ e transformar a soma numa integral.

% ---------------------------------------------------------------------
\section{Áreas entre curvas}\index{area entre curvas@área entre curvas}

Se $f(x)\ge g(x)$ em $[a,b]$, uma faixa vertical de largura $\Delta x$ tem altura
$f(x)-g(x)$. Somando e passando ao limite:
\[
A=\int_a^b\big[f(x)-g(x)\big]\,dx \qquad(\text{``de cima menos de baixo''}).
\]
Se as curvas se cruzam, divida o intervalo nos pontos de interseção ou use
$\int_a^b|f(x)-g(x)|\,dx$. Quando a região é limitada à direita e à esquerda por curvas
$x=u(y)$ e $x=v(y)$, pode ser mais simples fatiar na horizontal:
$A=\int_c^d\big[u(y)-v(y)\big]\,dy$ (``da direita menos da esquerda'').

\figtikz{\begin{tikzpicture}
\begin{axis}[painel,width=6.6cm,height=5cm,xmin=-1.6,xmax=2.7,ymin=-0.6,ymax=4.8,xtick={-1,1,2},ytick={1,4}]
\addplot[fill=azul!20,draw=none,domain=-1:2,samples=60] {x+2} \closedcycle;
\addplot[fill=white,draw=none,domain=-1:2,samples=60] {x^2} \closedcycle;
\addplot[azul,thick,domain=-1.5:2.5]{x+2}; \addplot[laranja,thick,domain=-1.6:2.2,samples=60]{x^2};
\draw[fill=verde!80,draw=azul] (axis cs:0.9,0.81) rectangle (axis cs:1.1,3.0);
\node[azul,font=\small] at (axis cs:2.35,3.8) {$y=x+2$}; \node[laranja,font=\small] at (axis cs:-1.2,2.6) {$y=x^2$};
\end{axis}\end{tikzpicture}}{Região entre $y=x+2$ e $y=x^2$, com uma faixa típica de altura
$(x+2)-x^2$ e largura $\Delta x$.}

\begin{exemplo}{6.1}{Área entre duas curvas}
Calcule a área entre $y=x+2$ e $y=x^2$.
\solucao
Interseções: $x^2=x+2\Rightarrow x=-1$ ou $x=2$. Em $[-1,2]$ a reta fica por cima:
\[
A=\int_{-1}^2(x+2-x^2)\,dx=\Big[\tfrac{x^2}2+2x-\tfrac{x^3}3\Big]_{-1}^2
=\tfrac{10}3-\big(-\tfrac76\big)=\tfrac92 .
\]
\end{exemplo}

\begin{exemplo}{6.2}{Fatiando na horizontal}
Calcule a área entre $x=y^2$ e $y=x-2$.
\solucao
Na vertical, a curva de baixo muda em $x=1$ e seriam necessárias duas integrais. Em $y$:
$x=y+2$ (direita) e $x=y^2$ (esquerda), que se cruzam em $y=-1$ e $y=2$:
$A=\int_{-1}^2(y+2-y^2)\,dy=\frac92$.
\end{exemplo}

% ---------------------------------------------------------------------
\section{Volumes por fatiamento}\index{volume!por fatiamento}

Se um sólido ocupa $a\le x\le b$ e a seção transversal em $x$ tem área $A(x)$, uma fatia de
espessura $\Delta x$ tem volume $\approx A(x)\,\Delta x$. Logo
\[
V=\int_a^bA(x)\,dx .
\]

\textbf{Discos.}\index{volume!por discos} Girando a região sob $y=f(x)\ge0$, $a\le x\le b$, em torno do eixo $x$, cada
seção é um disco de raio $f(x)$:
$V=\pi\int_a^b[f(x)]^2\,dx$.

\textbf{Arruelas.}\index{volume!por arruelas} Se a região está entre $f\ge g\ge0$, a seção é um anel de raio externo
$R=f(x)$ e raio interno $r=g(x)$:
$V=\pi\int_a^b\big([f(x)]^2-[g(x)]^2\big)\,dx$. Para giros em torno do eixo $y$ (ou de outra
reta), troque os papéis: os raios são distâncias ao eixo de rotação.

\figtikz{\begin{tikzpicture}
\begin{axis}[painel,width=7cm,height=5cm,xmin=-0.3,xmax=4.6,ymin=-2.4,ymax=2.6,xtick={4},ytick={2},axis equal image=false]
\addplot[azul,thick,domain=0:4,samples=60]{sqrt(x)}; \addplot[azul,thick,dashed,domain=0:4,samples=60]{-sqrt(x)};
\draw[azul,thick] (axis cs:4,0) ellipse [x radius=0.35, y radius=2];
\draw[fill=laranja!35,draw=laranja] (axis cs:2.5,0) ellipse [x radius=0.25, y radius=1.581];
\draw[fill=azul!15,draw=none] (axis cs:0,0) -- plot[domain=0:4,samples=40] (axis cs:\x,{sqrt(\x)}) -- (axis cs:4,0) -- cycle;
\addplot[azul,thick,domain=0:4,samples=60]{sqrt(x)};
\node[font=\small] at (axis cs:2.9,2.1) {$y=\sqrt x$}; \node[laranja,font=\small] at (axis cs:2.5,-1.95) {disco de raio $\sqrt x$};
\end{axis}\end{tikzpicture}}{Girando a região sob $y=\sqrt x$ em torno do eixo $x$: cada fatia
é um disco de raio $\sqrt x$.}

\begin{exemplo}{6.3}{Discos e arruelas}
(a) Girando a região sob $y=\sqrt x$, $0\le x\le4$, em torno do eixo $x$:
$V=\pi\int_0^4x\,dx=8\pi$.

(b) Girando a região entre $y=x$ e $y=x^2$ em torno do eixo $x$ ($R=x$, $r=x^2$):
$V=\pi\int_0^1(x^2-x^4)\,dx=\pi\big(\tfrac13-\tfrac15\big)=\tfrac{2\pi}{15}$.

(c) A esfera de raio $r$ é o giro de $y=\sqrt{r^2-x^2}$:
$V=\pi\int_{-r}^r(r^2-x^2)\,dx=\tfrac43\pi r^3$.
\end{exemplo}

Sólidos que não são de revolução também funcionam: basta conhecer $A(x)$. Uma pirâmide de
base quadrada de lado $L$ e altura $h$ tem, à profundidade $y$ abaixo do vértice, seção
quadrada de lado $\frac{L}{h}y$; logo $V=\int_0^h\frac{L^2}{h^2}y^2\,dy=\frac13L^2h$.

% ---------------------------------------------------------------------
\section{Volumes por cascas cilíndricas}\index{volume!por cascas cilindricas@por cascas cilíndricas}

Girando em torno do eixo $y$ a região sob $y=f(x)\ge0$, $0\le a\le x\le b$, uma faixa
vertical de largura $\Delta x$ a distância $x$ do eixo gera uma \emph{casca cilíndrica} de
raio $x$, altura $f(x)$ e espessura $\Delta x$. ``Desenrolada'', ela é quase uma placa de
volume $2\pi x\,f(x)\,\Delta x$:
\[
V=2\pi\int_a^b x\,f(x)\,dx \qquad(\text{raio}\times\text{altura}\times\text{espessura}).
\]
Em geral: $V=2\pi\int(\text{raio da casca})(\text{altura da casca})\,d(\text{variável})$.

\figtikz{\begin{tikzpicture}
\begin{axis}[painel,width=7.4cm,height=4.8cm,xmin=-1.25,xmax=1.3,ymin=-0.05,ymax=0.36,xtick={-1,1},ytick={0.25}]
\addplot[fill=azul!15,draw=none,domain=0:1,samples=60]{x-x^2} \closedcycle;
\addplot[azul,thick,domain=0:1,samples=60]{x-x^2}; \addplot[azul,dashed,domain=-1:0,samples=60]{-x-x^2};
\draw[fill=laranja!40,draw=laranja] (axis cs:0.62,0) rectangle (axis cs:0.68,0.2356);
\draw[laranja,dashed] (axis cs:-0.65,0) rectangle (axis cs:-0.62,0.2356);
\draw[laranja] (axis cs:0,0.2356) ellipse [x radius=0.65, y radius=0.025];
\draw[<->,azul] (axis cs:0,0.29) -- node[above,font=\scriptsize]{raio $x$} (axis cs:0.65,0.29);
\end{axis}\end{tikzpicture}}{Método das cascas: a faixa em $x$, girando em torno do eixo $y$,
gera um cilindro fino de raio $x$ e altura $x-x^2$.}

\begin{exemplo}{6.4}{Cascas}
Gire a região sob $y=x-x^2$, $0\le x\le1$, em torno do eixo $y$.
\solucao
Com arruelas teríamos de resolver $y=x-x^2$ para $x$ (dois ramos). Com cascas:
$V=2\pi\int_0^1x(x-x^2)\,dx=2\pi\big(\tfrac13-\tfrac14\big)=\tfrac\pi6$.
\end{exemplo}

\begin{nota}
Escolha o método pela geometria: \textbf{discos/arruelas} usam fatias
\emph{perpendiculares} ao eixo de rotação; \textbf{cascas} usam fatias \emph{paralelas} a
ele. Prefira o que deixa os raios como funções simples da variável de integração.
\end{nota}

% ---------------------------------------------------------------------
\section{Comprimento de arco e área de superfície}\index{comprimento de arco}

Entre $x$ e $x+\Delta x$, o gráfico de $f$ é quase um segmento de comprimento
$\sqrt{\Delta x^2+\Delta y^2}=\sqrt{1+(\Delta y/\Delta x)^2}\,\Delta x$. Se $f'$ é contínua,
\[
L=\int_a^b\sqrt{1+[f'(x)]^2}\,dx .
\]
Girando o gráfico em torno do eixo $x$, cada pedaço gera uma faixa (tronco de cone) de raio
$f(x)$; a \textbf{área de superfície}\index{area de superficie@área de superfície} é
\[
S=2\pi\int_a^b f(x)\sqrt{1+[f'(x)]^2}\,dx .
\]

\begin{exemplo}{6.5}{Comprimento e superfície}
(a) $y=x^{3/2}$, $0\le x\le1$: $f'=\frac32x^{1/2}$ e
$L=\int_0^1\sqrt{1+\tfrac94x}\,dx=\tfrac8{27}\Big[\big(1+\tfrac94x\big)^{3/2}\Big]_0^1
=\tfrac{13\sqrt{13}-8}{27}\approx1{,}440$.

(b) $y=\sqrt x$, $1\le x\le4$, girando em torno do eixo $x$: como
$f\sqrt{1+f'^2}=\sqrt x\sqrt{1+\frac1{4x}}=\sqrt{x+\tfrac14}$,
\[
S=2\pi\int_1^4\sqrt{x+\tfrac14}\,dx=\tfrac{4\pi}3\Big[\big(x+\tfrac14\big)^{3/2}\Big]_1^4
=\tfrac\pi6\big(17\sqrt{17}-5\sqrt5\big)\approx30{,}85 .
\]
\end{exemplo}

Poucas funções dão integrandos de comprimento de arco elementares; em geral recorremos a
métodos numéricos. As ``boas'' são justamente as construídas para que $1+f'^2$ seja um
quadrado perfeito (veja os exercícios).

% ---------------------------------------------------------------------
\section{Aplicações físicas}

\textbf{Massa.}\index{massa} Uma barra em $[a,b]$ com densidade linear $\rho(x)$ (kg/m) tem massa
$m=\int_a^b\rho(x)\,dx$. Um disco de raio $r$ com densidade radial $\rho(r)$ (kg/m$^2$) é
fatiado em anéis de área $2\pi r\,\Delta r$: $m=\int_0^R2\pi r\,\rho(r)\,dr$.

\textbf{Trabalho.}\index{trabalho} Uma força constante $F$ ao longo de uma distância $d$ realiza trabalho
$W=Fd$. Se a força varia, $W=\int_a^bF(x)\,dx$. Para uma mola, a lei de Hooke dá $F=kx$
($x$ = deformação a partir do comprimento natural).

\textbf{Bombeamento.}\index{bombeamento} Para esvaziar um tanque, fatie o \emph{líquido} na horizontal: a camada
na altura $y$, de área $A(y)$ e espessura $\Delta y$, pesa $\rho g A(y)\,\Delta y$ e precisa
subir uma distância $D(y)$. Então $W=\int\rho g\,A(y)\,D(y)\,dy$.

\textbf{Força hidrostática.}\index{forca hidrostatica@força hidrostática} À profundidade $h$, a pressão é $p=\rho gh$. Numa placa
vertical submersa com largura $w(h)$ na profundidade $h$: $F=\int\rho g\,h\,w(h)\,dh$.
(Água: $\rho=1000$ kg/m$^3$, $g=9{,}8$ m/s$^2$, $\rho g=9800$ N/m$^3$.)

\begin{exemplo}{6.6}{Mola, tanque e comporta}
(a) Uma força de $10$ N estica uma mola $0{,}2$ m: $k=50$ N/m. O trabalho para esticá-la
$0{,}5$ m a partir do repouso é $\int_0^{0{,}5}50x\,dx=6{,}25$ J.

(b) Um tanque cilíndrico vertical de raio $2$ m e altura $5$ m está cheio de água. Para
bombear tudo pela borda superior, a camada na altura $y$ sobe $5-y$:
$W=\int_0^5 9800\cdot4\pi(5-y)\,dy=9800\cdot4\pi\cdot12{,}5\approx1{,}54\times10^6$ J.

(c) Uma comporta retangular vertical de $2$ m de largura e $3$ m de altura tem a borda
superior na superfície: $F=\int_0^3 9800\,h\cdot2\,dh=88\,200$ N.
\end{exemplo}

% ---------------------------------------------------------------------
\section{Momentos e centros de massa}

Massas $m_i$ nos pontos $x_i$ de uma reta se equilibram em $\bar x=\frac{\sum m_ix_i}{\sum m_i}$
(o numerador é o \emph{momento} em relação à origem). Para uma placa fina homogênea (lâmina)
ocupando a região entre $f\ge g$ em $[a,b]$, com área $A$, o \textbf{centroide}\index{centroide} é
\[
\bar x=\frac1A\int_a^bx\,\big[f(x)-g(x)\big]\,dx,\qquad
\bar y=\frac1A\int_a^b\frac{[f(x)]^2-[g(x)]^2}2\,dx .
\]
(Em $\bar y$, cada faixa vertical tem seu centro na altura $\frac{f+g}2$ e área $(f-g)\Delta x$.)
Simetria ajuda: se a região é simétrica em relação a uma reta, o centroide está nela.

\begin{teorema}[Teorema de Pappus]{6.1}\index{teorema@Teorema!de pappus@de Pappus}
Se uma região plana de área $A$ gira em torno de uma reta do seu plano que não a corta, o
volume do sólido gerado é $V=2\pi d\,A$, onde $d$ é a distância do centroide da região ao
eixo.
\end{teorema}

\begin{exemplo}{6.7}{Centroide e toro}
(a) Região sob $y=\sqrt x$, $0\le x\le4$: $A=\int_0^4\sqrt x\,dx=\frac{16}3$,
$\bar x=\frac3{16}\int_0^4x^{3/2}dx=\frac3{16}\cdot\frac{64}5=\frac{12}5$ e
$\bar y=\frac3{16}\int_0^4\frac x2\,dx=\frac34$.
Por Pappus, o giro em torno do eixo $x$ tem $V=2\pi\cdot\frac34\cdot\frac{16}3=8\pi$,
como no Exemplo~6.3(a).

(b) Um disco de raio $r$ com centro a distância $R>r$ de um eixo gera um toro:
$V=2\pi R\cdot\pi r^2=2\pi^2Rr^2$.
\end{exemplo}

% ---------------------------------------------------------------------
\section{Integrais, funções exponenciais e logaritmos}

O TFC permite \emph{definir} o logaritmo natural por uma integral, sem falar em expoentes:
\[
\ln x=\int_1^x\frac{dt}t,\qquad x>0 .
\]
Pelo TFC, $\frac{d}{dx}\ln x=\frac1x>0$: $\ln$ é crescente, com $\ln1=0$. A propriedade
$\ln(ab)=\ln a+\ln b$ sai de uma substituição: em $\int_a^{ab}\frac{dt}t$, faça $t=au$ e
obtenha $\int_1^b\frac{du}u=\ln b$. Daí também $\ln(a^r)=r\ln a$ para $r$ racional.

O número $e$ é definido por $\ln e=1$ (a área sob $1/t$ de $1$ a $e$ vale $1$), e a
exponencial $e^x$ é a inversa de $\ln x$. Pela derivada da inversa,
$\frac{d}{dx}e^x=\frac1{1/e^x}=e^x$. Para $a>0$, define-se $a^x=e^{x\ln a}$, de onde
$\frac{d}{dx}a^x=a^x\ln a$ e $\log_ax=\frac{\ln x}{\ln a}$. Essa construção justifica com
rigor tudo o que usamos sobre exponenciais desde o capítulo~1.

% ---------------------------------------------------------------------
\section{Crescimento e decaimento exponencial}\index{crescimento exponencial}

Quando a taxa de variação de uma quantidade é proporcional à própria quantidade,
$y'=ky$, temos $\frac{y'}y=k$; integrando, $\ln|y|=kt+C$ e
\[
y(t)=y_0\,e^{kt}\qquad(y_0=y(0)).
\]
Com $k>0$ há \textbf{crescimento}, com \textbf{tempo de duplicação}\index{tempo de duplicacao@tempo de duplicação} $T_2=\frac{\ln2}k$. Com
$k<0$ há \textbf{decaimento}, com \textbf{meia-vida}\index{meia-vida} $T_{1/2}=\frac{\ln2}{|k|}$. Juros compostos
continuamente seguem o mesmo modelo: $P(t)=P_0e^{rt}$.

\begin{exemplo}{6.8}{Datação por carbono-14}
O C-14 tem meia-vida de $5730$ anos. Um fóssil tem $30\%$ do C-14 original. Qual a sua idade?
\solucao
$k=-\frac{\ln2}{5730}$. De $0{,}30=e^{kt}$: $t=\frac{\ln(1/0{,}3)}{\ln2}\cdot5730
\approx\frac{1{,}2040}{0{,}6931}\cdot5730\approx9950$ anos.
\end{exemplo}

% ---------------------------------------------------------------------
\section{Cálculo das funções hiperbólicas}

Das definições $\senh x=\frac{e^x-e^{-x}}2$ e $\cosh x=\frac{e^x+e^{-x}}2$ (capítulo~1):
\[
\frac{d}{dx}\senh x=\cosh x,\qquad \frac{d}{dx}\cosh x=\senh x,\qquad
\frac{d}{dx}\tgh x=\operatorname{sech}^2x .
\]
Repare na ausência do sinal de menos em $(\cosh x)'$, ao contrário de $(\cos x)'$. As inversas
se escrevem com logaritmos, por exemplo $\arcsenh x=\ln\big(x+\sqrt{x^2+1}\big)$, e dão
\[
\int\frac{dx}{\sqrt{1+x^2}}=\arcsenh x+C,\qquad
\int\frac{dx}{\sqrt{x^2-1}}=\arccosh x+C\ (x>1),\qquad
\int\frac{dx}{1-x^2}=\arctgh x+C\ (|x|<1).
\]

\textbf{A catenária.}\index{catenaria@catenária} Um cabo homogêneo suspenso pelas pontas assume a forma
$y=a\cosh\frac xa$. Como $1+\senh^2=\cosh^2$, o comprimento de arco é elementar:
\[
L=\int_{-b}^b\sqrt{1+\senh^2\tfrac xa}\,dx=\int_{-b}^b\cosh\tfrac xa\,dx=2a\senh\frac ba .
\]

\begin{exemplo}{6.9}{Hiperbólicas}
(a) $\frac{d}{dx}\cosh(x^2)=2x\senh(x^2)$.

(b) $\int\tgh x\,dx=\int\frac{\senh x}{\cosh x}\,dx=\ln(\cosh x)+C$.

(c) $\int_0^1\frac{dx}{\sqrt{1+x^2}}=\arcsen h1=\ln(1+\sqrt2)\approx0{,}8814$.
\end{exemplo}

% ---------------------------------------------------------------------
\section*{Resumo do capítulo}\markright{Resumo do capítulo}
\addcontentsline{toc}{section}{Resumo do capítulo}
\begin{itemize}
\item Área entre curvas: $\int_a^b(\text{cima}-\text{baixo})\,dx$ ou
  $\int_c^d(\text{direita}-\text{esquerda})\,dy$.
\item Volume: $\int A(x)\,dx$; discos $\pi\int R^2$; arruelas $\pi\int(R^2-r^2)$; cascas
  $2\pi\int(\text{raio})(\text{altura})$.
\item Comprimento $\int\sqrt{1+f'^2}$; superfície $2\pi\int f\sqrt{1+f'^2}$.
\item Massa $\int\rho$; trabalho $\int F$ (mola $F=kx$; bombeamento $\int\rho gAD$); força
  hidrostática $\int\rho g\,h\,w(h)\,dh$.
\item Centroide: $\bar x=\frac1A\int x(f-g)$, $\bar y=\frac1A\int\frac{f^2-g^2}2$; Pappus
  $V=2\pi dA$.
\item $\ln x=\int_1^x\frac{dt}t$; $e^x$ é sua inversa; $y'=ky\Rightarrow y=y_0e^{kt}$,
  $T_2=\frac{\ln2}k$.
\item $(\senh)'=\cosh$, $(\cosh)'=\senh$; catenária $y=a\cosh\frac xa$.
\end{itemize}

% Gerado por ferramentas/gerar.py; não edite à mão.
\section*{Exercícios}
\addcontentsline{toc}{section}{Exercícios}
Os pontos antes de cada número indicam a dificuldade: \nivel{1} fácil, \nivel{2} médio, \nivel{3} difícil e \nivel{4} desafio. Os exercícios marcados com \estrela{} têm solução completa, com animação do gráfico, no repositório do livro: \url{https://github.com/AleffAzevedo/calculo-1-pt/blob/main/exercicios/cap6/README.md}. As respostas estão no fim do capítulo.

\begin{itemize}[leftmargin=5.2em,labelsep=0.6em,itemsep=4pt]
\item[\exnum{1}{6.1}] Esboce a região entre $y=x^2$ e $y=2x-x^2$ e calcule a sua área.
\item[\exnum{1}{6.2}] Calcule a área entre $y=\sin x$ e $y=\cos x$ para $0\le x\le\tfrac\pi2$ (atenção ao cruzamento).
\item[\exnum{1}{6.3}] Calcule a área total da região limitada por $y=x^3$ e $y=x$.
\item[\exnum{1}{6.4}] Calcule a área entre $x=y^2$ e $x=4$ integrando em $y$.
\item[\exnum{1}{6.5}] Calcule a área entre $y=e^x$, $y=e^{-x}$ e a reta $x=1$.
\item[\exnum{1}{6.6}] Gire em torno do eixo $x$ a região sob $y=x^2$, $0\le x\le2$. Calcule o volume.
\item[\exnum{1}{6.7}] Gire em torno do eixo $x$ a região entre $y=\sqrt x$ e $y=x$. Calcule o volume (arruelas).
\item[\exnum{1}{6.8}] Gire em torno do eixo $y$ a região sob $y=\sqrt x$, $0\le x\le4$. Calcule o volume por cascas.
\item[\exnum{1}{6.9}] Deduza o volume de um cone circular reto de raio $r$ e altura $h$ girando uma reta adequada.
\item[\exnum{1}{6.10}] Calcule o comprimento de arco de $y=2x^{3/2}$ para $0\le x\le\tfrac13$.
\item[\exnum{1}{6.11}] São necessários $6$ J para esticar uma mola $0{,}3$ m além do comprimento natural. Quanto trabalho é preciso para esticá-la de $0{,}3$ m até $0{,}5$ m?
\item[\exnum{1}{6.12}] Uma barra ocupa $0\le x\le4$ (em metros) e tem densidade $\rho(x)=2+\sqrt x$ kg/m. Calcule a massa e diga se o centro de massa fica à esquerda ou à direita de $x=2$.
\item[\exnum{1}{6.13}] Calcule o centroide da região sob $y=2x$, $0\le x\le3$, e confira com o centroide do triângulo (média dos vértices).
\item[\exnum{1}{6.14}] Uma população de bactérias dobra a cada $3$ h. Quanto tempo leva para ficar $10$ vezes maior?
\item[\exnum{1}{6.15}] Um isótopo tem meia-vida de $8$ dias. De $200$ mg, quanto resta após $30$ dias?
\item[\exnum{1}{6.16}] Mostre que $\dfrac{d}{dx}(x\sinh x-\cosh x)=x\cosh x$ e use isso para calcular $\int_0^1x\cosh x\,dx$.
\item[\exnum{1}{6.17}] Calcule $\displaystyle\int\sinh x\cosh^2x\,dx$ e $\displaystyle\int_0^{\ln2}\cosh x\,dx$.
\item[\exnum{1}{6.18}] Usando $\ln2=\int_1^2\frac{dt}t$, estime $\ln2$ pela regra do ponto médio com $n=4$ e compare com $0{,}69315$.
\item[\exnum{1}{6.19}] Uma placa quadrada de $1$ m de lado está submersa na vertical, com a borda superior a $2$ m de profundidade. Calcule a força hidrostática numa face.
\item[\exnum{1}{6.20}] A base de um sólido é o disco $x^2+y^2\le4$; as seções perpendiculares ao eixo $x$ são quadrados. Calcule o volume.
\item[\exnum{2}{6.21}] \estrela\ \textbf{Área quando as curvas se cruzam.} Calcule a área total da região limitada por $y=x^3-3x$ e $y=x$.
\item[\exnum{2}{6.22}] Calcule a área entre $y=4-x^2$ e $y=x+2$.
\item[\exnum{2}{6.23}] Calcule a área da região limitada por $y=|x|$ e $y=2-x^2$.
\item[\exnum{2}{6.24}] Calcule a área da região limitada por $y=\ln x$, $y=0$ e $x=e$ integrando em $y$.
\item[\exnum{2}{6.25}] Calcule a área total entre $y=\sin x$ e $y=\cos x$ em $[0,2\pi]$.
\item[\exnum{2}{6.26}] Encontre $k$ tal que a reta $y=k$ divida a região entre $y=4-x^2$ e o eixo $x$ em duas partes de mesma área.
\item[\exnum{2}{6.27}] \estrela\ \textbf{Arruelas em torno de uma reta horizontal.} Gire a região entre $y=x$ e $y=x^2$ ($0\le x\le1$) em torno da reta $y=2$. Calcule o volume.
\item[\exnum{2}{6.28}] Gire em torno do eixo $x$ a região sob $y=\sec x$, $0\le x\le\tfrac\pi4$. Calcule o volume.
\item[\exnum{2}{6.29}] A base de um sólido é a região entre $y=x^2$ e $y=1$; as seções perpendiculares ao eixo $x$ são semicírculos. Calcule o volume.
\item[\exnum{2}{6.30}] \estrela\ \textbf{Cascas e arruelas: dois caminhos.} Gire a região entre $y=x$ e $y=x^2$ ($0\le x\le1$) em torno da reta $x=-1$. Calcule o volume pelo
método das cascas e confira com arruelas.
\item[\exnum{2}{6.31}] Gire em torno do eixo $y$ a região sob $y=e^{-x^2}$, $0\le x\le1$.
\item[\exnum{2}{6.32}] Para a região sob $y=\frac1x$, $1\le x\le3$, calcule o volume do giro em torno do eixo $x$ e em torno do eixo $y$.
\item[\exnum{2}{6.33}] \estrela\ \textbf{Um comprimento de arco que dá certo.} Calcule o comprimento de arco de $y=\dfrac{x^3}6+\dfrac1{2x}$ para $1\le x\le3$.
\item[\exnum{2}{6.34}] Calcule o comprimento de arco de $y=\ln(\cos x)$ para $0\le x\le\tfrac\pi4$.
\item[\exnum{2}{6.35}] Calcule a área da superfície obtida girando $y=x^3$, $0\le x\le1$, em torno do eixo $x$.
\item[\exnum{2}{6.36}] Uma corrente de $10$ m com $2$ kg/m pende do alto de um prédio. Quanto trabalho é preciso para içá-la toda até o topo?
\item[\exnum{2}{6.37}] Um disco de raio $2$ m tem densidade $\rho(r)=\frac1{1+r^2}$ kg/m² ($r$ = distância ao centro). Calcule a massa.
\item[\exnum{2}{6.38}] \estrela\ \textbf{Centroide e Pappus.} Encontre o centroide da região entre $y=x$ e $y=x^2$ e use o Teorema de Pappus para recalcular o
volume do giro em torno do eixo $x$ (Exemplo 6.3b).
\item[\exnum{2}{6.39}] \estrela\ \textbf{Crescimento de uma cultura.} Uma cultura de bactérias cresce a uma taxa proporcional ao seu tamanho e passa de $1000$ para $3000$
indivíduos em $2$ h. Encontre $k$, o tempo de duplicação, o tamanho após $5$ h e quando ela chega a $10\,000$.
\item[\exnum{2}{6.40}] R\$ $1000$ são aplicados a $8\%$ ao ano com juros compostos continuamente. Em quanto tempo o capital dobra? Quanto há após $10$ anos?
\item[\exnum{3}{6.41}] \estrela\ \textbf{Esvaziando um tanque cônico.} Um tanque cônico (vértice para baixo) tem $3$ m de raio no topo e $6$ m de altura e está cheio
de água. Calcule o trabalho para bombear toda a água até $1$ m acima do topo.
\item[\exnum{3}{6.42}] \estrela\ \textbf{Força numa comporta triangular.} Uma comporta vertical tem a forma de um triângulo com base de $4$ m na superfície da água e vértice
$3$ m abaixo. Calcule a força hidrostática sobre ela.
\item[\exnum{3}{6.43}] \estrela\ \textbf{O teorema de Arquimedes da faixa esférica.} Mostre que a área da faixa de uma esfera de raio $R$ entre dois planos paralelos a distância $h$
(ambos cortando a esfera) é $2\pi Rh$, qualquer que seja a posição dos planos. Que fração da superfície da
Terra fica acima da latitude $60^\circ$?
\item[\exnum{3}{6.44}] \estrela\ \textbf{O cabo suspenso.} Um cabo pende entre dois postes de mesma altura, separados por $40$ m, na forma $y=a\cosh\frac xa$
com $a=20$ m ($-20\le x\le20$). Mostre que o comprimento é $L=2a\sinh\frac{20}a$ e calcule $L$ e a flecha
(desnível entre os apoios e o ponto mais baixo).
\item[\exnum{3}{6.45}] \estrela\ \textbf{Volume do toro: Pappus × cascas.} Gire o disco $(x-3)^2+y^2\le1$ em torno do eixo $y$. Calcule o volume do toro por Pappus e
confira pelo método das cascas.
\item[\exnum{3}{6.46}] Calcule a área da região limitada por $y=x^2-2x$ e $y=4x-x^2$.
\item[\exnum{3}{6.47}] Calcule a área da região limitada por $y=\dfrac1{1+x^2}$ e $y=\dfrac{x^2}2$.
\item[\exnum{3}{6.48}] Gire em torno da reta $x=4$ a região limitada por $y=\sqrt x$, $y=0$ e $x=4$.
\item[\exnum{3}{6.49}] Gire em torno da reta $y=4$ a região limitada por $y=x^2$ e $y=4$.
\item[\exnum{3}{6.50}] Um furo cilíndrico é feito pelo centro de uma esfera de raio $R$, deixando um anel de altura $h$. Mostre que o volume que sobra é $\tfrac16\pi h^3$, independentemente de $R$.
\item[\exnum{3}{6.51}] Calcule o comprimento de arco de $y=\cosh x$ para $0\le x\le\ln2$.
\item[\exnum{3}{6.52}] Calcule o comprimento de arco de $y=\dfrac{x^2}4-\dfrac{\ln x}2$ para $1\le x\le e$.
\item[\exnum{3}{6.53}] Girando um segmento de reta, mostre que a área lateral de um cone de raio $r$ e geratriz $\ell$ é $\pi r\ell$.
\item[\exnum{3}{6.54}] Uma tigela hemisférica de raio $2$ m está cheia de água. Calcule o trabalho para bombear toda a água até a borda.
\item[\exnum{3}{6.55}] Uma janela circular de raio $0{,}5$ m está na parede vertical de um aquário, com o centro a $2$ m de profundidade. Calcule a força na janela por integral e verifique que $F=\rho g\,\bar h\,A$ ($\bar h$ = profundidade do centroide).
\item[\exnum{3}{6.56}] Use o Teorema de Pappus e o volume da esfera para mostrar que o centroide do semidisco de raio $r$ está a $\frac{4r}{3\pi}$ do diâmetro.
\item[\exnum{3}{6.57}] A força gravitacional sobre um satélite de $1000$ kg a distância $r$ do centro da Terra é $F=\frac{gR^2m}{r^2}$ ($R=6{,}4\times10^6$ m). Calcule o trabalho para elevá-lo da superfície até $1000$ km de altitude.
\item[\exnum{3}{6.58}] Pela lei de resfriamento de Newton, $T'=-k(T-20)$. Um café a $90$\textdegree C esfria para $60$\textdegree C em $10$ min. Quando chega a $40$\textdegree C? (Dica: $y=T-20$.)
\item[\exnum{3}{6.59}] Usando $\ln x=\int_1^x\frac{dt}t$, prove que $\ln x\le x-1$ para todo $x>0$ e conclua que $\tfrac12\le\ln2\le1$.
\item[\exnum{3}{6.60}] Calcule $\displaystyle\int\frac{dx}{\sqrt{x^2+4x+5}}$ completando o quadrado.
\item[\exnum{4}{6.61}] Um plano corta uma esfera de raio $R$. Mostre que uma calota de altura $h$ tem volume $\tfrac13\pi h^2(3R-h)$ e encontre $h$ para que o plano divida a esfera em volumes na razão $1:3$.
\item[\exnum{4}{6.62}] Mostre que para $y=\cosh x$ o comprimento de arco de $0$ a $b$ é igual à área sob a curva de $0$ a $b$. Prove que, entre as funções com $f(0)=1$ e $f'\ge0$, só $\cosh$ tem essa propriedade para todo $b$.
\item[\exnum{4}{6.63}] Calcule o volume comum a dois cilindros de raio $r$ cujos eixos se cruzam em ângulo reto (sólido de Steinmetz). (Dica: as seções paralelas aos dois eixos são quadrados.)
\item[\exnum{4}{6.64}] Encontre $f>0$ derivável com $f(0)=1$ tal que, para todo $x>0$: (a) a área sob $f$ de $0$ a $x$ seja $f(x)-1$; (b) essa área seja $f(x)^2-1$.
\item[\exnum{4}{6.65}] Prove que $\displaystyle\lim_{n\to\infty}\Big(\frac1{n+1}+\frac1{n+2}+\dots+\frac1{2n}\Big)=\ln2$.
\end{itemize}
\subsection*{Respostas}
{\small\begin{multicols}{2}\raggedright
\begin{itemize}[leftmargin=2.7em,labelsep=0.5em,itemsep=1pt]
\item[\textbf{6.1}] $\tfrac13$.
\item[\textbf{6.2}] $2\sqrt2-2\approx0{,}828$.
\item[\textbf{6.3}] $\tfrac12$.
\item[\textbf{6.4}] $\tfrac{32}3$.
\item[\textbf{6.5}] $e+e^{-1}-2\approx1{,}086$.
\item[\textbf{6.6}] $\tfrac{32\pi}5$.
\item[\textbf{6.7}] $\tfrac\pi6$.
\item[\textbf{6.8}] $\tfrac{128\pi}5$.
\item[\textbf{6.9}] Girando $y=\tfrac rhx$, $0\le x\le h$: $V=\tfrac13\pi r^2h$.
\item[\textbf{6.10}] $\tfrac{14}{27}$.
\item[\textbf{6.11}] $k=\tfrac{400}3$ N/m; $W=\tfrac{32}3\approx10{,}7$ J.
\item[\textbf{6.12}] $\tfrac{40}3$ kg; à direita (a barra é mais densa na ponta $x=4$).
\item[\textbf{6.13}] $(2,\,2)$.
\item[\textbf{6.14}] $3\tfrac{\ln10}{\ln2}\approx9{,}97$ h.
\item[\textbf{6.15}] $200\cdot2^{-30/8}\approx14{,}9$ mg.
\item[\textbf{6.16}] $1-e^{-1}\approx0{,}632$.
\item[\textbf{6.17}] $\tfrac13\cosh^3x+C$; $\tfrac34$.
\item[\textbf{6.18}] $M_4\approx0{,}69122$ (erro $\approx0{,}002$).
\item[\textbf{6.19}] $9800\int_2^3h\,dh=24\,500$ N.
\item[\textbf{6.20}] $\int_{-2}^24(4-x^2)\,dx=\tfrac{128}3$.
\item[\textbf{6.21}] $8$.
\item[\textbf{6.22}] $\tfrac92$.
\item[\textbf{6.23}] $\tfrac73$.
\item[\textbf{6.24}] $\int_0^1(e-e^y)\,dy=1$.
\item[\textbf{6.25}] $4\sqrt2$.
\item[\textbf{6.26}] $(4-k)^{3/2}=4$: $k=4-2^{4/3}\approx1{,}480$.
\item[\textbf{6.27}] $\tfrac{8\pi}{15}$.
\item[\textbf{6.28}] $\pi$.
\item[\textbf{6.29}] $\int_{-1}^1\tfrac\pi8(1-x^2)^2dx=\tfrac{2\pi}{15}$.
\item[\textbf{6.30}] $\tfrac\pi2$.
\item[\textbf{6.31}] $\pi(1-e^{-1})\approx1{,}986$.
\item[\textbf{6.32}] $\tfrac{2\pi}3$ e $4\pi$.
\item[\textbf{6.33}] $\tfrac{14}3$.
\item[\textbf{6.34}] $\int_0^{\pi/4}\sec x\,dx=\ln(1+\sqrt2)\approx0{,}881$.
\item[\textbf{6.35}] $\tfrac\pi{27}(10\sqrt{10}-1)\approx3{,}563$.
\item[\textbf{6.36}] $\int_0^{10}2\cdot9{,}8\,y\,dy=980$ J.
\item[\textbf{6.37}] $\int_0^2\tfrac{2\pi r}{1+r^2}dr=\pi\ln5\approx5{,}06$ kg.
\item[\textbf{6.38}] $(\bar x,\bar y)=\big(\tfrac12,\tfrac25\big)$; $V=\tfrac{2\pi}{15}$.
\item[\textbf{6.39}] $k=\tfrac{\ln3}2\approx0{,}549$/h; $T_2\approx1{,}26$ h; $\approx15\,588$; $\approx4{,}19$ h.
\item[\textbf{6.40}] $\tfrac{\ln2}{0{,}08}\approx8{,}66$ anos; $1000e^{0{,}8}\approx$ R\$ $2225{,}54$.
\item[\textbf{6.41}] $441\,000\pi\approx1{,}39\times10^6$ J.
\item[\textbf{6.42}] $58\,800$ N.
\item[\textbf{6.43}] $S=2\pi Rh$; cerca de $6{,}7\%$.
\item[\textbf{6.44}] $L=40\sinh1\approx47{,}0$ m; flecha $20(\cosh1-1)\approx10{,}9$ m.
\item[\textbf{6.45}] $6\pi^2\approx59{,}22$.
\item[\textbf{6.46}] $9$.
\item[\textbf{6.47}] $\tfrac\pi2-\tfrac13\approx1{,}237$.
\item[\textbf{6.48}] Cascas de raio $4-x$: $\tfrac{256\pi}{15}$.
\item[\textbf{6.49}] Discos de raio $4-x^2$: $\tfrac{512\pi}{15}$.
\item[\textbf{6.50}] Cascas: $V=2\pi\int_r^R2x\sqrt{R^2-x^2}\,dx=\tfrac43\pi(R^2-r^2)^{3/2}$ com $h=2\sqrt{R^2-r^2}$.
\item[\textbf{6.51}] $\sinh(\ln2)=\tfrac34$.
\item[\textbf{6.52}] $\tfrac{e^2+1}4\approx2{,}097$.
\item[\textbf{6.53}] $y=\tfrac rhx$: $S=2\pi\int_0^h\tfrac rhx\tfrac{\ell}{h}\,dx=\pi r\ell$.
\item[\textbf{6.54}] $9800\pi\int_0^2(4-u^2)u\,du=39\,200\pi\approx1{,}23\times10^5$ J.
\item[\textbf{6.55}] $4900\pi\approx15\,394$ N.
\item[\textbf{6.56}] $\tfrac43\pi r^3=2\pi\bar y\cdot\tfrac{\pi r^2}2$.
\item[\textbf{6.57}] $gmR\big(1-\tfrac R{R+h}\big)\approx8{,}48\times10^9$ J.
\item[\textbf{6.58}] $10\tfrac{\ln3{,}5}{\ln1{,}75}\approx22{,}4$ min.
\item[\textbf{6.59}] Compare $\tfrac1t$ com $1$ (e com $\tfrac12$ em $[1,2]$).
\item[\textbf{6.60}] $\sinh^{-1}(x+2)+C=\ln\big(x+2+\sqrt{x^2+4x+5}\big)+C$.
\item[\textbf{6.61}] $t=\tfrac hR$ satisfaz $t^3-3t^2+1=0$: $h\approx0{,}6527R$.
\item[\textbf{6.62}] Derivando: $\sqrt{1+f'^2}=f$, logo $f'=\sqrt{f^2-1}$, cuja solução com $f(0)=1$ é $\cosh x$.
\item[\textbf{6.63}] $\int_{-r}^r4(r^2-x^2)\,dx=\tfrac{16}3r^3$.
\item[\textbf{6.64}] (a) $f=e^x$; (b) $f=1+\tfrac x2$.
\item[\textbf{6.65}] É uma soma de Riemann de $\int_0^1\tfrac{dx}{1+x}=\ln2$.
\end{itemize}\end{multicols}}


% =====================================================================
\chapter{Técnicas de integração}

A substituição resolve muitas integrais, mas não todas: $\int xe^x\,dx$, $\int\sen^4x\,dx$,
$\int\sqrt{9-x^2}\,dx$ e $\int\frac{dx}{x^2-1}$ pedem outras ferramentas. Este capítulo, adaptado
do capítulo~3 de \emph{Calculus Volume~2}, reúne as quatro técnicas clássicas (integração por
partes, integrais trigonométricas, substituição trigonométrica e frações parciais), os métodos
numéricos para quando nenhuma antiderivada elementar existe e as integrais impróprias, em que
o intervalo ou o integrando são ilimitados.

% ---------------------------------------------------------------------
\section{Integração por partes}\index{integracao por partes@integração por partes}

Integrando a regra do produto $(uv)'=u'v+uv'$, obtemos
\[
\int u\,dv=uv-\int v\,du,\qquad
\int_a^bu\,dv=\Big[uv\Big]_a^b-\int_a^bv\,du .
\]
A ideia é trocar uma integral difícil por outra mais simples. Escolha para $u$ a função que
\emph{simplifica ao derivar} e para $dv$ a que se integra com facilidade. Uma ordem prática
de preferência para $u$ é \textbf{LIATE}\index{liate@LIATE}: \textbf{L}ogaritmos, \textbf{I}nversas trigonométricas,
\textbf{A}lgébricas (polinômios), \textbf{T}rigonométricas, \textbf{E}xponenciais.

\begin{exemplo}{7.1}{Partes}
(a) $\int xe^x\,dx$: $u=x$, $dv=e^xdx$, então $du=dx$, $v=e^x$ e
$\int xe^x\,dx=xe^x-\int e^x\,dx=(x-1)e^x+C$.

(b) $\int\ln x\,dx$: $u=\ln x$, $dv=dx$:
$\ x\ln x-\int x\cdot\frac1x\,dx=x\ln x-x+C$.

(c) $\int x^2\sen x\,dx$ pede partes duas vezes:
$-x^2\cos x+\int2x\cos x\,dx=-x^2\cos x+2x\sen x+2\cos x+C$.

(d) $\int_0^\pi x\sen x\,dx=\big[-x\cos x\big]_0^\pi+\int_0^\pi\cos x\,dx=\pi+0=\pi$.
\end{exemplo}

\textbf{Integrais cíclicas.}\index{integral!ciclica@cíclica} Às vezes a integral original reaparece. Em
$I=\int e^x\sen x\,dx$, partes duas vezes (sempre com $dv=e^xdx$) dá
$I=e^x\sen x-e^x\cos x-I$, e daí $I=\tfrac12e^x(\sen x-\cos x)+C$.

\figtikz{\begin{tikzpicture}
\begin{axis}[painel,width=6.4cm,height=4.6cm,xmin=-0.2,xmax=3.1,ymin=-0.15,ymax=1.3,xtick={1,2.718},xticklabels={$1$,$e$},ytick={1}]
\addplot[fill=azul!25,draw=none,domain=1:2.718,samples=50]{ln(x)} \closedcycle;
\addplot[fill=laranja!30,draw=none,samples=50,domain=0:1] ({exp(x)},{x}) -- (axis cs:0,1) -- (axis cs:0,0) -- cycle;
\addplot[azul,thick,domain=0.4:3,samples=60]{ln(x)};
\draw[dashed] (axis cs:0,1) -- (axis cs:2.718,1) -- (axis cs:2.718,0);
\node[azul,font=\small] at (axis cs:2.25,0.33) {$\int_1^e\!\ln x\,dx$};
\node[laranja!80!black,font=\small] at (axis cs:0.8,0.55) {$\int_0^1\! e^y\,dy$};
\end{axis}\end{tikzpicture}}{Partes, geometricamente: as duas áreas somam o retângulo
$e\times1$, ou seja, $\int_1^e\ln x\,dx=e-\int_0^1e^y\,dy=e-(e-1)=1$.}

% ---------------------------------------------------------------------
\section{Integrais trigonométricas}\index{integrais trigonometricas@integrais trigonométricas}

\textbf{Potências de seno e cosseno}, $\int\sen^mx\cos^nx\,dx$:
\begin{itemize}
\item se uma das potências é \emph{ímpar}, separe um fator dela e converta o resto com
  $\sen^2+\cos^2=1$. Por exemplo, $n$ ímpar: guarde $\cos x\,dx$ e use $u=\sen x$;
\item se ambas são \emph{pares}, reduza o grau com
  $\sen^2x=\frac{1-\cos2x}2$, $\cos^2x=\frac{1+\cos2x}2$, $\sen x\cos x=\frac12\sen2x$.
\end{itemize}

\textbf{Potências de tangente e secante}, $\int\tg^mx\sec^nx\,dx$: se $n$ é par, guarde
$\sec^2x\,dx$ e use $u=\tg x$; se $m$ é ímpar, guarde $\sec x\tg x\,dx$ e use $u=\sec x$
(com $\tg^2=\sec^2-1$). Os demais casos costumam exigir partes, como $\int\sec^3x\,dx$.

\textbf{Produtos com frequências diferentes} usam as fórmulas de produto em soma, por exemplo
$\sen A\cos B=\frac12[\sen(A-B)+\sen(A+B)]$.

\begin{exemplo}{7.2}{Integrais trigonométricas}
(a) $\int\sen^3x\cos^2x\,dx=\int(1-\cos^2x)\cos^2x\sen x\,dx
=-\frac{\cos^3x}3+\frac{\cos^5x}5+C$ \quad ($u=\cos x$).

(b) $\int\cos^2x\,dx=\int\frac{1+\cos2x}2\,dx=\frac x2+\frac{\sen2x}4+C$.

(c) $\int\sen^2x\cos^2x\,dx=\frac14\int\sen^22x\,dx=\frac18\int(1-\cos4x)\,dx
=\frac x8-\frac{\sen4x}{32}+C$.

(d) $\int\sec^3x\,dx$: partes com $u=\sec x$, $dv=\sec^2x\,dx$ dá
$\sec x\tg x-\int\sec x\tg^2x\,dx=\sec x\tg x-\int\sec^3x\,dx+\int\sec x\,dx$, logo
$\int\sec^3x\,dx=\frac12\big(\sec x\tg x+\ln|\sec x+\tg x|\big)+C$.
\end{exemplo}

% ---------------------------------------------------------------------
\section{Substituição trigonométrica}\index{substituicao trigonometrica@substituição trigonométrica}

Raízes de expressões quadráticas desaparecem com a substituição certa, porque
$1-\sen^2=\cos^2$, $1+\tg^2=\sec^2$ e $\sec^2-1=\tg^2$:

\begin{center}\renewcommand{\arraystretch}{1.3}
\begin{tabular}{lll}
\textbf{Expressão} & \textbf{Substituição} & \textbf{Identidade}\\\hline
$\sqrt{a^2-x^2}$ & $x=a\sen\theta$, $-\frac\pi2\le\theta\le\frac\pi2$ & $a^2-a^2\sen^2\theta=a^2\cos^2\theta$\\
$\sqrt{a^2+x^2}$ & $x=a\tg\theta$, $-\frac\pi2<\theta<\frac\pi2$ & $a^2+a^2\tg^2\theta=a^2\sec^2\theta$\\
$\sqrt{x^2-a^2}$ & $x=a\sec\theta$, $0\le\theta<\frac\pi2$ & $a^2\sec^2\theta-a^2=a^2\tg^2\theta$
\end{tabular}\end{center}

Para voltar à variável $x$, desenhe o triângulo retângulo que traduz a substituição.

\figtikz{\begin{tikzpicture}[scale=0.85,font=\small]
\draw[thick] (0,0) -- (3,0) -- (3,2) -- cycle; \draw (2.75,0) |- (3,0.25);
\node[below] at (1.5,0) {$\sqrt{a^2-x^2}$}; \node[right] at (3,1) {$x$}; \node[above left] at (1.5,1) {$a$};
\node at (0.6,0.17) {$\theta$}; \node at (1.5,-0.9) {$x=a\sen\theta$};
\begin{scope}[xshift=4.8cm]
\draw[thick] (0,0) -- (3,0) -- (3,2) -- cycle; \draw (2.75,0) |- (3,0.25);
\node[below] at (1.5,0) {$a$}; \node[right] at (3,1) {$x$}; \node[above left] at (1.5,1) {$\sqrt{a^2+x^2}$};
\node at (0.6,0.17) {$\theta$}; \node at (1.5,-0.9) {$x=a\tg\theta$};
\end{scope}
\begin{scope}[xshift=9.6cm]
\draw[thick] (0,0) -- (3,0) -- (3,2) -- cycle; \draw (2.75,0) |- (3,0.25);
\node[below] at (1.5,0) {$a$}; \node[right] at (3,1) {$\sqrt{x^2-a^2}$}; \node[above left] at (1.5,1) {$x$};
\node at (0.6,0.17) {$\theta$}; \node at (1.5,-0.9) {$x=a\sec\theta$};
\end{scope}
\end{tikzpicture}}{Os três triângulos de referência da substituição trigonométrica.}

\begin{exemplo}{7.3}{Substituição trigonométrica}
(a) $\int\sqrt{9-x^2}\,dx$: com $x=3\sen\theta$, $dx=3\cos\theta\,d\theta$ e
$\sqrt{9-x^2}=3\cos\theta$:
\[
\int9\cos^2\theta\,d\theta=\tfrac92(\theta+\sen\theta\cos\theta)+C
=\tfrac92\arcsen\tfrac x3+\tfrac x2\sqrt{9-x^2}+C .
\]
(b) $\int\frac{dx}{x^2\sqrt{x^2+4}}$: com $x=2\tg\theta$, $dx=2\sec^2\theta\,d\theta$, a
integral vira $\frac14\int\frac{\cos\theta}{\sen^2\theta}\,d\theta=-\frac1{4\sen\theta}+C$.
Pelo triângulo, $\sen\theta=\frac x{\sqrt{x^2+4}}$, e o resultado é
$-\frac{\sqrt{x^2+4}}{4x}+C$.
\end{exemplo}

% ---------------------------------------------------------------------
\section{Frações parciais}

Para integrar uma função racional $\frac{P(x)}{Q(x)}$:
\begin{estrategia}{frações parciais}\index{fracoes parciais@frações parciais}
\begin{enumerate}[label=\arabic*.]
\item Se $\operatorname{grau}P\ge\operatorname{grau}Q$, divida primeiro (polinômio + fração
  própria).
\item Fatore $Q$ em fatores lineares $(ax+b)$ e quadráticos irredutíveis $(ax^2+bx+c)$.
\item Escreva uma parcela para cada fator: $\frac A{ax+b}$; para um fator repetido $(ax+b)^k$,
  as parcelas $\frac{A_1}{ax+b}+\dots+\frac{A_k}{(ax+b)^k}$; para um quadrático,
  $\frac{Bx+C}{ax^2+bx+c}$.
\item Encontre as constantes (substituindo raízes ou comparando coeficientes) e integre cada
  parcela: logaritmos, potências e arco-tangentes.
\end{enumerate}
\end{estrategia}

\begin{exemplo}{7.4}{Os três casos}
(a) $\frac{x+5}{x^2+x-2}=\frac A{x-1}+\frac B{x+2}$. Então $x+5=A(x+2)+B(x-1)$; com $x=1$,
$A=2$; com $x=-2$, $B=-1$. Logo
$\int\frac{x+5}{x^2+x-2}\,dx=2\ln|x-1|-\ln|x+2|+C$.

(b) $\frac1{x^2(x+1)}=\frac Ax+\frac B{x^2}+\frac C{x+1}$ dá $B=1$, $C=1$, $A=-1$:
$\int\frac{dx}{x^2(x+1)}=-\ln|x|-\frac1x+\ln|x+1|+C$.

(c) $\frac{2x^2-x+4}{x^3+4x}=\frac Ax+\frac{Bx+C}{x^2+4}$ dá $A=1$, $B=1$, $C=-1$:
\[
\int\frac{2x^2-x+4}{x^3+4x}\,dx=\ln|x|+\tfrac12\ln(x^2+4)-\tfrac12\arctg\tfrac x2+C .
\]
\end{exemplo}

% ---------------------------------------------------------------------
\section{Integração numérica}

Funções como $e^{-x^2}$, $\frac{\sen x}x$ e $\sqrt{1+x^3}$ não têm antiderivada elementar. Com
$\Delta x=\frac{b-a}n$ e $x_i=a+i\Delta x$:
\begin{itemize}
\item \textbf{Regra do trapézio:}\index{regra do trapezio@regra do trapézio}
  $T_n=\frac{\Delta x}2\big[f(x_0)+2f(x_1)+\dots+2f(x_{n-1})+f(x_n)\big]$.
\item \textbf{Regra de Simpson}\index{regra de simpson@regra de Simpson} ($n$ par), que aproxima $f$ por parábolas a cada dois
  subintervalos:
  $S_n=\frac{\Delta x}3\big[f(x_0)+4f(x_1)+2f(x_2)+4f(x_3)+\dots+4f(x_{n-1})+f(x_n)\big]$.
\end{itemize}

\begin{teorema}[Limitantes de erro]{7.1}\index{integracao numerica@integração numérica}
Se $|f''|\le K$ em $[a,b]$, então $|E_T|\le\dfrac{K(b-a)^3}{12n^2}$. Se $|f^{(4)}|\le M$, então
$|E_S|\le\dfrac{M(b-a)^5}{180n^4}$.
\end{teorema}

Dobrar $n$ divide o erro do trapézio por cerca de $4$ e o de Simpson por cerca de $16$.
Simpson é exata para polinômios de grau até $3$.

\figtikz{\begin{tikzpicture}
\begin{axis}[painel,width=7.4cm,height=4.6cm,xmin=0.8,xmax=2.15,ymin=0,ymax=1.15,xtick={1,1.25,1.5,1.75,2},ytick={0.5,1},xticklabel style={font=\tiny}]
\addplot[fill=laranja!25,draw=laranja] coordinates {(1,0) (1,1) (1.25,0.8) (1.5,0.6667) (1.75,0.5714) (2,0.5) (2,0)} -- cycle;
\addplot[azul,thick,domain=0.85:2.1,samples=50]{1/x};
\foreach \a in {1.25,1.5,1.75}{\edef\tmp{\noexpand\draw[laranja] (axis cs:\a,0) -- (axis cs:\a,{1/\a});}\tmp}
\end{axis}\end{tikzpicture}}{Regra do trapézio com $n=4$ para $\int_1^2\frac{dx}x$: os
trapézios ficam acima da curva convexa e superestimam a integral.}

\begin{exemplo}{7.5}{Trapézio $\times$ Simpson}
Para $\int_1^2\frac{dx}x=\ln2\approx0{,}693147$, com $n=4$:
$T_4\approx0{,}697024$ (erro $\approx3{,}9\times10^{-3}$) e $S_4\approx0{,}693254$ (erro
$\approx1{,}1\times10^{-4}$). O teorema garante $|E_T|\le\frac{2\cdot1}{12\cdot16}\approx0{,}0104$
(pois $f''=\frac2{x^3}\le2$), e o erro real é bem menor que essa cota.
\end{exemplo}

% ---------------------------------------------------------------------
\section{Integrais impróprias}\index{integral!impropria@imprópria}

\begin{definicao}
(i) Intervalo infinito: $\displaystyle\int_a^\infty f(x)\,dx=\lim_{t\to\infty}\int_a^tf(x)\,dx$
(analogamente em $-\infty$). Para $\int_{-\infty}^\infty$, divida em um ponto $c$ e exija que
\emph{as duas} partes convirjam.

(ii) Integrando ilimitado: se $f$ é descontínua em $b$,
$\displaystyle\int_a^bf(x)\,dx=\lim_{t\to b^-}\int_a^tf(x)\,dx$ (analogamente em $a$ ou num ponto
interior, dividindo o intervalo).

Se o limite existe e é finito, a integral \textbf{converge}; caso contrário, \textbf{diverge}.
\end{definicao}

\begin{exemplo}{7.6}{Convergência e divergência}
(a) $\int_1^\infty\frac{dx}{x^2}=\lim_{t\to\infty}\big(1-\frac1t\big)=1$, mas
$\int_1^\infty\frac{dx}x=\lim\ln t=\infty$. Em geral, $\int_1^\infty\frac{dx}{x^p}$ converge
(para $\frac1{p-1}$) se, e só se, $p>1$.

(b) $\int_0^1\frac{dx}{\sqrt x}=\lim_{t\to0^+}(2-2\sqrt t)=2$; já $\int_0^1\frac{dx}{x^p}$
converge se, e só se, $p<1$.

(c) $\int_{-\infty}^\infty\frac{dx}{1+x^2}=\frac\pi2+\frac\pi2=\pi$.

(d) \emph{Armadilha:} $\int_{-1}^1\frac{dx}{x^2}$ não vale $\big[-\frac1x\big]_{-1}^1=-2$
(uma integral de função positiva não pode ser negativa!). O integrando explode em $x=0$, e
$\int_0^1\frac{dx}{x^2}$ diverge.
\end{exemplo}

\begin{teorema}[Teste da comparação]{7.2}\index{teste!da comparacao@da comparação}
Se $0\le f(x)\le g(x)$ para $x\ge a$: se $\int_a^\infty g$ converge, então $\int_a^\infty f$
converge; se $\int_a^\infty f$ diverge, então $\int_a^\infty g$ diverge.
\end{teorema}
Por exemplo, $\int_1^\infty e^{-x^2}dx$ converge, pois $e^{-x^2}\le e^{-x}$ para $x\ge1$; e
$\int_1^\infty\frac{1+\sen^2x}{\sqrt x}\,dx$ diverge, pois o integrando é $\ge\frac1{\sqrt x}$.

% ---------------------------------------------------------------------
\section*{Resumo do capítulo}\markright{Resumo do capítulo}
\addcontentsline{toc}{section}{Resumo do capítulo}
\begin{itemize}
\item Partes: $\int u\,dv=uv-\int v\,du$; $u$ pela ordem LIATE; integrais cíclicas voltam a si
  mesmas.
\item Trigonométricas: potência ímpar → guarde um fator; pares → ângulo duplo;
  $\tg^m\sec^n$ → $u=\tg x$ ($n$ par) ou $u=\sec x$ ($m$ ímpar).
\item Substituição trigonométrica: $a\sen\theta$, $a\tg\theta$, $a\sec\theta$ para
  $\sqrt{a^2-x^2}$, $\sqrt{a^2+x^2}$, $\sqrt{x^2-a^2}$; volte pelo triângulo.
\item Frações parciais: divida se for imprópria, fatore, uma parcela por fator (repetidos e
  quadráticos têm forma própria).
\item Trapézio (erro $\sim\frac1{n^2}$) e Simpson (erro $\sim\frac1{n^4}$).
\item Impróprias: limite de integrais próprias; $\int_1^\infty x^{-p}$ converge sse $p>1$;
  $\int_0^1x^{-p}$ converge sse $p<1$; teste da comparação.
\end{itemize}

% Gerado por ferramentas/gerar.py; não edite à mão.
\section*{Exercícios}
\addcontentsline{toc}{section}{Exercícios}
Os pontos antes de cada número indicam a dificuldade: \nivel{1} fácil, \nivel{2} médio, \nivel{3} difícil e \nivel{4} desafio. Os exercícios marcados com \estrela{} têm solução completa, com animação do gráfico, no repositório do livro: \url{https://github.com/AleffAzevedo/calculo-1-pt/blob/main/exercicios/cap7/README.md}. As respostas estão no fim do capítulo.

\begin{itemize}[leftmargin=5.2em,labelsep=0.6em,itemsep=4pt]
\item[\exnum{1}{7.1}] Calcule $\displaystyle\int x\cos x\,dx$ e $\displaystyle\int xe^{-x}\,dx$.
\item[\exnum{1}{7.2}] Calcule $\displaystyle\int x\ln x\,dx$.
\item[\exnum{1}{7.3}] Calcule $\displaystyle\int\tan^{-1}x\,dx$ (use $dv=dx$).
\item[\exnum{1}{7.4}] Calcule $\displaystyle\int_0^{\pi/2}x\sin x\,dx$.
\item[\exnum{1}{7.5}] Calcule $\displaystyle\int\sin^3x\,dx$.
\item[\exnum{1}{7.6}] Calcule $\displaystyle\int\cos^2x\,dx$ e $\displaystyle\int_0^\pi\sin^2x\,dx$.
\item[\exnum{1}{7.7}] Calcule $\displaystyle\int\sin^2x\cos^3x\,dx$.
\item[\exnum{1}{7.8}] Calcule $\displaystyle\int\tan^2x\,dx$ e $\displaystyle\int\sec^4x\,dx$.
\item[\exnum{1}{7.9}] Calcule $\displaystyle\int\sin3x\cos2x\,dx$ usando produto em soma.
\item[\exnum{1}{7.10}] Calcule $\displaystyle\int_0^2\sqrt{4-x^2}\,dx$ com $x=2\sin\theta$ e confira com a área de um quarto de círculo.
\item[\exnum{1}{7.11}] Calcule $\displaystyle\int\frac{dx}{(x^2+9)^{3/2}}$ com $x=3\tan\theta$.
\item[\exnum{1}{7.12}] Calcule $\displaystyle\int\frac{dx}{(x-1)(x+2)}$.
\item[\exnum{1}{7.13}] Calcule $\displaystyle\int\frac{3x+1}{x^2-1}\,dx$.
\item[\exnum{1}{7.14}] Calcule $\displaystyle\int\frac{x^3}{x-1}\,dx$ (divida primeiro).
\item[\exnum{1}{7.15}] Calcule $T_4$ e $S_4$ para $\int_0^2x^3\,dx$ e compare com o valor exato. Por que $S_4$ acerta?
\item[\exnum{1}{7.16}] Decida se convergem e calcule quando for o caso: $\int_1^\infty\frac{dx}{x^3}$ e $\int_1^\infty\frac{dx}{\sqrt x}$.
\item[\exnum{1}{7.17}] Calcule $\int_0^\infty e^{-2x}\,dx$ e $\int_0^1\ln x\,dx$.
\item[\exnum{1}{7.18}] Calcule $\int_0^4\frac{dx}{\sqrt{4-x}}$ e decida se $\int_0^2\frac{dx}{(x-1)^2}$ converge.
\item[\exnum{1}{7.19}] Calcule $\displaystyle\int x^2\sin x\,dx$ e confira derivando a resposta.
\item[\exnum{1}{7.20}] Calcule $\displaystyle\int e^x\cos x\,dx$ (integral cíclica).
\item[\exnum{2}{7.21}] \estrela\ \textbf{Partes com logaritmo (e a figura do retângulo).} Calcule $\int_1^e\ln x\,dx$ e $\int_1^ex\ln x\,dx$ por partes. Interprete o primeiro resultado como
uma diferença de áreas num retângulo.
\item[\exnum{2}{7.22}] \estrela\ \textbf{Uma integral que volta a si mesma.} Calcule $\int e^x\sin x\,dx$ e depois $\int_0^\pi e^x\sin x\,dx$.
\item[\exnum{2}{7.23}] \estrela\ \textbf{A integral da secante ao cubo.} Mostre que $\int\sec^3x\,dx=\frac12\big(\sec x\tan x+\ln|\sec x+\tan x|\big)+C$ e calcule
$\int_0^{\pi/4}\sec^3x\,dx$.
\item[\exnum{2}{7.24}] \estrela\ \textbf{A área da elipse.} Use substituição trigonométrica para mostrar que a elipse $\frac{x^2}{a^2}+\frac{y^2}{b^2}\le1$ tem área $\pi ab$.
\item[\exnum{2}{7.25}] \estrela\ \textbf{Frações parciais com fator repetido.} Calcule $\displaystyle\int\frac{5x^2+3x-2}{x^3+2x^2}\,dx$ e depois a integral de $1$ a $2$.
\item[\exnum{2}{7.26}] \estrela\ \textbf{Trapézio e Simpson para $e^{-x^2}$.} A função $e^{-x^2}$ não tem antiderivada elementar. Calcule $T_4$ e $S_4$ para $\int_0^1e^{-x^2}dx$ e
use os limitantes de erro (com $|f''|\le2$ e $|f^{(4)}|\le12$ em $[0,1]$) para dizer quantas casas são confiáveis.
\item[\exnum{2}{7.27}] Calcule $\displaystyle\int x^2\cos x\,dx$.
\item[\exnum{2}{7.28}] Calcule $\displaystyle\int x^3e^{x^2}\,dx$ (substituição e depois partes).
\item[\exnum{2}{7.29}] Calcule $\displaystyle\int\sin^{-1}x\,dx$ e $\displaystyle\int(\ln x)^2\,dx$.
\item[\exnum{2}{7.30}] Calcule $\displaystyle\int\sin^4x\,dx$.
\item[\exnum{2}{7.31}] Calcule $\displaystyle\int\tan^3x\sec^3x\,dx$.
\item[\exnum{2}{7.32}] Calcule $\displaystyle\int\frac{\sqrt{x^2-9}}{x}\,dx$ ($x>3$).
\item[\exnum{2}{7.33}] Calcule $\displaystyle\int\frac{x^2}{\sqrt{9-x^2}}\,dx$.
\item[\exnum{2}{7.34}] Calcule $\displaystyle\int\frac{x^2+1}{x(x-1)^2}\,dx$.
\item[\exnum{2}{7.35}] Calcule $\displaystyle\int\frac{4}{x^3+4x}\,dx$.
\item[\exnum{2}{7.36}] Calcule $\displaystyle\int\frac{dx}{e^x+1}$ (faça $u=e^x$ e use frações parciais).
\item[\exnum{2}{7.37}] Calcule $S_4$ para $\int_0^\pi\sin x\,dx$. Qual o menor $n$ que garante, pelo teorema, erro de Simpson menor que $10^{-6}$?
\item[\exnum{2}{7.38}] Calcule $\int_0^\infty xe^{-x}\,dx$ e $\int_{-\infty}^\infty\frac{dx}{x^2+4x+5}$.
\item[\exnum{2}{7.39}] Calcule $\int_0^3\frac{dx}{(x-1)^{2/3}}$ (cuidado com $x=1$).
\item[\exnum{2}{7.40}] Use comparação para decidir: $\int_1^\infty\frac{dx}{x^2+\sqrt x}$ e $\int_2^\infty\frac{dx}{\ln x}$.
\item[\exnum{3}{7.41}] \estrela\ \textbf{O comprimento da parábola.} Calcule o comprimento de arco de $y=x^2$ para $0\le x\le1$.
\item[\exnum{3}{7.42}] \estrela\ \textbf{A trombeta de Gabriel.} Gire $y=\frac1x$, $x\ge1$, em torno do eixo $x$. Mostre que o sólido tem volume finito, mas a
superfície tem área infinita.
\item[\exnum{3}{7.43}] \estrela\ \textbf{A função gama.} Seja $\Gamma(x)=\int_0^\infty t^{x-1}e^{-t}\,dt$ para $x>0$. Mostre que $\Gamma(1)=1$, que
$\Gamma(x+1)=x\,\Gamma(x)$ e conclua que $\Gamma(n+1)=n!$.
\item[\exnum{3}{7.44}] \estrela\ \textbf{A substituição universal $t=\tan(x/2)$.} Com $t=\tan\frac x2$, mostre que $\cos x=\frac{1-t^2}{1+t^2}$ e $dx=\frac{2\,dt}{1+t^2}$, e calcule
$\displaystyle\int_0^{\pi/2}\frac{dx}{2+\cos x}$.
\item[\exnum{3}{7.45}] \estrela\ \textbf{Crescimento logístico.} Uma população segue $\dfrac{dP}{dt}=0{,}5\,P\Big(1-\dfrac P{1000}\Big)$, com $P(0)=100$. Separe as variáveis,
use frações parciais e encontre $P(t)$. Quando a população chega a $500$? Quando cresce mais rápido?
\item[\exnum{3}{7.46}] Calcule $\displaystyle\int e^{2x}\sin3x\,dx$.
\item[\exnum{3}{7.47}] Calcule $\displaystyle\int\cos\sqrt x\,dx$.
\item[\exnum{3}{7.48}] Deduza a fórmula de redução $\int\sin^nx\,dx=-\frac1n\sin^{n-1}x\cos x+\frac{n-1}n\int\sin^{n-2}x\,dx$ e calcule $\int_0^{\pi/2}\sin^6x\,dx$.
\item[\exnum{3}{7.49}] Calcule $\displaystyle\int x\tan^{-1}x\,dx$.
\item[\exnum{3}{7.50}] Calcule $\displaystyle\int\frac{dx}{(1+x^2)^2}$ e $\displaystyle\int_0^1\frac{dx}{(1+x^2)^2}$.
\item[\exnum{3}{7.51}] Calcule $\displaystyle\int\frac{e^x}{e^{2x}+3e^x+2}\,dx$.
\item[\exnum{3}{7.52}] Calcule $\displaystyle\int\frac{x^3+x+1}{(x^2+1)^2}\,dx$.
\item[\exnum{3}{7.53}] Calcule $\displaystyle\int\frac{dx}{x^3+1}$.
\item[\exnum{3}{7.54}] Calcule $\displaystyle\int\frac{\sqrt x}{1+x}\,dx$ (racionalize com $u=\sqrt x$).
\item[\exnum{3}{7.55}] Quantos subintervalos garantem, pela regra do trapézio, erro menor que $10^{-4}$ em $\int_1^2\frac{dx}x$?
\item[\exnum{3}{7.56}] Calcule $\int_0^\infty e^{-x}\sin x\,dx$.
\item[\exnum{3}{7.57}] Calcule $\int_1^\infty\frac{\ln x}{x^2}\,dx$.
\item[\exnum{3}{7.58}] Para quais $p$ a integral $\int_0^1\frac{dx}{x^p}$ converge, e para quanto? E $\int_0^\infty\frac{dx}{x^p}$?
\item[\exnum{3}{7.59}] Calcule $\displaystyle\int_0^\infty\frac{dx}{(1+x)\sqrt x}$ (é imprópria nas duas pontas).
\item[\exnum{3}{7.60}] A transformada de Laplace de $f$ é $F(s)=\int_0^\infty e^{-st}f(t)\,dt$. Calcule-a para $f(t)=1$ e $f(t)=t$ ($s>0$).
\item[\exnum{4}{7.61}] Prove que $\displaystyle\int_0^1\frac{x^4(1-x)^4}{1+x^2}\,dx=\frac{22}7-\pi$ e conclua que $\pi<\frac{22}7$.
\item[\exnum{4}{7.62}] Calcule $\displaystyle\int_0^{\pi/2}\ln(\sin x)\,dx$. (Dica: compare com $\int_0^{\pi/2}\ln(\cos x)\,dx$ e use $\sin2x=2\sin x\cos x$.)
\item[\exnum{4}{7.63}] Prove que $\displaystyle\int_1^\infty\frac{\sin x}x\,dx$ converge (partes com $u=\frac1x$), mas $\int_1^\infty\frac{|\sin x|}x\,dx$ diverge.
\item[\exnum{4}{7.64}] Sabendo que $\int_0^\infty e^{-x^2}dx=\frac{\sqrt\pi}2$, mostre que $\Gamma(\tfrac12)=\sqrt\pi$ e calcule $\Gamma(\tfrac52)$.
\item[\exnum{4}{7.65}] Sejam $I_n=\int_0^{\pi/2}\sin^nx\,dx$. Mostre que $I_{2n+1}\le I_{2n}\le I_{2n-1}$ e use a fórmula de redução para provar o produto de Wallis $\frac\pi2=\lim_{n\to\infty}\frac{2\cdot2\cdot4\cdot4\cdots2n\cdot2n}{1\cdot3\cdot3\cdot5\cdots(2n-1)(2n+1)}$.
\end{itemize}
\subsection*{Respostas}
{\small\begin{multicols}{2}\raggedright
\begin{itemize}[leftmargin=2.7em,labelsep=0.5em,itemsep=1pt]
\item[\textbf{7.1}] $x\sin x+\cos x+C$; $-(x+1)e^{-x}+C$.
\item[\textbf{7.2}] $\tfrac{x^2}2\ln x-\tfrac{x^2}4+C$.
\item[\textbf{7.3}] $x\tan^{-1}x-\tfrac12\ln(1+x^2)+C$.
\item[\textbf{7.4}] $1$.
\item[\textbf{7.5}] $-\cos x+\tfrac13\cos^3x+C$.
\item[\textbf{7.6}] $\tfrac x2+\tfrac{\sin2x}4+C$; $\tfrac\pi2$.
\item[\textbf{7.7}] $\tfrac13\sin^3x-\tfrac15\sin^5x+C$.
\item[\textbf{7.8}] $\tan x-x+C$; $\tan x+\tfrac13\tan^3x+C$.
\item[\textbf{7.9}] $-\tfrac1{10}\cos5x-\tfrac12\cos x+C$.
\item[\textbf{7.10}] $\pi$.
\item[\textbf{7.11}] $\tfrac{x}{9\sqrt{x^2+9}}+C$.
\item[\textbf{7.12}] $\tfrac13\ln\Big|\tfrac{x-1}{x+2}\Big|+C$.
\item[\textbf{7.13}] $2\ln|x-1|+\ln|x+1|+C$.
\item[\textbf{7.14}] $\tfrac{x^3}3+\tfrac{x^2}2+x+\ln|x-1|+C$.
\item[\textbf{7.15}] $T_4=4{,}25$; $S_4=4$ (exato): Simpson é exata para cúbicas.
\item[\textbf{7.16}] $\tfrac12$; diverge.
\item[\textbf{7.17}] $\tfrac12$; $-1$.
\item[\textbf{7.18}] $4$; diverge (o integrando explode em $x=1$).
\item[\textbf{7.19}] $-x^2\cos x+2x\sin x+2\cos x+C$.
\item[\textbf{7.20}] $\tfrac12e^x(\sin x+\cos x)+C$.
\item[\textbf{7.21}] $1$; $\tfrac{e^2+1}4\approx2{,}097$.
\item[\textbf{7.22}] $\tfrac12e^x(\sin x-\cos x)+C$; $\tfrac{e^\pi+1}2\approx12{,}070$.
\item[\textbf{7.23}] $\tfrac12\big(\sqrt2+\ln(1+\sqrt2)\big)\approx1{,}1478$.
\item[\textbf{7.24}] $\pi ab$.
\item[\textbf{7.25}] $2\ln|x|+\frac1x+3\ln|x+2|+C$; $2\ln2-\tfrac12+3\ln\tfrac43\approx1{,}749$.
\item[\textbf{7.26}] $T_4\approx0{,}742984$; $S_4\approx0{,}746855$; valor $\approx0{,}746824$.
\item[\textbf{7.27}] $x^2\sin x+2x\cos x-2\sin x+C$.
\item[\textbf{7.28}] $\tfrac12e^{x^2}(x^2-1)+C$.
\item[\textbf{7.29}] $x\sin^{-1}x+\sqrt{1-x^2}+C$; $x(\ln x)^2-2x\ln x+2x+C$.
\item[\textbf{7.30}] $\tfrac{3x}8-\tfrac{\sin2x}4+\tfrac{\sin4x}{32}+C$.
\item[\textbf{7.31}] $\tfrac15\sec^5x-\tfrac13\sec^3x+C$.
\item[\textbf{7.32}] $\sqrt{x^2-9}-3\sec^{-1}\tfrac x3+C$.
\item[\textbf{7.33}] $\tfrac92\sin^{-1}\tfrac x3-\tfrac x2\sqrt{9-x^2}+C$.
\item[\textbf{7.34}] $\ln|x|-\tfrac2{x-1}+C$.
\item[\textbf{7.35}] $\ln|x|-\tfrac12\ln(x^2+4)+C$.
\item[\textbf{7.36}] $x-\ln(1+e^x)+C$.
\item[\textbf{7.37}] $S_4=\tfrac\pi{12}(4\sqrt2+2)\approx2{,}00456$; $n=38$.
\item[\textbf{7.38}] $1$; $\pi$.
\item[\textbf{7.39}] $3+3\sqrt[3]2\approx6{,}780$.
\item[\textbf{7.40}] Converge ($\le\tfrac1{x^2}$); diverge ($\tfrac1{\ln x}\ge\tfrac1x$).
\item[\textbf{7.41}] $\tfrac14\big(2\sqrt5+\ln(2+\sqrt5)\big)\approx1{,}4789$.
\item[\textbf{7.42}] $V=\pi$; $S=\infty$.
\item[\textbf{7.43}] $\Gamma(n+1)=n!$.
\item[\textbf{7.44}] $\tfrac{\pi}{3\sqrt3}\approx0{,}6046$.
\item[\textbf{7.45}] $P(t)=\dfrac{1000}{1+9e^{-t/2}}$; $t=2\ln9\approx4{,}39$; crescimento máximo quando $P=500$.
\item[\textbf{7.46}] $\tfrac1{13}e^{2x}(2\sin3x-3\cos3x)+C$.
\item[\textbf{7.47}] $2\big(\sqrt x\sin\sqrt x+\cos\sqrt x\big)+C$.
\item[\textbf{7.48}] $\tfrac56\cdot\tfrac34\cdot\tfrac12\cdot\tfrac\pi2=\tfrac{5\pi}{32}$.
\item[\textbf{7.49}] $\tfrac12\big[(x^2+1)\tan^{-1}x-x\big]+C$.
\item[\textbf{7.50}] $\tfrac12\big(\tan^{-1}x+\tfrac x{1+x^2}\big)+C$; $\tfrac\pi8+\tfrac14$.
\item[\textbf{7.51}] $\ln\tfrac{e^x+1}{e^x+2}+C$.
\item[\textbf{7.52}] $\tfrac12\ln(x^2+1)+\tfrac12\big(\tan^{-1}x+\tfrac x{1+x^2}\big)+C$.
\item[\textbf{7.53}] $\tfrac13\ln|x+1|-\tfrac16\ln(x^2-x+1)+\tfrac1{\sqrt3}\tan^{-1}\tfrac{2x-1}{\sqrt3}+C$.
\item[\textbf{7.54}] $2\sqrt x-2\tan^{-1}\sqrt x+C$.
\item[\textbf{7.55}] $\tfrac{2}{12n^2}<10^{-4}$: $n\ge41$.
\item[\textbf{7.56}] $\tfrac12$.
\item[\textbf{7.57}] $1$.
\item[\textbf{7.58}] $p<1$, vale $\tfrac1{1-p}$; a segunda diverge para todo $p$.
\item[\textbf{7.59}] $\pi$ (com $u=\sqrt x$).
\item[\textbf{7.60}] $\tfrac1s$; $\tfrac1{s^2}$.
\item[\textbf{7.61}] Divida: o quociente é $x^6-4x^5+5x^4-4x^2+4$ e o resto dá $-4\tan^{-1}x$; o integrando é positivo.
\item[\textbf{7.62}] $-\tfrac\pi2\ln2$.
\item[\textbf{7.63}] Partes: $\cos1-\int_1^\infty\tfrac{\cos x}{x^2}dx$, absolutamente convergente; e $|\sin x|\ge\sin^2x=\tfrac{1-\cos2x}2$.
\item[\textbf{7.64}] $t=x^2$; $\Gamma(\tfrac52)=\tfrac32\cdot\tfrac12\sqrt\pi=\tfrac{3\sqrt\pi}4$.
\item[\textbf{7.65}] $\tfrac{I_{2n}}{I_{2n+1}}\to1$ pelo sanduíche, e essa razão é $\tfrac\pi2$ dividido pelo produto parcial.
\end{itemize}\end{multicols}}



\appendix
\renewcommand{\chaptername}{Apêndice}
\titleformat{\chapter}[display]{\normalfont\barlow\bfseries\color{azul}}
  {\large\color{laranja}APÊNDICE \thechapter}{2pt}{\Huge}[{\color{verde}\titlerule[2.5pt]}]
% =====================================================================
\chapter{Formulário}
\raggedbottom
\index{formulário}

Resumo das fórmulas usadas ao longo do livro, para consulta rápida durante os exercícios.

% ---------------------------------------------------------------------
\section{Álgebra}
\begin{multicols}{2}\small
\textbf{Fatoração}
\begin{align*}
a^2-b^2&=(a-b)(a+b)\\
a^3\pm b^3&=(a\pm b)(a^2\mp ab+b^2)\\
(a\pm b)^2&=a^2\pm2ab+b^2\\
ax^2+bx+c&=a(x-x_1)(x-x_2)
\end{align*}
com $x_{1,2}=\dfrac{-b\pm\sqrt{b^2-4ac}}{2a}$.

\textbf{Completar o quadrado}
\[
ax^2+bx+c=a\Big(x+\frac b{2a}\Big)^2+c-\frac{b^2}{4a}
\]

\columnbreak
\textbf{Expoentes e logaritmos}
\begin{gather*}
a^xa^y=a^{x+y},\quad \frac{a^x}{a^y}=a^{x-y},\quad (a^x)^y=a^{xy}\\
\ln(ab)=\ln a+\ln b,\quad \ln\frac ab=\ln a-\ln b\\
\ln a^r=r\ln a,\quad \log_ax=\frac{\ln x}{\ln a}\\
a^x=e^{x\ln a},\quad e^{\ln x}=x,\quad \ln e^x=x
\end{gather*}

\textbf{Valor absoluto}
\[
|x|<a\iff-a<x<a,\qquad |a+b|\le|a|+|b|
\]
\end{multicols}

% ---------------------------------------------------------------------
\section{Trigonometria}
\begin{multicols}{2}\small
\textbf{Valores notáveis}

\begin{tabular}{c|ccccc}
$\theta$ & $0$ & $\frac\pi6$ & $\frac\pi4$ & $\frac\pi3$ & $\frac\pi2$\\\hline
$\sin\theta$ & $0$ & $\frac12$ & $\frac{\sqrt2}2$ & $\frac{\sqrt3}2$ & $1$\\
$\cos\theta$ & $1$ & $\frac{\sqrt3}2$ & $\frac{\sqrt2}2$ & $\frac12$ & $0$\\
$\tan\theta$ & $0$ & $\frac{\sqrt3}3$ & $1$ & $\sqrt3$ & ---
\end{tabular}

\textbf{Pitagóricas}
\begin{gather*}
\sin^2x+\cos^2x=1\\
1+\tan^2x=\sec^2x,\qquad 1+\cot^2x=\csc^2x
\end{gather*}

\textbf{Paridade e complementares}
\begin{gather*}
\sin(-x)=-\sin x,\qquad \cos(-x)=\cos x\\
\sin\big(\tfrac\pi2-x\big)=\cos x,\qquad \cos\big(\tfrac\pi2-x\big)=\sin x
\end{gather*}

\columnbreak
\textbf{Soma de arcos}
\begin{gather*}
\sin(a\pm b)=\sin a\cos b\pm\cos a\sin b\\
\cos(a\pm b)=\cos a\cos b\mp\sin a\sin b
\end{gather*}

\textbf{Arco duplo e metade}
\begin{gather*}
\sin2x=2\sin x\cos x\\
\cos2x=\cos^2x-\sin^2x=2\cos^2x-1=1-2\sin^2x\\
\sin^2x=\frac{1-\cos2x}2,\qquad \cos^2x=\frac{1+\cos2x}2
\end{gather*}

\textbf{Produto em soma}
\begin{gather*}
\sin a\cos b=\tfrac12[\sin(a-b)+\sin(a+b)]\\
\sin a\sin b=\tfrac12[\cos(a-b)-\cos(a+b)]\\
\cos a\cos b=\tfrac12[\cos(a-b)+\cos(a+b)]
\end{gather*}
\end{multicols}

{\small\textbf{Hiperbólicas:} $\cosh x=\frac{e^x+e^{-x}}2$, $\sinh x=\frac{e^x-e^{-x}}2$,
$\cosh^2x-\sinh^2x=1$, $\sinh^{-1}x=\ln\big(x+\sqrt{x^2+1}\big)$,
$\cosh^{-1}x=\ln\big(x+\sqrt{x^2-1}\big)$ ($x\ge1$), $\tanh^{-1}x=\frac12\ln\frac{1+x}{1-x}$ ($|x|<1$).}

% ---------------------------------------------------------------------
\section{Limites}
{\small
\[
\lim_{x\to0}\frac{\sin x}x=1,\qquad \lim_{x\to0}\frac{1-\cos x}x=0,\qquad
\lim_{x\to\infty}\Big(1+\frac1x\Big)^x=e,\qquad \lim_{x\to0}\frac{e^x-1}x=1,
\]
\[
\lim_{x\to\infty}\frac{x^n}{e^x}=0,\qquad \lim_{x\to\infty}\frac{\ln x}{x^p}=0\ (p>0),\qquad
\lim_{x\to0^+}x^p\ln x=0\ (p>0).
\]}

% ---------------------------------------------------------------------
\section{Derivadas}
\begin{multicols}{2}\small
\textbf{Regras} ($c$ constante)
\begin{gather*}
(cf)'=cf',\qquad (f\pm g)'=f'\pm g'\\
(fg)'=f'g+fg'\\
\Big(\frac fg\Big)'=\frac{f'g-fg'}{g^2}\\
\big(f(g(x))\big)'=f'(g(x))\,g'(x)\\
(f^{-1})'(x)=\frac1{f'(f^{-1}(x))}
\end{gather*}

\textbf{Potências, exponenciais, logaritmos}
\begin{gather*}
(x^n)'=nx^{n-1},\qquad (e^x)'=e^x,\qquad (a^x)'=a^x\ln a\\
(\ln|x|)'=\frac1x,\qquad (\log_ax)'=\frac1{x\ln a}
\end{gather*}

\columnbreak
\textbf{Trigonométricas}
\begin{gather*}
(\sin x)'=\cos x,\qquad (\cos x)'=-\sin x\\
(\tan x)'=\sec^2x,\qquad (\cot x)'=-\csc^2x\\
(\sec x)'=\sec x\tan x,\qquad (\csc x)'=-\csc x\cot x
\end{gather*}

\textbf{Inversas trigonométricas}
\begin{gather*}
(\sin^{-1}x)'=\frac1{\sqrt{1-x^2}},\qquad (\cos^{-1}x)'=-\frac1{\sqrt{1-x^2}}\\
(\tan^{-1}x)'=\frac1{1+x^2},\qquad (\sec^{-1}x)'=\frac1{|x|\sqrt{x^2-1}}
\end{gather*}

\textbf{Hiperbólicas}
\begin{gather*}
(\sinh x)'=\cosh x,\qquad (\cosh x)'=\sinh x,\qquad (\tanh x)'=\operatorname{sech}^2x
\end{gather*}
\end{multicols}

% ---------------------------------------------------------------------
\clearpage
\section{Integrais}
{\small Em todas, some $+C$. Substituição: $\int f(g(x))g'(x)\,dx=\int f(u)\,du$. Partes:
$\int u\,dv=uv-\int v\,du$.}

{\small\noindent\begin{minipage}[t]{0.44\linewidth}\vspace{0pt}
\[\begin{aligned}
&\int x^n\,dx=\frac{x^{n+1}}{n+1}\quad(n\ne-1)\\
&\int\frac{dx}x=\ln|x|\\
&\int e^x\,dx=e^x,\qquad \int a^x\,dx=\frac{a^x}{\ln a}\\
&\int\ln x\,dx=x\ln x-x\\
&\int\sin x\,dx=-\cos x,\qquad \int\cos x\,dx=\sin x\\
&\int\sec^2x\,dx=\tan x,\qquad \int\csc^2x\,dx=-\cot x\\
&\int\sec x\tan x\,dx=\sec x\\
&\int\tan x\,dx=\ln|\sec x|\\
&\int\cot x\,dx=\ln|\sin x|\\
&\int\sec x\,dx=\ln|\sec x+\tan x|\\
&\int\csc x\,dx=-\ln|\csc x+\cot x|
\end{aligned}\]
\end{minipage}\hfill\begin{minipage}[t]{0.54\linewidth}\vspace{0pt}
\[\begin{aligned}
&\int\sec^3x\,dx=\tfrac12\big(\sec x\tan x+\ln|\sec x+\tan x|\big)\\
&\int\sin^2x\,dx=\frac x2-\frac{\sin2x}4\\
&\int\cos^2x\,dx=\frac x2+\frac{\sin2x}4\\
&\int\frac{dx}{\sqrt{a^2-x^2}}=\sin^{-1}\frac xa\\
&\int\frac{dx}{a^2+x^2}=\frac1a\tan^{-1}\frac xa\\
&\int\frac{dx}{x\sqrt{x^2-a^2}}=\frac1a\sec^{-1}\frac{|x|}a\\
&\int\frac{dx}{\sqrt{x^2\pm a^2}}=\ln\big|x+\sqrt{x^2\pm a^2}\big|\\
&\int\frac{dx}{a^2-x^2}=\frac1{2a}\ln\Big|\frac{a+x}{a-x}\Big|\\
&\int\sqrt{a^2-x^2}\,dx=\frac x2\sqrt{a^2-x^2}+\frac{a^2}2\sin^{-1}\frac xa\\
&\int\sinh x\,dx=\cosh x,\qquad \int\cosh x\,dx=\sinh x
\end{aligned}\]
\end{minipage}\par}

% ---------------------------------------------------------------------
\section{Somatórios e geometria}
\begin{multicols}{2}\small
\textbf{Somatórios}
\begin{gather*}
\sum_{i=1}^n i=\frac{n(n+1)}2,\qquad \sum_{i=1}^n i^2=\frac{n(n+1)(2n+1)}6\\
\sum_{i=1}^n i^3=\Big[\frac{n(n+1)}2\Big]^2,\qquad \sum_{i=0}^{n-1}r^i=\frac{1-r^n}{1-r}
\end{gather*}

\textbf{Aplicações da integral}
\begin{gather*}
A=\int_a^b(f-g),\quad V=\pi\int R^2,\quad V=2\pi\int r\,h\\
L=\int_a^b\sqrt{1+f'^2},\quad S=2\pi\int_a^bf\sqrt{1+f'^2}\\
f_{\text{méd}}=\frac1{b-a}\int_a^bf
\end{gather*}

\columnbreak
\textbf{Áreas e volumes}
\begin{gather*}
\text{círculo: }A=\pi r^2,\ C=2\pi r;\qquad \text{elipse: }A=\pi ab\\
\text{setor (ângulo $\theta$): }A=\tfrac12r^2\theta,\ s=r\theta\\
\text{esfera: }V=\tfrac43\pi r^3,\ S=4\pi r^2\\
\text{cilindro: }V=\pi r^2h;\qquad \text{cone: }V=\tfrac13\pi r^2h\\
\text{cone (lateral): }S=\pi r\ell
\end{gather*}

\textbf{Métodos numéricos}
\begin{gather*}
x_{n+1}=x_n-\frac{f(x_n)}{f'(x_n)}\ \text{(Newton)}\\
|E_T|\le\frac{K(b-a)^3}{12n^2},\qquad |E_S|\le\frac{M(b-a)^5}{180n^4}
\end{gather*}
\end{multicols}


\backmatter
\printindex

\end{document}
